% Les villes — Géographie 1re Tech — Cours signé OnBuch+
% Programme officiel MINESEC. Compiler avec : tectonic N03-villes.tex
\newcommand{\DOCMATIERE}{Géographie}
\newcommand{\DOCNIVEAU}{1re Tech}
\newcommand{\DOCMODULE}{Géographie – classe de Première technique}
\newcommand{\DOCLECON}{N03}
\newcommand{\DOCTITRE}{Les villes}
% OnBuch+ — préambule « cours signés » ENRICHI (compilé avec tectonic / XeLaTeX)
% Système visuel « Pop » d'OnBuch+ : fond crème (jamais blanc), bordures encre
% épaisses, ombres « dures » décalées, titres Archivo Black en capitales.
% Version MATHÉMATIQUES : graphes (pgfplots/tikz), boîtes pédagogiques
% (définition, propriété, méthode, attention, le savais-tu, exercice résolu,
% exercice, TP, bilan), page de garde et signature OnBuch+ ancrée en bas.
%
% Macros attendues dans main.tex AVANT \input{preamble} :
%   \DOCMATIERE  \DOCNIVEAU  \DOCMODULE  \DOCLECON (numéro)  \DOCTITRE
\documentclass[11pt]{article}

\usepackage{fontspec}
\usepackage[french]{babel}
\usepackage[a4paper,margin=2cm,top=3.4cm,bottom=2.8cm,headheight=1.4cm,footskip=1.3cm]{geometry}
\usepackage{amsmath,amssymb,amsfonts}
\usepackage{mathtools} % \xmapsto, \coloneqq... souvent utilisés par les modèles
\usepackage{cancel} % simplifications barrées (bilans de forces)
% \abs{z} (module, valeur absolue) et \norm{u} (norme), utilisés par les modèles
\providecommand{\abs}{}\let\abs\relax\DeclarePairedDelimiter{\abs}{\lvert}{\rvert}
\providecommand{\norm}{}\let\norm\relax\DeclarePairedDelimiter{\norm}{\lVert}{\rVert}
\usepackage[table]{xcolor} % [table] : fournit aussi \rowcolors (alternance de lignes)
\usepackage{pagecolor}
\usepackage{tikz}
\usetikzlibrary{arrows.meta,positioning,calc,shapes.geometric,shapes.misc,shapes.arrows,decorations.pathmorphing,decorations.pathreplacing,decorations.markings,patterns,shadows,mindmap,trees,fit,backgrounds,matrix,intersections,angles,quotes}
\usepackage{pgfplots}
\usepackage[version=4]{mhchem} % équations nucléaires \ce{^{235}_{92}U}
\usepackage[european,siunitx]{circuitikz} % schémas électriques
\pgfplotsset{compat=1.18}
\usepgfplotslibrary{fillbetween,groupplots} % aires entre courbes (calcul intégral)
\usepackage{enumitem}
\usepackage{fancyhdr}
% [most] charge toutes les bibliothèques tcolorbox (listings, minted, raster,
% documentation...) et gonfle inutilement la mémoire TeX sur un document long
% (~300 boîtes) : on ne charge que ce qui est réellement utilisé.
\usepackage[skins,breakable]{tcolorbox}
\usepackage{graphicx}
\usepackage{colortbl}
\usepackage{booktabs}
\usepackage{tabularx}
\usepackage{multirow}
\usepackage{diagbox} % case d'en-tête diagonale des tableaux à double entrée
% Colonnes extensibles centrée / gauche / droite (C, L, R), souvent utilisées par les modèles
\newcolumntype{C}{>{\centering\arraybackslash}X}
\newcolumntype{L}{>{\raggedright\arraybackslash}X}
\newcolumntype{R}{>{\raggedleft\arraybackslash}X}
\usepackage{array}
\usepackage{multicol}
\usepackage{siunitx}
% parse-numbers=false : accepte aussi \SI{2,5 \times 10^{-3}}{...}
\sisetup{locale=FR,output-decimal-marker={,},group-minimum-digits=4,parse-numbers=false}
\DeclareSIUnit\ohm{\ensuremath{\symup{\Omega}}} % U+2126 absent de la police
\DeclareSIUnit\division{div} % réglages d'oscilloscope (ms/div, V/div)
\DeclareSIUnit\tr{tr} % tours (vitesse de rotation, ex. \SI{1500}{\tr\per\minute})
\DeclareSIUnit\molar{M} % concentration molaire (ex. \SI{3}{\molar}, \SI{1}{\micro\molar})
\DeclareSIUnit\an{an} % année (le modèle confond parfois avec \year, un compteur TeX) — ex. \SI{5}{\kilo\meter\per\an}
\DeclareSIUnit\FCFA{FCFA} % franc CFA (ex. \SI{500}{\FCFA\per\kilo\gram})
\DeclareSIUnit\degreeBrix{\textdegree Brix} % teneur en sucre d'un sirop/jus (réfractométrie)
\DeclareSIUnit\month{mois} % le modèle utilise parfois \month, un compteur TeX, comme unité de durée
\DeclareSIUnit\microliter{\micro\liter} % le modèle écrit parfois \microliter en un seul token
\DeclareSIUnit\million{millions} % dénombrement (ex. \SI{15}{\million\per\mL} spermatozoïdes)
\usepackage{qrcode}
% \degree : commande fréquemment utilisée par les modèles pour un angle ou une
% température en dehors de \SI{}{} (non fournie par les paquets chargés).
\providecommand{\degree}{^{\circ}}
% \qed (fin de démonstration), utilisé par les modèles sans amsthm
\providecommand{\qed}{\hfill$\square$}
% \barre : double barre des valeurs interdites dans un tableau de variations (tkz-tab)
\providecommand{\barre}{\ensuremath{\|}}
% environnement proof (amsthm non chargé) utilisé par les modèles
\ifdefined\proof\else\newenvironment{proof}[1][Démonstration]{\par\noindent\textit{#1.}\ }{\qed\par}\fi
\usepackage{lastpage}
\usepackage{hyperref}
\hypersetup{hidelinks}

% ============================================================
% Palette « Pop » OnBuch+ — mêmes hex que la classe P (plus_ui.dart)
% ============================================================
\definecolor{poporange}{HTML}{F59321}
\definecolor{popdark}{HTML}{E07A0C}
\definecolor{popcream}{HTML}{FFF6E8}
\definecolor{popink}{HTML}{1C1714}
\definecolor{poppurple}{HTML}{7A5AE0}
\definecolor{popgreen}{HTML}{2E9E6A}
\definecolor{poppink}{HTML}{E0457B}
\definecolor{popblue}{HTML}{2F7DE1}
\definecolor{popgold}{HTML}{C98A12}
\definecolor{popmuted}{HTML}{8A7F76}
\definecolor{popline}{HTML}{EADFD2}
\definecolor{popgray}{HTML}{8A7F76}
\colorlet{popgrey}{popgray}
% Alias tolérants (le modèle nomme parfois les couleurs différemment)
\colorlet{popwhite}{white}
\colorlet{popblack}{popink}
\colorlet{popred}{poppink}
\colorlet{popyellow}{popgold}
% Teintes claires (fonds de tableaux/encadrés) — le modèle nomme parfois
% les couleurs avec un suffixe L (ex. popblueL) sans les avoir définies.
\colorlet{poporangeL}{poporange!20}
\colorlet{popdarkL}{popdark!20}
\colorlet{poppurpleL}{poppurple!20}
\colorlet{popgreenL}{popgreen!20}
\colorlet{poppinkL}{poppink!20}
\colorlet{popblueL}{popblue!20}
\colorlet{popgoldL}{popgold!20}
\colorlet{popredL}{popred!20}
\colorlet{popgrayL}{popgray!20}
% Teintes claires pour fonds de tableaux / zones
\colorlet{popblueL}{popblue!10!white}
\colorlet{poporangeL}{poporange!14!white}
\colorlet{popgreenL}{popgreen!12!white}
\colorlet{poppurpleL}{poppurple!12!white}
\colorlet{poppinkL}{poppink!12!white}
\colorlet{popgoldL}{popgold!14!white}

\pagecolor{popcream}

% ============================================================
% Typographie
% ============================================================
\defaultfontfeatures{Ligatures=TeX}
\setmainfont{PlusJakartaSans-Medium.ttf}[Path=../../../fonts/,BoldFont=PlusJakartaSans-Bold.ttf]
% Maths en sans-serif (Fira Math) avec chiffres et lettres droites de Plus
% Jakarta, pour que les formules se fondent dans le texte.
\usepackage[math-style=ISO,bold-style=ISO]{unicode-math}
\setmathfont{FiraMath-Regular.otf}[Path=../../../fonts/]
\setmathfont{PlusJakartaSans-Medium.ttf}[Path=../../../fonts/,range={up/num,up/{Latin,latin},\mathup}]
\usepackage{icomma} % virgule décimale sans espace en mode math
% Caractères mathématiques Unicode absents des polices de texte (Plus Jakarta,
% Archivo Black) : rendus par leur équivalent mathématique, en texte comme en
% formule — sinon ils s'affichent en carré vide (ex. « Arithmétique dans ℤ »).
\usepackage{newunicodechar}
\newunicodechar{ℝ}{\ensuremath{\mathbb{R}}}
\newunicodechar{ℤ}{\ensuremath{\mathbb{Z}}}
\newunicodechar{ℂ}{\ensuremath{\mathbb{C}}}
\newunicodechar{ℕ}{\ensuremath{\mathbb{N}}}
\newunicodechar{ℚ}{\ensuremath{\mathbb{Q}}}
\newunicodechar{𝔻}{\ensuremath{\mathbb{D}}}
\newunicodechar{α}{\ensuremath{\alpha}}
\newunicodechar{β}{\ensuremath{\beta}}
\newunicodechar{γ}{\ensuremath{\gamma}}
\newunicodechar{δ}{\ensuremath{\delta}}
\newunicodechar{ε}{\ensuremath{\varepsilon}}
\newunicodechar{θ}{\ensuremath{\theta}}
\newunicodechar{λ}{\ensuremath{\lambda}}
\newunicodechar{μ}{\ensuremath{\mu}}
\newunicodechar{σ}{\ensuremath{\sigma}}
\newunicodechar{τ}{\ensuremath{\tau}}
\newunicodechar{φ}{\ensuremath{\varphi}}
\newunicodechar{ψ}{\ensuremath{\psi}}
\newunicodechar{ω}{\ensuremath{\omega}}
\newunicodechar{Σ}{\ensuremath{\Sigma}}
% majuscules produites par \MakeUppercase dans les titres (α-aminés -> Α)
\newunicodechar{Α}{\ensuremath{\alpha}}
\newunicodechar{Β}{\ensuremath{\beta}}
\newunicodechar{Γ}{\ensuremath{\Gamma}}
\newunicodechar{Δ}{\ensuremath{\Delta}}
\newunicodechar{Ε}{\ensuremath{\varepsilon}}
\newunicodechar{Θ}{\ensuremath{\Theta}}
\newunicodechar{Λ}{\ensuremath{\Lambda}}
\newunicodechar{Μ}{\ensuremath{\mu}}
\newunicodechar{Τ}{\ensuremath{\tau}}
\newunicodechar{Φ}{\ensuremath{\Phi}}
\newunicodechar{Ψ}{\ensuremath{\Psi}}
\newunicodechar{Ω}{\ensuremath{\Omega}}
\newunicodechar{↦}{\ensuremath{\mapsto}}
\newunicodechar{∈}{\ensuremath{\in}}
\newunicodechar{∉}{\ensuremath{\notin}}
\newunicodechar{∧}{\ensuremath{\wedge}}
\newunicodechar{⇔}{\ensuremath{\Leftrightarrow}}
\newunicodechar{⟺}{\ensuremath{\Longleftrightarrow}}
\newunicodechar{⇒}{\ensuremath{\Rightarrow}}
\newunicodechar{⟹}{\ensuremath{\Longrightarrow}}
\newunicodechar{∘}{\ensuremath{\circ}}
\newunicodechar{∪}{\ensuremath{\cup}}
\newunicodechar{∩}{\ensuremath{\cap}}
\newunicodechar{≡}{\ensuremath{\equiv}}
\newunicodechar{⊂}{\ensuremath{\subset}}
\newunicodechar{⊥}{\ensuremath{\perp}}
\newunicodechar{∃}{\ensuremath{\exists}}
\newunicodechar{∀}{\ensuremath{\forall}}
\newunicodechar{∛}{\ensuremath{\sqrt[3]{\,}}}
\newunicodechar{∜}{\ensuremath{\sqrt[4]{\,}}}
\newunicodechar{⃗}{\ensuremath{{}^{\to}}}
\newunicodechar{ⁿ}{\ifmmode^{n}\else\textsuperscript{\ensuremath{n}}\fi}
\newunicodechar{ˣ}{\ifmmode^{x}\else\textsuperscript{\ensuremath{x}}\fi}
\newunicodechar{ᵇ}{\ifmmode^{b}\else\textsuperscript{\ensuremath{b}}\fi}
\newunicodechar{ʳ}{\ifmmode^{r}\else\textsuperscript{\ensuremath{r}}\fi}
\newunicodechar{ᵘ}{\ifmmode^{u}\else\textsuperscript{\ensuremath{u}}\fi}
\newunicodechar{ʸ}{\ifmmode^{y}\else\textsuperscript{\ensuremath{y}}\fi}
\newunicodechar{ᵃ}{\ifmmode^{a}\else\textsuperscript{\ensuremath{a}}\fi}
\newunicodechar{ᵉ}{\ifmmode^{e}\else\textsuperscript{\ensuremath{e}}\fi}
\newunicodechar{ᵏ}{\ifmmode^{k}\else\textsuperscript{\ensuremath{k}}\fi}
\newunicodechar{ᵖ}{\ifmmode^{p}\else\textsuperscript{\ensuremath{p}}\fi}
\newunicodechar{ᴺ}{\ifmmode^{N}\else\textsuperscript{\ensuremath{N}}\fi}
\newunicodechar{ᶜ}{\ifmmode^{c}\else\textsuperscript{\ensuremath{c}}\fi}
\newunicodechar{ʰ}{\ifmmode^{h}\else\textsuperscript{\ensuremath{h}}\fi}
\newunicodechar{ᵠ}{\ifmmode^{q}\else\textsuperscript{\ensuremath{q}}\fi}
\newunicodechar{ₙ}{\ifmmode_{n}\else\textsubscript{\ensuremath{n}}\fi}
\newunicodechar{ₐ}{\ifmmode_{a}\else\textsubscript{\ensuremath{a}}\fi}
\newunicodechar{ᵢ}{\ifmmode_{i}\else\textsubscript{\ensuremath{i}}\fi}
\newunicodechar{ᵣ}{\ifmmode_{r}\else\textsubscript{\ensuremath{r}}\fi}
\newunicodechar{ₖ}{\ifmmode_{k}\else\textsubscript{\ensuremath{k}}\fi}
\newunicodechar{ᵦ}{\ifmmode_{\beta}\else\textsubscript{\ensuremath{\beta}}\fi}
\newunicodechar{ₓ}{\ifmmode_{x}\else\textsubscript{\ensuremath{x}}\fi}
\newfontfamily\popdisplay{ArchivoBlack-Regular.ttf}[Path=../../../fonts/]
\newfontfamily\poplogo{SpaceGrotesk-Bold.ttf}[Path=../../../fonts/]
\newfontfamily\popsemi{PlusJakartaSans-SemiBold.ttf}[Path=../../../fonts/]
\newfontfamily\popxbold{PlusJakartaSans-ExtraBold.ttf}[Path=../../../fonts/]
\newfontfamily\popxboldls{PlusJakartaSans-ExtraBold.ttf}[Path=../../../fonts/,LetterSpace=3.0]

\setlist[enumerate]{leftmargin=1.7em,itemsep=5pt,topsep=4pt}
\setlist[itemize]{leftmargin=1.7em,itemsep=4pt,topsep=4pt,label={\color{poporange}\textbullet}}
\setlength{\parindent}{0pt}
\setlength{\parskip}{5pt}
\renewcommand{\arraystretch}{1.25}

% pgfplots : style Pop par défaut
\pgfplotsset{
  every axis/.append style={axis line style={popink,line width=1pt},tick style={popink},
    grid style={popline,line width=0.5pt},label style={font=\small},tick label style={font=\footnotesize}},
  popcurve/.style={poporange,line width=1.8pt,no markers,smooth},
}
% Même style disponible dans un \draw TikZ simple (hors pgfplots)
\tikzset{popcurve/.style={poporange,line width=1.8pt}}
\tikzset{popgrid/.style={popline,very thin}}
\colorlet{popcurve}{poporange} % le nom du style est parfois utilisé comme couleur

% ============================================================
% Pill / squiggle
% ============================================================
\newcommand{\poppill}[3]{%
  \tikz[baseline=-0.55ex]\node[draw=popink,line width=1.2pt,rounded corners=2.6pt,%
    fill=#2,inner xsep=6.5pt,inner ysep=3pt]{%
    \popxboldls\fontsize{7}{7}\selectfont\color{#3}\MakeUppercase{#1}};%
}
\newcommand{\popsquiggle}[1]{%
  \tikz[baseline=-0.4ex]{
    \draw[#1,line width=2.6pt,line cap=round]
      plot [smooth] coordinates {(0,0.09) (0.34,0.01) (0.68,0.21) (1.02,0.01) (1.36,0.10)};
  }%
}
% Mot-clé surligné (marqueur orange sous le texte)
\newcommand{\cle}[1]{{\popxbold\color{popdark}#1}}

% ============================================================
% En-tête / pied
% ============================================================
\pagestyle{fancy}
\fancyhf{}
\renewcommand{\headrulewidth}{0pt}
\renewcommand{\footrulewidth}{0pt}
\lhead{\raisebox{-0.15ex}{\poplogo\fontsize{13}{13}\selectfont\color{popink}OnBuch\color{poporange}+}}
\rhead{\poppill{\DOCMATIERE\ \textbullet\ \DOCNIVEAU\ \textbullet\ Leçon \DOCLECON}{white}{popink}}
\lfoot{\popxbold\fontsize{7.6}{7.6}\selectfont\color{popmuted}\MakeUppercase{Cours signé OnBuch+}\ \textbullet\ {\popsemi\DOCTITRE}}
\rfoot{\tikz[baseline=-0.6ex]\node[rounded corners=8.5pt,fill=popink,minimum height=17pt,minimum width=17pt,inner xsep=5pt,inner sep=0]{\popxbold\fontsize{8}{8}\selectfont\color{popcream}\thepage\,/\,\pageref*{LastPage}};}

% ============================================================
% Titres
% ============================================================
\newcommand{\chaptitre}[2]{%
  \begin{tcolorbox}[enhanced,colback=poporange,colframe=popink,boxrule=2.6pt,arc=11pt,
      drop shadow=popink,left=15pt,right=15pt,top=12pt,bottom=11pt,before skip=3pt,after skip=18pt]
    \poppill{#2}{popink}{popcream}\par\vspace{9pt}
    {\raggedright\hyphenpenalty=10000\exhyphenpenalty=10000\popdisplay\fontsize{22}{24}\selectfont\color{popcream}\MakeUppercase{#1}\par}
  \end{tcolorbox}
}
% Section numérotée : pastille orange avec le numéro + titre Archivo
\newcounter{popsec}
\newcommand{\coursec}[1]{%
  \stepcounter{popsec}\setcounter{popsub}{0}%
  \par\vspace{16pt}\noindent
  \begin{minipage}{\linewidth}
  \tikz[baseline=-0.7ex]\node[draw=popink,line width=1.6pt,fill=poporange,rounded corners=4pt,minimum size=22pt,inner sep=0]{\popdisplay\fontsize{12}{12}\selectfont\color{popcream}\thepopsec};%
  \hspace{8pt}{\popdisplay\fontsize{15}{17}\selectfont\color{popink}\MakeUppercase{#1}}\par
  \vspace{4pt}\noindent{\color{poporange}\rule{36pt}{3.6pt}}
  \end{minipage}\par\nopagebreak\vspace{8pt}
}
\newcounter{popsub}
\newcommand{\courssub}[1]{%
  \stepcounter{popsub}%
  \par\vspace{9pt}\noindent{\popxbold\fontsize{11.5}{13}\selectfont\color{popink}{\color{poporange}\thepopsec.\thepopsub}\hspace{6pt}#1}\par\nopagebreak\vspace{4pt}
}

% ============================================================
% Boîtes pédagogiques — même recette : fond blanc, bordure épaisse colorée,
% ombre dure encre, badge plein. Titre optionnel [#1].
% ============================================================
\tcbset{popbase/.style={breakable,enhanced,colback=white,boxrule=2.3pt,arc=9pt,
  shadow={3pt}{-3pt}{0pt}{popink},left=14pt,right=14pt,top=11pt,bottom=11pt,before skip=11pt,after skip=15pt,
  fontupper=\popsemi\fontsize{10.6}{14.5}\selectfont\color{popink},
  fontlower=\popsemi\fontsize{10.6}{14.5}\selectfont\color{popink}}}
% Coche dessinée (le glyphe ✓ n'existe pas dans les polices du document)
\newcommand{\popcheck}[1]{\tikz[baseline=-0.5ex]\draw[#1,line width=1.6pt,line cap=round,line join=round](-0.13,0.0)--(-0.04,-0.09)--(0.13,0.1);}
\providecommand{\coche}{\popcheck{popgreen}}
\providecommand{\resultat}[1]{\par\smallskip\noindent\fcolorbox{popgreen}{popgreenL}{\parbox{\dimexpr\linewidth-2\fboxsep-2\fboxrule\relax}{\textbf{Résultat.} #1}}\par\smallskip}
\newcommand{\popboxhead}[3]{\poppill{#1}{#2}{white}\ifx\relax#3\relax\else\hspace{6pt}{\popxbold\fontsize{10.6}{12}\selectfont\color{popink}#3}\fi\par\vspace{7pt}}

\newtcolorbox{aretenir}[1][]{popbase,colframe=popgreen,before upper={\popboxhead{À retenir}{popgreen}{#1}}}
\newtcolorbox{exemplebox}[1][]{popbase,colframe=poporange,before upper={\popboxhead{Exemple}{poporange}{#1}}}
\newtcolorbox{definition}[1][]{popbase,colframe=popblue,before upper={\popboxhead{Définition}{popblue}{#1}}}
\newtcolorbox{propriete}[1][]{popbase,colframe=poppurple,before upper={\popboxhead{Propriété}{poppurple}{#1}}}
\newtcolorbox{methode}[1][]{popbase,colframe=popgold,before upper={\popboxhead{Méthode}{popgold}{#1}}}
\newtcolorbox{attention}[1][]{popbase,colframe=poppink,before upper={\popboxhead{Attention}{poppink}{#1}}}
\newtcolorbox{experience}[1][]{popbase,colframe=popblue,colback=popblueL,before upper={\popboxhead{Expérience}{popblue}{#1}}}
% « Le savais-tu ? » : carte violette pleine (contexte, culture, Cameroun)
\newtcolorbox{savaistu}[1][]{popbase,colback=poppurple,colframe=popink,
  fontupper=\popsemi\fontsize{10.4}{14.2}\selectfont\color{white},
  before upper={\poppill{Le savais-tu\,?}{white}{poppurple}\ifx\relax#1\relax\else\hspace{6pt}{\popxbold\color{white}#1}\fi\par\vspace{7pt}}}
% Exercice résolu : énoncé en haut, \tcblower puis la solution sur fond crème
\newcounter{popexo}\newcounter{popexr}
\newtcolorbox{exoresolu}[1][]{popbase,colframe=poporange,colbacklower=popcream,
  segmentation style={popink,line width=1.2pt,dashed},
  before upper={\stepcounter{popexr}\popboxhead{Exercice résolu \thepopexr}{poporange}{#1}},
  before lower={\poppill{Solution}{popink}{popcream}\par\vspace{6pt}}}
% Exercice (non résolu en place) — la correction est dans la partie Corrigés
\newtcolorbox{exercice}[1][]{popbase,colframe=popink,
  before upper={\stepcounter{popexo}\popboxhead{Exercice \thepopexo}{popink}{#1}}}
% Corrigé
\newtcolorbox{corrige}[1][]{popbase,colframe=popgreen,colback=popgreenL,
  before upper={\popboxhead{Corrigé}{popgreen}{#1}}}
% Objectifs / prérequis / bilan
\newtcolorbox{objectifs}{popbase,colframe=popink,colback=poporangeL,
  before upper={\popboxhead{Objectifs de la leçon}{popink}{}}}
\newtcolorbox{prerequis}{popbase,colframe=popink,colback=popblueL,
  before upper={\popboxhead{Prérequis}{popblue}{}}}
\newtcolorbox{bilan}{popbase,colframe=popink,colback=white,boxrule=2.8pt,
  before upper={\popboxhead{Fiche bilan}{poporange}{L'essentiel à connaître pour l'examen}}}
% Figure encadrée avec légende
\newtcolorbox{popfigure}[1][]{enhanced,breakable=false,colback=white,colframe=popline,boxrule=1.4pt,arc=8pt,
  left=10pt,right=10pt,top=10pt,bottom=8pt,before skip=10pt,after skip=12pt,halign=center,
  lower separated=false,halign lower=center,
  fontlower=\popsemi\fontsize{9}{11.5}\selectfont\color{popmuted},
  #1}
\newcommand{\legende}[1]{\par\vspace{6pt}{\popsemi\fontsize{9}{11.5}\selectfont\color{popmuted}#1}}

% ============================================================
% Signature OnBuch+
% ============================================================
\newcommand{\poplogoinline}[1]{{\poplogo\fontsize{#1}{#1}\selectfont\color{popink}OnBuch\color{poporange}+}}
\newcommand{\signatureonbuch}{%
  \begin{tcolorbox}[enhanced,colback=popink,colframe=popink,boxrule=2.2pt,arc=10pt,
      drop shadow=poporange,left=14pt,right=14pt,top=10pt,bottom=10pt,before skip=0pt,after skip=0pt]
    \tikz[baseline=-0.7ex]\node[circle,fill=popgreen,draw=popcream,line width=1.3pt,minimum size=22pt,inner sep=0]{\popcheck{white}};%
    \hspace{9pt}\begin{minipage}[c]{0.62\linewidth}
      {\popxbold\fontsize{12}{13}\selectfont\color{popcream}Signé\ \textbullet\ {\poplogo OnBuch\color{poporange}+}}\par\vspace{2pt}
      {\raggedright\popsemi\fontsize{8}{10.5}\selectfont\color{popline}\MakeUppercase{Équipe pédagogique OnBuch+}\\\MakeUppercase{Conforme au programme officiel MINESEC}\par}
    \end{minipage}\hfill
    {\poplogo\fontsize{22}{22}\selectfont\color{popcream}OnBuch\color{poporange}+}
  \end{tcolorbox}%
}

% ============================================================
% Page de garde
% #1 titre · #2 matière · #3 niveau · #4 module · #5 n° leçon · #6 durée ·
% #7 description / objectifs (texte ou itemize)
% ============================================================
\newcommand{\pagegarde}[7]{%
  \thispagestyle{empty}
  \newgeometry{a4paper,margin=1.7cm,top=1.6cm,bottom=1.4cm}%
  \begin{tikzpicture}[remember picture,overlay]
    \fill[poporange!18] ([xshift=-3.2cm,yshift=-3.6cm]current page.north east) circle (4.6cm);
    \fill[poppurple!14] ([xshift=2.4cm,yshift=7.6cm]current page.south west) circle (3.8cm);
  \end{tikzpicture}%
  {\poplogo\fontsize{30}{30}\selectfont\color{popink}OnBuch\color{poporange}+}\hfill
  \raisebox{0.6ex}{\poppill{Cours signé \textbullet\ Premium}{popink}{popcream}}\par
  \vspace{1pt}\popsquiggle{poppurple}\par
  \vspace{8pt}
  \begin{tcolorbox}[enhanced,colback=poporange,colframe=popink,boxrule=2.8pt,arc=13pt,
      drop shadow=popink,left=18pt,right=18pt,top=12pt,bottom=13pt,after skip=10pt]
    \poppill{#2 \textbullet\ #3}{popink}{popcream}\hspace{5pt}\poppill{#4}{white}{popink}\par\vspace{8pt}
    {\popdisplay\fontsize{14}{14}\selectfont\color{popink}LEÇON #5}\par\vspace{4pt}
    {\raggedright\hyphenpenalty=10000\exhyphenpenalty=10000\popdisplay\fontsize{25}{27}\selectfont\color{popcream}\MakeUppercase{#1}\par}
    \vspace{8pt}
    {\popxbold\fontsize{9.5}{9.5}\selectfont\color{popink}\MakeUppercase{Durée conseillée : #6}}
  \end{tcolorbox}
  \begin{tcolorbox}[enhanced,colback=white,colframe=popink,boxrule=2.2pt,arc=10pt,
      drop shadow=popink,left=16pt,right=16pt,top=10pt,bottom=10pt,after skip=8pt]
    \poppill{Dans cette leçon}{poppurple}{white}\par\vspace{5pt}
    \popsemi\fontsize{9.3}{12.6}\selectfont\color{popink}#7
  \end{tcolorbox}
  \noindent\begin{minipage}[t]{0.487\linewidth}
    \begin{tcolorbox}[enhanced,colback=popgreen,colframe=popink,boxrule=2.2pt,arc=10pt,
        drop shadow=popink,left=11pt,right=11pt,top=10pt,bottom=10pt,height=3.1cm,valign=center]
      \tcbox[colback=white,colframe=popink,boxrule=1.3pt,arc=5pt,boxsep=0pt,nobeforeafter,
        left=3pt,right=3pt,top=3pt,bottom=3pt,valign=center]{\qrcode[height=1.9cm]{https://play.google.com/store/apps/details?id=com.luvvix.onbuch&hl=fr}}%
      \hspace{8pt}\begin{minipage}[c]{0.5\linewidth}
        {\popdisplay\fontsize{11}{12.5}\selectfont\color{white}\MakeUppercase{Télécharge OnBuch}}\par\vspace{4pt}
        {\raggedright\popsemi\fontsize{8.4}{11}\selectfont\color{white}Gratuit \textbullet\ Google Play\\Tuteur IA, annales, concours, résultats\par}
      \end{minipage}
    \end{tcolorbox}
  \end{minipage}\hfill
  \begin{minipage}[t]{0.487\linewidth}
    \begin{tcolorbox}[enhanced,colback=poppurple,colframe=popink,boxrule=2.2pt,arc=10pt,
        drop shadow=popink,left=11pt,right=11pt,top=10pt,bottom=10pt,height=3.1cm,valign=center]
      \tikz[baseline=-0.7ex]\node[circle,fill=white,draw=popink,line width=1.3pt,minimum size=1.6cm,inner sep=0]{\popdisplay\fontsize{24}{24}\selectfont\color{poppurple}+};%
      \hspace{8pt}\begin{minipage}[c]{0.6\linewidth}
        {\popdisplay\fontsize{11}{12.5}\selectfont\color{white}\MakeUppercase{Passe à OnBuch+}}\par\vspace{4pt}
        {\raggedright\popsemi\fontsize{8.4}{11}\selectfont\color{white}Tous les cours signés, Tuteur IA illimité, fiches PDF — onglet OnBuch+ de l'app\par}
      \end{minipage}
    \end{tcolorbox}
  \end{minipage}
  \vspace{8pt}
  \popcontenu
  \vfill
  \signatureonbuch
  \restoregeometry
  \newpage
}

% Rangée « Ce cours contient » (page de garde)
\newcommand{\poptile}[3]{% #1 couleur #2 gros texte #3 légende
  \begin{tcolorbox}[enhanced,colback=#1,colframe=popink,boxrule=2pt,arc=9pt,shadow={2.5pt}{-2.5pt}{0pt}{popink},
      left=6pt,right=6pt,top=9pt,bottom=9pt,halign=center,valign=center,height=1.75cm,width=\linewidth,nobeforeafter]
    {\popdisplay\fontsize{12.5}{13}\selectfont\color{white}\MakeUppercase{#2}}\par\vspace{3pt}
    {\centering\popsemi\fontsize{7.8}{9.5}\selectfont\color{white}#3\par}
  \end{tcolorbox}}
\newcommand{\popcontenu}{%
  \par\noindent{\popxboldls\fontsize{7.5}{7.5}\selectfont\color{popmuted}CE COURS SIGNÉ CONTIENT}\par\vspace{7pt}
  \noindent\begin{minipage}[t]{0.235\linewidth}\poptile{popblue}{Cours}{complet, illustré, pas à pas}\end{minipage}\hfill
  \begin{minipage}[t]{0.235\linewidth}\poptile{poporange}{Méthodes}{et exercices résolus}\end{minipage}\hfill
  \begin{minipage}[t]{0.235\linewidth}\poptile{poppink}{Exercices}{type examen + corrigés}\end{minipage}\hfill
  \begin{minipage}[t]{0.235\linewidth}\poptile{popgreen}{Fiche bilan}{l'essentiel à retenir}\end{minipage}\par}

% Sommaire
\newtcolorbox{sommaire}{enhanced,colback=white,colframe=popink,boxrule=2.2pt,arc=10pt,
  drop shadow=popink,left=18pt,right=18pt,top=13pt,bottom=13pt,before skip=6pt,after skip=20pt,
  before upper={\poppill{Sommaire}{poppurple}{white}\par\vspace{9pt}\popsemi\fontsize{10.6}{14.5}\selectfont\color{popink}}}

% Signature de fin de cours — poussée tout en bas de la dernière page
\newcommand{\findecours}{%
  \par\vfill\nopagebreak
  \begin{center}\popsquiggle{poporange}\end{center}\vspace{4pt}
  \signatureonbuch
}
\AtBeginDocument{\def\llbracket{\mathopen{[\mkern-3.5mu[}}\def\rrbracket{\mathclose{]\mkern-3.5mu]}}} % intervalles d'entiers (glyphe ⟦ absent de la police)
% Glyphes absents de Fira Math : redéfinis à partir de glyphes disponibles
\AtBeginDocument{%
  \def\setminus{\mathbin{\backslash}}\let\smallsetminus\setminus
  \def\vdots{\mathord{\vbox{\baselineskip3.2pt\lineskiplimit0pt\kern4pt\hbox{.}\hbox{.}\hbox{.}}}}%
  \def\ddots{\mathinner{\mkern1mu\raise6.5pt\hbox{.}\mkern2mu\raise3.6pt\hbox{.}\mkern2mu\raise0.7pt\hbox{.}\mkern1mu}}%
  \def\triangle{\mathord{\scalebox{0.9}{$\symup{\Delta}$}}}%
  \def\nabla{\mathord{\raisebox{\depth}{\scalebox{1}[-1]{$\symup{\Delta}$}}}}%
  \def\checkmark{\popcheck{popdark}}%
  \def\star{\mathbin{\ast}}\def\bigstar{\mathord{\ast}}%
  \def\diamond{\mathbin{\scalebox{0.6}{$\Diamond$}}}%
  \def\bigcup{\mathop{\mathchoice{\raisebox{-0.3em}{\scalebox{1.7}{$\cup$}}}{\scalebox{1.25}{$\cup$}}{\cup}{\cup}}\displaylimits}%
  \def\bigcap{\mathop{\mathchoice{\raisebox{-0.3em}{\scalebox{1.7}{$\cap$}}}{\scalebox{1.25}{$\cap$}}{\cap}{\cap}}\displaylimits}%
  \def\lesseqqgtr{\mathrel{\lessgtr}}\def\bigoplus{\mathop{\oplus}\displaylimits}%
}
\providecommand{\ssi}{si et seulement si} % abréviation fréquente
\AtBeginDocument{\providecommand{\wideparen}{\overparen}} % arc de cercle

%% FIGFIX-BEGIN v4 — correctifs globaux des figures (docs/onbuch-generation/tools/figfix)
\usepackage{adjustbox}\usepackage{etoolbox}
% toute figure TikZ/pgfplots plus large que la ligne est réduite à la largeur disponible
\BeforeBeginEnvironment{tikzpicture}{\begin{adjustbox}{max width=\linewidth}}
\AfterEndEnvironment{tikzpicture}{\end{adjustbox}}
% légendes pgfplots lisibles par-dessus les courbes
\pgfplotsset{every axis/.append style={legend style={fill=white,fill opacity=0.92,text opacity=1,draw=black!15,inner sep=3pt}}}
% étiquettes d'arêtes (syntaxe edge["..."]) avec fond blanc
\tikzset{every edge quotes/.append style={fill=white,inner sep=1.2pt}}
%% FIGFIX-END
\begin{document}

\pagegarde{Les villes}{Géographie}{1re Tech}{Géographie – classe de Première technique}{N03}{4 séances (fiche de progression harmonisée)}{Cette leçon étudie la croissance urbaine au Cameroun, les problèmes qu'elle engendre et les solutions envisageables. Elle analyse ensuite les relations ville-campagne et le rôle de l'agriculture et de l'élevage dans l'approvisionnement des villes et le développement du pays.
\vspace{4pt}
\begin{multicols}{2}\small
\begin{itemize}
\item À la fin de cette leçon, je sais définir les notions de ville, d'urbanisation et de croissance urbaine
\item À la fin de cette leçon, je sais identifier et expliquer les principaux problèmes urbains au Cameroun à partir de documents
\item À la fin de cette leçon, je sais proposer des solutions aux problèmes urbains en justifiant ma démarche
\item À la fin de cette leçon, je sais décrire les relations d'échanges et de dépendance entre villes et campagnes
\item À la fin de cette leçon, je sais distinguer les types d'agriculture pratiqués au Cameroun et leurs caractéristiques
\item À la fin de cette leçon, je sais présenter les formes d'élevage au Cameroun et leurs difficultés
\end{itemize}\end{multicols}}

\begin{sommaire}
\begin{multicols}{2}
\begin{enumerate}
\item La ville et l'urbanisation au Cameroun
\item Les problèmes urbains
\item Dossier 2 : Les relations ville-campagne
\item L'agriculture au Cameroun
\item L'élevage au Cameroun
\item Synthèse : villes, campagnes et développement
\item Activité d'intégration
\item Méthodes et exercices résolus
\item Exercices
\item Corrigés des exercices
\item Fiche bilan
\end{enumerate}
\end{multicols}
\end{sommaire}

\begin{objectifs}
\begin{itemize}
\item À la fin de cette leçon, je sais définir les notions de ville, d'urbanisation et de croissance urbaine
\item À la fin de cette leçon, je sais identifier et expliquer les principaux problèmes urbains au Cameroun à partir de documents
\item À la fin de cette leçon, je sais proposer des solutions aux problèmes urbains en justifiant ma démarche
\item À la fin de cette leçon, je sais décrire les relations d'échanges et de dépendance entre villes et campagnes
\item À la fin de cette leçon, je sais distinguer les types d'agriculture pratiqués au Cameroun et leurs caractéristiques
\item À la fin de cette leçon, je sais présenter les formes d'élevage au Cameroun et leurs difficultés
\item À la fin de cette leçon, je sais construire un croquis, un diagramme ou une carte simple à partir de données chiffrées
\item À la fin de cette leçon, je sais appliquer ces notions à des situations concrètes de la vie courante et rédiger une conclusion argumentée
\end{itemize}
\end{objectifs}

\begin{prerequis}
\begin{itemize}
\item Notion de population et de répartition spatiale (classe de 4e)
\item Notion de milieu naturel et d'activités humaines au Cameroun
\item Lecture simple d'une carte et d'un tableau statistique
\item Notion de migration et d'exode rural (vie courante et cours de 2de)
\end{itemize}
\end{prerequis}

\newpage

\coursec{La ville et l'urbanisation au Cameroun}

Chaque matin, des milliers de jeunes quittent les villages du Noun ou de la Menoua pour rejoindre Douala et Yaoundé, villes qui comptent désormais plus de 3 millions d'habitants chacune. Embouteillages interminables à Akwa, quartiers précaires comme Nkolbisson ou Makepe-Missoke, insécurité alimentaire alors que les campagnes environnantes regorgent de produits vivriers : la ville camerounaise fascine autant qu'elle inquiète. Comment expliquer cette croissance urbaine rapide ? Quels problèmes pose-t-elle ? Et comment les villes et les campagnes peuvent-elles coopérer au lieu de s'opposer ? C'est à ces questions que cette leçon apporte des réponses.

\courssub{Définitions et concepts fondamentaux}

\begin{definition}[Ville, urbanisation et taux d'urbanisation]
\textbf{Ville} : agglomération qui se distingue par une population nombreuse (généralement $\ge$ 5\,000 habitants au Cameroun), une densité de bâti élevée, une prédominance des activités secondaires et tertiaires, et la présence d'équipements collectifs (administration, santé, éducation, marchés, transports).

\textbf{Urbanisation} : processus d'augmentation de la part de la population urbaine dans la population totale d'un pays. Elle résulte de la croissance naturelle des villes, de l'exode rural et de la création de nouveaux centres urbains (érigement en chefs-lieux).

\textbf{Taux d'urbanisation} : pourcentage de la population urbaine par rapport à la population totale du pays à une date donnée.
\[
\text{Taux d'urbanisation (\%)} = \frac{\text{Population urbaine}}{\text{Population totale}} \times 100
\]

\textbf{Croissance urbaine} : augmentation absolue du nombre d'habitants dans les villes sur une période.

\textbf{Métropole} : ville de rang supérieur (souvent $>$ 1 million d'habitants) qui exerce des fonctions de commande à l'échelle nationale ou régionale (siège d'institutions, pôles économiques, nœuds de transport, universités).
\end{definition}

\begin{propriete}[Critères de définition d'une ville au Cameroun]
L'Institut national de la statistique (INS) retient trois critères cumulatifs pour qualifier une localité de « urbaine » :
\begin{enumerate}
\item \textbf{Seuil démographique} : population $\ge$ 5\,000 habitants (recensement).
\item \textbf{Critère fonctionnel} : prédominance des activités non agricoles (commerce, industrie, services, administration).
\item \textbf{Critère d'équipement} : existence d'infrastructures urbaines (voirie, eau, électricité, établissements scolaires et sanitaires, marché structuré, siège d'une commune d'arrondissement ou de district).
\end{enumerate}
Une localité qui ne remplit qu'un ou deux de ces critères reste classée « rurale » même si sa population dépasse 5\,000 habitants (ex. : gros villages agricoles du Nord).
\end{propriete}

\courssub{L'évolution de l'urbanisation au Cameroun : 1960–2020}

Le Cameroun a connu une transition urbaine spectaculaire en six décennies. En 1960, à l'indépendance, moins de 15 \% des Camerounais vivaient en ville. En 2020, ce taux dépasse 57 \%, soit plus de 14 millions d'urbains sur une population totale d'environ 26 millions.

\begin{popfigure}
\begin{tikzpicture}
\begin{axis}[
    ybar,
    bar width=12pt,
    width=\linewidth,
    height=8cm,
    enlarge x limits=0.25,
    ylabel={Taux d'urbanisation (\%)},
    xlabel={Année},
    symbolic x coords={1960,1970,1980,1990,2000,2010,2020},
    xtick=data,
    nodes near coords,
    nodes near coords align={vertical},
    ymin=0,ymax=70,
    ymajorgrids=true,
    grid style=dashed,
    tick label style={font=\small},
    label style={font=\small},
    title style={font=\normalsize},
    title={Évolution du taux d'urbanisation au Cameroun (1960–2020)},
    popblue,
    fill=popblue!60,
]
\addplot coordinates {
    (1960,14.5)
    (1970,20.2)
    (1980,32.1)
    (1990,40.5)
    (2000,46.8)
    (2010,52.3)
    (2020,57.7)
};
\end{axis}
\end{tikzpicture}
\legende{Source : INS, Recensements généraux de la population et de l'habitat (RGPH) 1976, 1987, 2005 ; projections ONU / World Urbanization Prospects 2018.}
\end{popfigure}

\begin{exemplebox}[Lecture du graphique]
Entre 1960 et 2020, le taux d'urbanisation a été multiplié par \textbf{4}. La pente est la plus forte entre 1970 et 1990 (période de boom pétrolier et d'extension administrative), puis elle s'infléchit légèrement après 2000, tout en restant soutenue. Cela traduit une urbanisation rapide mais dont le rythme tend à se stabiliser à mesure que le pays approche du seuil des 60 \%.
\end{exemplebox}

\begin{aretenir}[Causes de la croissance urbaine au Cameroun]
Trois moteurs principaux, souvent imbriqués :
\begin{enumerate}
\item \textbf{Exode rural} : départ massif des jeunes (15–35 ans) des zones agricoles vers les villes, poussés par la pauvreté rurale, le sous-emploi, l'insécurité foncière, l'attrait des salaires urbains et des services (santé, éducation). C'est la cause historique majeure (années 1970–1990).
\item \textbf{Accroissement naturel urbain} : taux de natalité encore élevé en ville (indice de fécondité $\approx$ 4,0 en milieu urbain contre 5,5 en milieu rural en 2018) et baisse de la mortalité grâce aux services de santé. Aujourd'hui, ce facteur contribue pour \textbf{plus de 60 \%} à la croissance urbaine annuelle.
\item \textbf{Création de chefs-lieux} : découpage administratif (nouvelles régions, départements, arrondissements) qui transforme des villages en centres urbains artificiels, attirant fonctionnaires et commerçants (ex. : création de la région du Nord-Ouest en 1972, découpage de 2008).
\end{enumerate}
\end{aretenir}

\begin{attention}[Ne pas confondre croissance naturelle et exode rural]
\textbf{Erreur fréquente} : attribuer toute la croissance des villes à l'arrivée de migrants ruraux. En réalité, depuis les années 2000, \textbf{l'accroissement naturel (naissances – décès) à l'intérieur même des villes} est le premier moteur. L'exode rural alimente encore les bidonvilles périphériques, mais la « ville qui fait des enfants » pèse plus lourd dans les chiffres globaux. Au Probatoire, précisez toujours la part respective des deux facteurs si des données sont fournies.
\end{attention}

\courssub{Les principales villes du Cameroun et leurs fonctions}

Le réseau urbain camerounais est fortement \textbf{macrocéphale} : deux métropoles (Douala, Yaoundé) concentrent près de 40 \% de la population urbaine totale, suivies d'un deuxième rang de capitales régionales (Garoua, Maroua, Bamenda, Bafoussam, Ngaoundéré, Bertoua, Ebolowa) et de villes moyennes.

\begin{popfigure}
\begin{tikzpicture}[scale=0.85]
% Contour très simplifié du Cameroun (polygone approximatif)
\draw[popdark, thick, fill=popcream!30]
(0,0) -- (1.2,0.2) -- (2.5,0.5) -- (3.8,0.3) -- (5.2,0.8) -- (6.5,1.5) -- (7.2,2.8) -- (7.5,4.2) -- (7.0,5.5) -- (6.0,6.8) -- (4.8,7.5) -- (3.5,7.8) -- (2.2,7.5) -- (1.0,6.8) -- (0.2,5.5) -- (0,4.0) -- (0.2,2.5) -- (0.5,1.2) -- cycle;

% Villes principales (coordonnées approximatives relatives)
% Douala
\fill[poporange] (1.3,2.0) circle (3pt);
\node[above right, font=\tiny, popdark] at (1.3,2.0) {Douala};
% Yaoundé
\fill[poporange] (3.2,3.2) circle (3pt);
\node[above right, font=\tiny, popdark] at (3.2,3.2) {Yaoundé};
% Garoua
\fill[popblue] (4.8,6.2) circle (2.5pt);
\node[above, font=\tiny, popdark] at (4.8,6.2) {Garoua};
% Maroua
\fill[popblue] (5.8,6.8) circle (2.5pt);
\node[above right, font=\tiny, popdark] at (5.8,6.8) {Maroua};
% Ngaoundéré
\fill[popblue] (5.5,4.8) circle (2.5pt);
\node[below right, font=\tiny, popdark] at (5.5,4.8) {Ngaoundéré};
% Bertoua
\fill[popblue] (5.8,3.5) circle (2.5pt);
\node[below right, font=\tiny, popdark] at (5.8,3.5) {Bertoua};
% Bafoussam
\fill[popblue] (2.0,4.5) circle (2.5pt);
\node[left, font=\tiny, popdark] at (2.0,4.5) {Bafoussam};
% Bamenda
\fill[popblue] (1.5,5.5) circle (2.5pt);
\node[left, font=\tiny, popdark] at (1.5,5.5) {Bamenda};
% Ebolowa
\fill[popblue] (2.8,1.8) circle (2.5pt);
\node[below left, font=\tiny, popdark] at (2.8,1.8) {Ebolowa};
% Kribi (port en développement)
\fill[popgreen] (1.0,1.0) circle (2pt);
\node[below left, font=\tiny, popdark] at (1.0,1.0) {Kribi};

% Flèche du nord
\draw[->, thick, popdark] (7.5,1.0) -- (7.5,1.8);
\node[right, font=\small, popdark] at (7.5,1.8) {N};

% Échelle indicative
\draw[thick, popdark] (0.2,0.3) -- (1.2,0.3);
\node[below, font=\tiny, popdark] at (0.7,0.25) {~200 km};

% Légende manuelle
\node[anchor=south west, font=\tiny, popdark, align=center] at (0.1,7.2){
\begin{tabular}{ll}
\multicolumn{2}{l}{\textbf{Légende :}} \\
\color{poporange}\rule{4pt}{4pt} & Métropoles ($>$ 3 M hab.) \\
\color{popblue}\rule{4pt}{4pt} & Capitales régionales (200–500 k hab.) \\
\color{popgreen}\rule{4pt}{4pt} & Ville portuaire émergente
\end{tabular}
};
\end{tikzpicture}
\legende{Carte schématique du Cameroun : localisation des 10 principales villes (contours et positions approximatifs).}
\end{popfigure}

\begin{exemplebox}[Douala et Yaoundé : deux métropoles de plus de 3 millions d'habitants]
\begin{itemize}
\item \textbf{Douala} (région du Littoral) : capitale économique, premier port d'Afrique centrale (12 millions de tonnes/an), poumon industriel (raffinerie SONARA à Limbe proche, brasseries, textiles, cimenteries), nœud autoroutier et ferroviaire (Transcam). Population estimée 2023 : \textbf{3,9 millions} (aire urbaine).
\item \textbf{Yaoundé} (région du Centre) : capitale politique et administrative (Présidence, Assemblée, ministères, ambassades), pôle universitaire (Université de Yaoundé I et II, grandes écoles), fonctions de services supérieurs. Population estimée 2023 : \textbf{3,6 millions} (aire urbaine).
\end{itemize}
Ces deux villes forment un \textbf{duopole} qui structure l'espace national : l'axe Douala–Yaoundé (300 km) concentre 60 \% du PIB industriel et 50 \% de la population urbaine.
\end{exemplebox}

\begin{propriete}[Fonctions urbaines principales]
Chaque ville assure une ou plusieurs fonctions dominantes qui justifient sa place dans le réseau :
\begin{center}
\begin{tabularx}{\linewidth}{|l|X|X|}\hline
\textbf{Fonction} & \textbf{Description} & \textbf{Villes emblématiques} \\ \hline
Administrative & Siège d'institutions (État, région, département), services publics & Yaoundé (nationale), toutes les capitales régionales \\ \hline
Commerciale & Marchés de gros, import-export, banques, grands magasins & Douala (Akwa, Bassa), Yaoundé (Mokolo, Mfoundi) \\ \hline
Industrielle & Usines, zones industrielles, transformation matières premières & Douala (Bassa, Bonabéri), Édéa (aluminium), Limbe (pétrole) \\ \hline
Universitaire / Recherche & Universités, grandes écoles, centres de recherche & Yaoundé, Ngaoundéré, Dschang, Buea, Maroua \\ \hline
Touristique / Balnéaire & Plages, parcs nationaux, patrimoine culturel & Kribi, Limbe, Foumban, Maroua (parc Waza) \\ \hline
Religieuse / Culturelle & Lieux de pèlerinage, chefferies traditionnelles & Foumban (sultanat), Mokolo, Garoua \\ \hline
\end{tabularx}
\end{center}
La plupart des villes sont \textbf{multifonctionnelles} (ex. : Bafoussam = administrative + commerciale + universitaire), mais une fonction domine souvent.
\end{propriete}

\courssub{Méthode : Calculer un taux d'urbanisation à partir de données de population}

\begin{methode}[Calcul du taux d'urbanisation]
\emph{Objectif} : déterminer le pourcentage d'urbains dans la population totale d'un pays ou d'une région à une date donnée.

\begin{enumerate}
\item \textbf{Identifier les données} : population urbaine ($P_u$) et population totale ($P_t$) pour la même année (source : RGPH, INS, ONU).
\item \textbf{Appliquer la formule} :
      \[
      T_u (\%) = \frac{P_u}{P_t} \times 100
      \]
\item \textbf{Exprimer le résultat} avec une décimale (ex. : 57,7 \%) et l'interpréter : « En 2020, 57,7 \% des Camerounais vivaient en ville, contre 14,5 \% en 1960. »
\item \textbf{Calculer l'évolution} (taux de variation) entre deux dates :
      \[
      \Delta T_u = T_{u,2} - T_{u,1} \quad \text{(en points de pourcentage)}
      \]
      ou taux de croissance annuel moyen :
      \[
      \text{TCAM} = \left( \frac{T_{u,2}}{T_{u,1}} \right)^{\frac{1}{n}} - 1
      \]
\end{enumerate}
\end{methode}

\begin{exoresolu}[Calcul du taux d'urbanisation 2020]
\textbf{Énoncé} : Selon le RGPH 2005 et les projections 2020, le Cameroun compte 26,5 millions d'habitants dont 15,3 millions d'urbains. Calculer le taux d'urbanisation 2020 et son augmentation en points depuis 2005 (taux 2005 = 46,8 \%).

\textbf{Solution détaillée}
\begin{align*}
P_t &= 26,5 \text{ millions} \\
P_u &= 15,3 \text{ millions} \\[4pt]
T_{u,2020} &= \frac{15,3}{26,5} \times 100 = 57,7358... \approx \textbf{57,7 \%} \\[4pt]
\Delta T_u &= 57,7 - 46,8 = \textbf{10,9 points de pourcentage en 15 ans} \\
\text{TCAM} &= \left( \frac{57,7}{46,8} \right)^{\frac{1}{15}} - 1 \approx 0,0136 = \textbf{1,36 \% par an}
\end{align*}
\textbf{Interprétation} : Le taux d'urbanisation a gagné près de 11 points en 15 ans, soit une progression moyenne de 1,36 \% par an. La transition urbaine se poursuit mais ralentit par rapport à la période 1976–1987 (TCAM $\approx$ 2,5 \%).
\tcblower
\textbf{Point méthode} : Toujours vérifier l'unité (millions) et arrondir au dixième de pourcentage. Ne pas confondre « points de pourcentage » (différence) et « pour-cent » (taux relatif).
\end{exoresolu}

\courssub{Repères historiques et pièges à éviter}

\begin{savaistu}[Yaoundé, capitale depuis 1922]
Yaoundé n'a pas toujours été la capitale. De 1884 à 1901, les Allemands administrent le Kamerun depuis \textbf{Douala} (puis Buea brièvement). En 1901, le gouverneur von Puttkamer déplace le siège à \textbf{Yaoundé} (Jaunde), plus sain (altitude 700 m, climat tempéré) et central. La France confirme ce choix en 1922 (mandat SDN) et l'État indépendant le maintient en 1960. Ce statut de \textbf{ville-capitale} pendant un siècle explique la concentration inédite de fonctions administratives, diplomatiques et de services supérieurs, moteur principal de sa croissance bien plus que l'industrie (quasi absente).
\end{savaistu}

\begin{attention}[Pièges classiques au Probatoire]
\begin{itemize}
\item \textbf{Confondre « urbanisation » et « urbanisme »} : l'urbanisation est un processus démographique ; l'urbanisme est l'art d'aménager les villes.
\item \textbf{Oublier le critère fonctionnel} : un village de 6\,000 habitants resté agricole n'est pas une ville.
\item \textbf{Mélanger « taux d'urbanisation » et « taux de croissance urbaine »} : le premier est une part (\%), le second est une variation annuelle (\%/an).
\item \textbf{Croire que Douala = capitale} : Douala est la capitale \emph{économique}, Yaoundé la capitale \emph{politique/administrative}.
\item \textbf{Négliger l'accroissement naturel} : depuis 2000, il dépasse l'exode rural comme facteur de croissance urbaine.
\end{itemize}
\end{attention}

\courssub{Synthèse de la section}

\begin{aretenir}[Essentiel à retenir sur la ville et l'urbanisation au Cameroun]
\begin{itemize}
\item \textbf{Ville} = population $\ge$ 5\,000 + fonctions non agricoles + équipements urbains.
\item \textbf{Taux d'urbanisation} = $\frac{P_{urbaine}}{P_{totale}} \times 100$ ; passé de \textbf{14,5 \% (1960)} à \textbf{57,7 \% (2020)}.
\item \textbf{Croissance urbaine} tirée aujourd'hui principalement par l'\textbf{accroissement naturel} (60 \%) puis par l'\textbf{exode rural} et la \textbf{création de chefs-lieux}.
\item \textbf{Réseau urbain macrocéphale} : Douala (économique, port, industrie) et Yaoundé (politique, administrative, universitaire) = 2 métropoles $>$ 3,5 M hab. chacune.
\item \textbf{Fonctions urbaines} : administrative, commerciale, industrielle, universitaire, touristique, religieuse — souvent combinées.
\end{itemize}
\end{aretenir}

\begin{exercice}[Application directe]
\begin{enumerate}
\item Définir : (a) urbanisation, (b) taux d'urbanisation, (c) métropole.
\item Citer trois critères qui permettent de définir une ville au Cameroun.
\item À partir des données ci-dessous, calculer le taux d'urbanisation du Cameroun en 2010 et son augmentation en points par rapport à 1987.
      \begin{center}
      \begin{tabular}{|l|c|c|}\hline
      Année & Population totale (millions) & Population urbaine (millions) \\ \hline
      1987 & 10,5 & 3,4 \\
      2010 & 19,6 & 10,3 \\ \hline
      \end{tabular}
      \end{center}
\item Expliquer pourquoi l'accroissement naturel est devenu le premier moteur de la croissance urbaine depuis 2000.
\item Sur une carte muette du Cameroun, localiser et nommer : Douala, Yaoundé, Garoua, Maroua, Ngaoundéré, Bafoussam, Bamenda, Bertoua, Ebolowa, Kribi.
\end{enumerate}
\end{exercice}

\begin{corrige}[Exercice n]
\begin{enumerate}
\item (a) Urbanisation : processus d'augmentation de la part de la population urbaine dans la population totale. (b) Taux d'urbanisation : $\%$ d'urbains dans la population totale. (c) Métropole : ville de rang supérieur ($>$ 1 M hab.) exerçant des fonctions de commande nationale/régionale.
\item Seuil démographique ($\ge$ 5\,000 hab.), prédominance activités non agricoles, équipements urbains (voirie, eau, électricité, santé, éducation, marché, siège administratif).
\item 1987 : $T_u = \frac{3,4}{10,5}\times100 = 32,4\%$. 2010 : $T_u = \frac{10,3}{19,6}\times100 = 52,6\%$. Augmentation = $52,6 - 32,4 = 20,2$ points en 23 ans.
\item Baisse de la mortalité urbaine (santé), fécondité encore élevée en ville (IF $\approx$ 4,0), base de population urbaine déjà large (effet d'inertie démographique). L'exode rural continue mais son flux relatif diminue.
\item Voir carte schématique en début de section (figure TikZ).
\end{enumerate}
\end{corrige}

\coursec{Les problèmes urbains}

La section précédente a montré que l'urbanisation au Cameroun est rapide, souvent non planifiée et concentrée sur quelques métropoles, notamment Douala et Yaoundé. Cette croissance démographique fulgurante, couplée à une incapacité des infrastructures à suivre le rythme, engendre une multitude de dysfonctionnements. Nous analysons ici les principaux problèmes qui minent le quotidien des citadins camerounais, leurs causes structurelles et les réponses apportées par les pouvoirs publics.

\courssub{Le chômage et l'explosion du secteur informel}

L'arrivée massive de jeunes sur le marché du travail, bien plus rapide que la création d'emplois formels, génère un chômage structurel élevé, touchant particulièrement les diplômés de l'enseignement technique et supérieur. Face à l'absence de filet de protection sociale, la survie passe par le \cle{secteur informel} : activités non déclarées, sans contrat de travail, sans couverture sociale, à faible productivité et à revenu aléatoire.

\begin{definition}[Secteur informel]
Ensemble des unités de production de biens et de services dont l'objectif principal est la création d'emplois et de revenus pour les personnes concernées, qui fonctionnent à un faible niveau d'organisation, avec une petite échelle d'opérations, un capital et des compétences faibles, et qui ne sont pas enregistrées auprès des administrations publiques (impôts, sécurité sociale, statistiques).
\end{definition}

Au Cameroun, ce secteur absorbe plus de \textbf{90\% de la population active occupée} (source : INS, EESI). Deux figures emblématiques illustrent cette réalité :
\begin{itemize}
    \item Le \cle{bayam-sellam} (acheter-revendre) : commerce de rue, étals sur les trottoirs, vente à la sauvette aux carrefours. Faible marge, forte concurrence, harcèlement municipal.
    \item Les \cle{motos-taxis} (\emph{benskin}) : réponse à l'insuffisance des transports en commun. Source de revenus pour des milliers de jeunes, mais vecteur d'insécurité routière (absence de casque, excès de vitesse) et de conflits avec les forces de l'ordre.
\end{itemize}

\begin{exemplebox}[Bayam-sellam et motos-taxis à Douala]
À Douala, on estime à plus de 100 000 le nombre de conducteurs de motos-taxis. Pour un revenu journalier moyen de 3 000 à 5 000 FCFA, ils doivent payer la location de la moto (environ 2 000 FCFA/jour), l'essence, les pots-de-vin aux contrôles routiers et les amendes. Le bayam-sellam, quant à lui, occupe les abords des marchés (Nkololoun, Congo, Mfoundi) et les artères principales, transformant les trottoirs en étals et compliquant la circulation piétonne et automobile.
\end{exemplebox}

\begin{attention}[Piège : l'explication monocausale]
Ne pas réduire le secteur informel à une simple « paresse » ou « manque d'esprit d'entreprise ». C'est une \textbf{réponse rationnelle} à l'absence d'alternatives formelles. Les causes sont systémiques : inadéquation formation-emploi, fiscalité étouffante pour les PME, coût du foncier, lenteur administrative. Toute analyse doit articuler conjoncture économique, politiques publiques et stratégies d'acteurs.
\end{attention}

\courssub{L'habitat précaire et les quartiers spontanés}

L'incapacité de l'État à produire du logement social abordable et la spéculation foncière poussent les populations à faible revenu vers l'auto-construction sur des terrains marginaux, souvent insalubres ou à risque.

\begin{definition}[Quartier spontané]
Agglomération d'habitations construites sans titre foncier, sans permis de construire, en dehors de tout plan d'urbanisme, sur des terrains occupés illégalement (zones marécageuses, flancs de collines, emprises routières/ferroviaires). Caractérisé par l'absence de voirie, de réseaux d'eau/électricité/assainissement, et une forte densité de population.
\end{definition}

Deux cas d'école à Yaoundé et Douala :
\begin{itemize}
    \item \textbf{Makepe-Missoke (Douala)} : Situé en zone marécageuse au bord du Wouri, sujet aux inondations chroniques. Habitat en matériaux de récupération (tôles, planches), promiscuité extrême, absence totale d'assainissement. Choléra endémique.
    \item \textbf{Nkolbikok (Yaoundé)} : Quartier de pente instable, risques d'éboulements. Réseau hydrographique détourné par les constructions, aggravant l'érosion. Accès difficile pour les engins de secours et de collecte des ordures.
\end{itemize}

\begin{methode}[Analyser une photographie de quartier précaire]
\begin{enumerate}
    \item \textbf{Décrire} : Type d'habitat (matériaux, hauteur, densité), état de la voirie (pistes, caniveaux), présence/absence de réseaux (poteaux électriques, bornes-fontaines), activités visibles (commerce, artisanat, lessive).
    \item \textbf{Localiser} : Indices de topographie (plat, pente, fond de vallée), proximité d'infrastructures (route goudronnée, voie ferrée, zone industrielle), végétation résiduelle (marécage, forêt).
    \item \textbf{Expliquer} : Pourquoi ici ? (Foncier libre/peu cher, proximité emploi informel, tolérance administrative). Quels risques ? (Inondation, éboulement, incendie, épidémie). Quelles traces d'action publique ? (Bornes-fontaines CAMWATER, poteaux ENEO, passage d'engins HYSACAM).
\end{enumerate}
\end{methode}

\begin{exemplebox}[Makepe-Missoke : le piège du foncier marécageux]
Makepe-Missoke s'est développé sur des terres gagnées sur la mangrove par remblaiement sauvage (sable, déchets). Lors des grandes marées ou pluies intenses, l'eau remonte par capillarité et par les caniveaux bouchés. Les latrines à ciel ouvert déversent dans le même milieu qui sert parfois de source d'eau pour la lessive. Le taux de prévalence du paludisme et des diarrhées y est significativement plus élevé que dans les quartiers planifiés (Bastos, Bonapriso).
\end{exemplebox}

\courssub{Les problèmes de transport : embouteillages, insuffisance des voies, accidents}

La mobilité urbaine est l'un des goulets d'étranglement majeurs. Le réseau routier, conçu pour des villes de 500 000 habitants, en supporte aujourd'hui plus de 3 à 4 millions.
\begin{itemize}
    \item \textbf{Embouteillages quotidiens} : Perte de temps économique énorme (2 à 3 h/jour pour certains trajets), surconsommation de carburant, stress, pollution atmosphérique.
    \item \textbf{Insuffisance des voies} : Absence de rocades, de voies rapides structurantes, de transports en commun de masse (pas de métro, tramway, BHNS opérationnel à grande échelle). Le réseau est radioconcentrique, tout converge vers le centre, créant des nœuds de congestion (Carrefour Warda, Rond-point Deido, Carrefour Mvog-Ada).
    \item \textbf{Accidents} : Vétusté du parc automobile (âge moyen > 15 ans), non-respect du code de la route, cohabitation forcée voitures/motos/piétons sur chaussées étroites sans trottoirs. Le Cameroun compte l'un des taux de mortalité routière les plus élevés d'Afrique (environ 27 morts pour 100 000 hab. selon l'OMS).
\end{itemize}

\begin{savaistu}[Le deuxième pont sur le Wouri, inauguré en 2018]
Long de 756 m, ce pont à haubans (financement Exim Bank Chine / État camerounais, coût ~130 milliards FCFA) double la capacité de franchissement du Wouri entre Bonabéri (zone industrielle/portuaire) et le centre-ville de Douala (Akwa, Bonanjo). Il a réduit le temps de trajet aux heures de pointe de 45-60 min à 10-15 min. Cependant, il a déplacé les bouchons vers les échangeurs d'entrée/sortie (Bonabéri, Deido), prouvant qu'une infrastructure ponctuelle ne résout pas un problème systémique sans plan de transport global.
\end{savaistu}

\courssub{La pollution : déchets, eaux usées, air et fleuve Wouri}

La gestion des déchets est le visage le plus visible de la crise urbaine.

\begin{exemplebox}[Douala : plus de 1 500 tonnes de déchets par jour]
La ville produit environ \textbf{1 500 à 2 000 tonnes de déchets ménagers par jour} (estimation HYSACAM / CUD). La société HYSACAM (Hygiène et Salubrité du Cameroun), concessionnaire, en collecte environ 60-70\% dans les quartiers formels. Le reste (30-40\%) finit dans les caniveaux, les cours d'eau, les terrains vagues ou est brûlé à l'air libre. Coût de la collecte : ~15 000 FCFA/tonne, financé par la taxe d'enlèvement d'ordures (TEOM) sur la facture d'électricité, au taux de recouvrement faible.
\end{exemplebox}

\begin{itemize}
    \item \textbf{Eaux usées} : Quasi absence de stations d'épuration. Les fosses septiques (quand elles existent) vidangées dans la nature ou le réseau pluvial. Rejet direct dans le Wouri et la Mfoundi. Eutrophisation, mortalité piscicole, risque sanitaire (choléra, typhoïde).
    \item \textbf{Pollution de l'air} : Parc automobile vétuste (carburant frelaté, huile moteur), brûlage des ordures, poussières des pistes non revêtues, industries (cimenteries, raffinerie SONARA à Limbé, usines de Bonabéri). Pic de PM2.5 et PM10 aux heures de pointe.
    \item \textbf{Pollution du Wouri} : Métaux lourds (plomb, cadmium), hydrocarbures, charge bactérienne colossale. Impact sur la pêche artisanale, la mangrove (frayère) et la qualité de l'eau de mer aux plages de Kribi/Limbé par drift côtier.
\end{itemize}

\begin{popfigure}
\begin{tikzpicture}
    % Pie chart data: Chômage 30, Habitat 25, Transport 20, Pollution 15, Eau/Elec 10
    \def\total{100}
    \def\startangle{90}
    \foreach \label/\value/\color in {
        Chômage/30/popblue,
        Habitat précaire/25/poporange,
        Transport/20/popgreen,
        Pollution/15/poppurple,
        Eau \& Électricité/10/poppink} {
        \pgfmathsetmacro{\endangle}{\startangle - \value/\total*360}
        \draw[fill=\color, thick] (0,0) -- (\startangle:3cm) arc (\startangle:\endangle:3cm) -- cycle;
        \pgfmathsetmacro{\midangle}{(\startangle + \endangle)/2}
        \node at (\midangle:3.8cm) {\textbf{\label} (\value\%)};
        \xdef\startangle{\endangle}
    }
    \node[circle, fill=white, draw=popdark, thick, minimum size=2cm, align=center] at (0,0) {Problèmes\\urbains};
\end{tikzpicture}
\legende{Répartition simulée des problèmes urbains cités par les habitants d'une métropole camerounaise (enquête type Afrobaromètre / INS). Le chômage et l'habitat concentrent plus de la moitié des doléances.}
\end{popfigure}

\courssub{L'insécurité, la délinquance juvénile et les tensions sociales}

La précarité économique, l'oisiveté forcée des jeunes non scolarisés/non insérés, la promiscuité et l'impunité relative alimentent l'insécurité.
\begin{itemize}
    \item \textbf{Feymen / Coupeurs de route} : Bandes organisées (souvent originaires des quartiers précaires) opérant aux abords des villes, sur les axes interurbains, parfois dans les quartiers la nuit. Racket, vol à main armée, agressions.
    \item \textbf{Délinquance juvénile} : Vols à l'arraché (téléphones), consommation de drogue (tramadol, cannabis, « kobolo »), conflits de territoire entre bandes rivales.
    \item \textbf{Tensions sociales} : Conflits locataires/propriétaires (loyers impayés, expulsions), conflits autochtones/allogènes (accès au foncier, représentation politique), grèves urbaines (transporteurs, commerçants, enseignants) paralysant l'activité.
\end{itemize}
La réponse sécuritaire (opérations « Épervier », « BIR » en ville) reste répressive et ne traite pas les causes profondes : exclusion sociale, échec scolaire, absence de perspectives.

\courssub{Les difficultés d'approvisionnement en eau potable et en électricité}

Paradoxe : le Cameroun est « château d'eau » de l'Afrique centrale (fleuves, nappes) et exportateur d'électricité (vers le Tchad, projet vers le Nigeria), mais les citadins subissent pénuries et délestages.
\begin{itemize}
    \item \textbf{Eau (CAMWATER)} : Réseau vétuste (fuites > 40\% d'eau non facturée), extension anarchique non raccordée, pompages sauvages, coupures récurrentes (pannes électriques aux stations de pompage, rationnement en saison sèche). Recours aux forages privés (coûteux, qualité non contrôlée), aux vendeurs d'eau (« push-push », bidons) au prix 10 à 20 fois plus cher que le réseau.
    \item \textbf{Électricité (ENEO/SONATREL)} : Déficit de production en saison sèche (hydroélectricité dépendante des pluies, centrales thermiques d'appoint coûteuses et polluantes). Délestages tournants programmés ou imprévus. Surtensions détruisant les appareils. Branchements illicites (« câbles volants ») causant incendies et électrocutions dans les quartiers précaires.
\end{itemize}

\courssub{Les solutions : planification urbaine, lotissements, grands travaux, décentralisation}

Face à ce tableau, les pouvoirs publics déploient des stratégies multiples, aux résultats inégaux.

\begin{itemize}
    \item \textbf{Planification urbaine} : Schémas Directeurs d'Aménagement Urbain (SDAU) pour Douala, Yaoundé, Bafoussam, Bamenda. Plans d'Occupation des Sols (POS). Objectif : anticiper l'extension, réserver des emprises (voiries, équipements, espaces verts), densifier le centre. Freins : manque de moyens pour l'exécution, non-respect par les acteurs (politiques, promoteurs), corruption foncière.
    \item \textbf{Lotissements et logement social} : Programme « 10 000 logements » (peu abouti), nouveaux lotissements (Nkolbisson, Nsimalen, Douala 5ème, Bepanda Omnisports). Problème : coût encore élevé pour la cible (foncier + viabilisation + construction), éloignement des bassins d'emploi, absence de transports dédiés → risque de « villes-dortoirs » vides le jour.
    \item \textbf{Grands travaux structurants} :
    \begin{itemize}
        \item Ponts : 2ème pont sur le Wouri (Douala), pont sur la Sanaga (Edéa), projets de 3ème pont (Douala), pont Kribi-Lolabé.
        \item Voiries : Programme d'asphaltage des voiries urbaines (PAVUR), réhabilitation axes structurants (Boulevard de la Liberté, Carrefour Warda-Échangeur).
        \item Équipements : Stades (Japoma, Olembé, Bafoussam, Garoua, Limbé) pour CAN 2021 → legs d'infrastructures sportives, hôtelières, routières, mais coût d'opportunité élevé et maintenance incertaine.
    \end{itemize}
    \item \textbf{Politique de décentralisation} : Loi de 2019, transfert de compétences (voirie, marchés, état civil, plan local d'urbanisme) et de ressources (centrées d'impôts locaux, dotation générale de décentralisation) aux Communes (Communes d'Arrondissement + Communes de Ville). Objectif : gestion de proximité, redevabilité, financement local du développement. Défi : capacité technique/financière des élus, tutelle administrative encore pesante.
\end{itemize}

\begin{popfigure}
\begin{tikzpicture}[
    node distance=1.2cm and 1.5cm,
    box/.style={draw=popdark, thick, rounded corners, fill=popcream, minimum width=3cm, minimum height=1.2cm, align=center, font=\small},
    arrow/.style={-Stealth, thick, popdark},
    circlebig/.style={draw=popblue, thick, circle, minimum size=#1}
]
% Central Cause
\node[box, fill=popblue!20, minimum width=4cm, align=center] (cause){Croissance urbaine rapide \\ non maîtrisée \\ (Exode rural, Natalité)};
% First Ring: Problems
\node[box, above left=of cause, align=center] (p1){Chômage \\ Secteur informel};
\node[box, above=of cause, align=center] (p2){Habitat précaire \\ Quartiers spontanés};
\node[box, above right=of cause, align=center] (p3){Transport \\ Embouteillages};
\node[box, right=of cause, align=center] (p4){Pollution \\ Déchets, Eau, Air};
\node[box, below right=of cause, align=center] (p5){Insécurité \\ Délinquance};
\node[box, below=of cause, align=center] (p6){Eau / Électricité \\ Pénuries récurrentes};
% Second Ring: Consequences
\node[box, fill=poporange!20, above left=of p1, align=center] (c1){Pauvreté \\ Exclusion sociale};
\node[box, fill=poporange!20, above=of p2, align=center] (c2){Maladies \\ Choléra, Paludisme};
\node[box, fill=poporange!20, above right=of p3, align=center] (c3){Perte de productivité \\ Compétitivité faible};
\node[box, fill=poporange!20, right=of p4, align=center] (c4){Dégradation \\ Environnementale};
\node[box, fill=poporange!20, below right=of p5, align=center] (c5){Instabilité \\ Tensions sociales};
\node[box, fill=poporange!20, below=of p6, align=center] (c6){Frein à l'investissement \\ Fuite des capitaux};
% Arrows Cause -> Problems
\foreach \p in {p1,p2,p3,p4,p5,p6} {
    \draw[arrow] (cause) -- (\p);
}
% Arrows Problems -> Consequences
\draw[arrow] (p1) -- (c1);
\draw[arrow] (p2) -- (c2);
\draw[arrow] (p3) -- (c3);
\draw[arrow] (p4) -- (c4);
\draw[arrow] (p5) -- (c5);
\draw[arrow] (p6) -- (c6);
% Cross links (problems interact)
\draw[arrow, popmuted, dashed, bend left=15] (p1) to (p2);
\draw[arrow, popmuted, dashed, bend left=15] (p2) to (p4);
\draw[arrow, popmuted, dashed, bend left=15] (p3) to (p4);
\draw[arrow, popmuted, dashed, bend left=15] (p5) to (p1);
% Legend boxes
\node[draw=popblue, thick, fill=popblue!10, rounded corners, anchor=north west, font=\footnotesize, align=left] at (current bounding box.north east) {
    \textbf{Légende :}\\
    \textcolor{popblue}{$\blacksquare$} Causes structurelles\\
    \textcolor{popdark}{$\blacksquare$} Problèmes urbains directs\\
    \textcolor{poporange}{$\blacksquare$} Conséquences systémiques\\
    \textcolor{popmuted}{---} Interactions entre problèmes
};
\end{tikzpicture}
\legende{Schéma systémique : causes structurelles → problèmes urbains observés → conséquences socio-économiques et environnementales. Les flèches pointillées indiquent les rétroactions (ex: habitat précaire aggrave la pollution).}
\end{popfigure}

\begin{exoresolu}[Analyse d'un problème urbain : la gestion des déchets à Douala]
\textbf{Énoncé :} À partir des données : production 1 800 t/j, collecte HYSACAM 1 200 t/j, population 3,5 millions d'habitants, budget annuel gestion déchets 25 milliards FCFA. Calculer la production par habitant/jour, le taux de collecte, le coût par tonne collectée. Proposer deux mesures pour améliorer le taux de collecte.

\textbf{Solution détaillée :}
\begin{enumerate}
    \item \textbf{Production par habitant/jour :} $1\,800\,000\,\text{kg} / 3\,500\,000\,\text{hab} \approx \mathbf{0,51\,\text{kg/hab/jour}}$. (Moyenne africaine ~0,5-0,8 kg, standard OCDE ~1,5 kg).
    \item \textbf{Taux de collecte :} $1\,200 / 1\,800 = \mathbf{66,7\%}$. Un tiers des déchets n'est pas collecté formellement.
    \item \textbf{Coût par tonne collectée :} Collecte annuelle $\approx 1\,200 \times 365 = 438\,000\,\text{t}$. Coût/t = $25 \times 10^9 / 438\,000 \approx \mathbf{57\,000\,\text{FCFA/t}}$ (coût réel souvent sous-estimé car n'incluant pas l'enfouissement/traitement final).
    \item \textbf{Mesures d'amélioration :}
    \begin{itemize}
        \item \textbf{Pré-collecte communautaire / tri à la source :} Organiser les jeunes des quartiers (GIC, associations) pour le ramassage porte-à-porte vers des points de transfert (bennes), rémunérés par la commune via la TEOM. Réduit le coût « dernier kilomètre » pour HYSACAM et crée de l'emploi.
        \item \textbf{Valorisation / Économie circulaire :} Installer des unités de compostage (matière organique ~60-70\% des déchets) et de tri recyclables (plastique, métal) au niveau des dépôts de transit (Pk 10, Nkolfoulou). Réduit le volume à enfouir, génère des revenus (compost agricole, matière première recyclage), allonge la durée de vie de la décharge.
    \end{itemize}
\end{enumerate}
\end{exoresolu}

\begin{aretenir}[Bilan pour le Probatoire : Les problèmes urbains]
\begin{itemize}
    \item \textbf{Chômage/Informel :} > 90\% pop. active occupée (bayam-sellam, benskin). Cause : inadéquation formation/emploi, fiscalité PME.
    \item \textbf{Habitat :} Quartiers spontanés (Makepe, Nkolbikok) : risque inondation/éboulement, insalubrité, absence réseaux.
    \item \textbf{Transport :} Réseau radioconcentrique saturé, 2-3h bouchons/jour, mortalité routière ~27/100 000 hab. 2ème pont Wouri (2018) : soulagement partiel.
    \item \textbf{Pollution :} Déchets 1500-2000 t/j (Douala), collecte 60-70\%. Eaux usées sans épuration → Wouri/Mfoundi. Air : PM2.5, poussières, industrie.
    \item \textbf{Eau/Élec :} Paradoxe « château d'eau » / exportateur élec. vs pénuries. CAMWATER (fuites >40\%), ENEO (délestages saison secs).
    \item \textbf{Solutions :} SDAU/POS (planification), Lotissements (accès difficile), Grands travaux (ponts, stades, PAVUR), Décentralisation 2019 (compétences/ressources insuffisantes).
\end{itemize}
\end{aretenir}

\begin{exercice}[Application : Diagnostic croisé]
À partir du schéma systémique (figure ci-dessus) et des données du cours, rédigez une réponse structurée (introduction, développement en 3 parties, conclusion) au sujet suivant :
\textbf{« Les problèmes urbains au Cameroun ne sont pas une fatalité, mais le résultat de choix de développement. » Discutez cette affirmation en vous appuyant sur l'analyse des causes, des manifestations et des solutions en cours.}
\end{exercice}

\begin{corrige}[Exercice n1]
\textbf{Introduction :} Accroche (urbanisation 3,7\%/an, 58\% urbains). Définition problème urbain. Thèse : problèmes structurels liés à des choix (modèle de croissance, priorités budgétaires, gouvernance foncière) mais solutions techniques connues, freinées par l'institutionnel.
\textbf{Développement :}
\begin{enumerate}
    \item \textbf{Causes structurelles : choix de développement.} Priorité aux infrastructures « visibles » (stades, ponts) vs infrastructures « invisibles » (assainissement, voiries secondaires, logement social). Politique foncière libérale non régulée → spéculation, quartiers spontanés. Modèle éducatif déconnecté du marché → chômage diplômés.
    \item \textbf{Manifestations systémiques :} Boucle de rétroaction : Habitat précaire (pente, marécage) $\rightarrow$ Inondations/Érosion $\rightarrow$ Pollution/Eaux usées $\rightarrow$ Maladies $\rightarrow$ Absentéisme/Travail informel précaire $\rightarrow$ Pauvreté $\rightarrow$ Incapacité à payer loyer/formel $\rightarrow$ Habitat précaire. Transport : absence de TCSP (choix bus vs tramway/métro) $\rightarrow$ Moto-taxi roi $\rightarrow$ Accidents/Embouteillages $\rightarrow$ Perte PIB.
    \item \textbf{Solutions : volontarisme vs contrainte.} SDAU/POS existent mais non opposables / non financés. Décentralisation : transfert compétences sans transfert ressources effectif (Dotation Générale de Décentralisation < 5\% budget national vs 15-20\% recommandé CGLU). Grands travaux : utiles mais « top-down », peu intégrés aux plans locaux.
\end{enumerate}
\textbf{Conclusion :} Non-fatalité technique (solutions connues : tri/compost, BHNS, lotissements viabilisés, police de l'urbanisme). Fatalité politique/institutionnelle tant que : gouvernance foncière opaque, budget communal insignifiant, planification déconnectée du financier, absence de reddition de comptes. Le Probatoire attend l'articulation \emph{diagnostic $\rightarrow$ causes $\rightarrow$ solutions $\rightarrow$ limites}.
\end{corrige}

\coursec{Dossier 2 : Les relations ville-campagne}

Les problèmes urbains étudiés précédemment (surpeuplement, habitat précaire, chômage) trouvent une partie de leur explication dans les échanges intenses qui lient la ville à la campagne. Comprendre ces relations permet d'identifier des leviers pour un développement territorial plus équilibré.

\courssub{L'exode rural : causes et conséquences}

\begin{definition}[Exode rural]
L'\cle{exode rural} désigne le déplacement massif et durable de populations depuis les zones rurales vers les villes, motivé par la recherche de meilleures conditions de vie.
\end{definition}

\textbf{Causes principales}
\begin{itemize}
\item \cle{Misère rurale} : faibles revenus agricoles, manque d'infrastructures (routes, électricité, santé, éducation), sols dégradés, aléas climatiques.
\item \cle{Attrait des lumières} : représentation idéalisée de la vie urbaine (emplois salariés, commerces, loisirs, modernité).
\item \cle{Recherche d'emploi} : concentration des activités industrielles, tertiaires et informelles dans les centres urbains.
\end{itemize}

\textbf{Conséquences}
\begin{itemize}
\item \textbf{Vieillissement des campagnes} : départ des jeunes actifs, maintien des personnes âgées, baisse de la main-d'œuvre agricole.
\item \textbf{Gonflement des villes} : croissance rapide, souvent non planifiée, des quartiers périphériques ; pression sur le logement, l'eau, l'assainissement, le transport.
\item \textbf{Perte de savoir-faire agricoles} et réduction de la production vivrière locale.
\end{itemize}

\begin{exemplebox}[Les bayam-sellam et l'approvisionnement de Douala]
Chaque semaine, plusieurs milliers de tonnes de vivriers (manioc, plantain, maïs, igname) quittent la région de l'Ouest (Bafoussam, Dschang, Foumban) vers les marchés de Douala. Ce flux est assuré par des centaines de commerçants et transporteurs appelés \emph{bayam-sellam} (de l'anglais \emph{buy and sell}, « acheter et vendre »), qui achètent les produits au village et les revendent en ville. Ce cas illustre l'ampleur du flux de produits agricoles vers la métropole économique et la dépendance alimentaire de la ville vis-à-vis de sa campagne environnante.
\end{exemplebox}

\courssub{Les flux entre ville et campagne}

Les relations ville-campagne s'expriment par trois grands types de flux qui s'entrecroisent en permanence.

\textbf{Flux de personnes}
\begin{itemize}
\item \textbf{Migrations quotidiennes} : navetteurs (travailleurs, élèves, commerçants) qui résident en périphérie rurale et travaillent en ville.
\item \textbf{Migrations saisonnières} : main-d'œuvre agricole qui rejoint les champs à la saison des pluies et revient en ville en saison sèche.
\item \textbf{Migrations définitives} : installation durable en ville (exode rural) ou, plus rarement, retour au village (retraite, investissement foncier).
\end{itemize}

\textbf{Flux de produits}
\begin{itemize}
\item \textbf{Vers la ville} : vivriers (manioc, plantain, maïs, riz, haricots), produits de rente (cacao, café, coton), produits forestiers, bétail sur pied.
\item \textbf{Vers la campagne} : produits manufacturés (textiles, quincaillerie, intrants agricoles, carburants), biens d'équipement (motopompes, motoculteurs), denrées alimentaires transformées (huile, sucre, conserves).
\end{itemize}

\textbf{Flux de capitaux et de services}
\begin{itemize}
\item \textbf{Transferts d'argent} : envois réguliers des citadins vers leurs familles restées au village (scolarité, santé, construction, cérémonies).
\item \textbf{Services concentrés en ville} : banques, hôpitaux de référence, lycées, universités, administrations, marchés de gros.
\item \textbf{Investissements ruraux} : achat de terres, création de petites exploitations, financement de groupements agricoles par des urbains d'origine rurale.
\end{itemize}

\begin{definition}[Aire d'influence urbaine]
L'\cle{aire d'influence} d'une ville est l'espace rural autour d'elle qui entretient des échanges quotidiens ou réguliers (travail, marché, services, approvisionnement) avec ce centre urbain. Elle se délimite par la portée des flux de personnes, de biens et de services.
\end{definition}

\begin{methode}[Construire un croquis de relations ville-campagne]
\begin{enumerate}
\item Tracer le contour simplifié de la ville (cercle ou polygone) et de la campagne environnante.
\item Identifier les principaux flux (personnes, produits, capitaux) à partir de données locales ou d'enquêtes.
\item Représenter chaque flux par une flèche dont l'épaisseur est proportionnelle à l'importance quantitative (tonnes, nombre de navetteurs, montant des transferts).
\item Utiliser un code couleur : \textcolor{popgreen}{vert} pour les produits agricoles vers la ville, \textcolor{poporange}{orange} pour les produits manufacturés vers la campagne, \textcolor{popblue}{bleu} pour les personnes, \textcolor{poppurple}{violet} pour les capitaux.
\item Ajouter une légende complète (échelle des flèches, code couleur, orientation nord).
\item Intituler le croquis : « Relations ville-campagne -- Aire d'influence de [Nom de la ville] ».
\end{enumerate}
\end{methode}

\begin{popfigure}
\begin{center}
\begin{tikzpicture}[scale=0.9, every node/.style={font=\small}]
% Ville
\node[draw, fill=popblue!20, rounded corners, minimum width=2.5cm, minimum height=1.2cm] (ville) at (0,0) {\textbf{VILLE}};
% Campagne
\node[draw, fill=popgreen!15, rounded corners, minimum width=5cm, minimum height=2.5cm] (campagne) at (0,-3.2) {\textbf{CAMPAGNE}};
% Flèche produits agricoles (vert)
\draw[->, line width=2.5pt, popgreen] (-0.8,-2.0) -- node[midway, right, align=left] {Vivriers\\ (manioc, plantain, maïs)} (-0.8,-0.65);
% Flèche produits manufacturés (orange)
\draw[->, line width=1.5pt, poporange] (0.8,-0.65) -- node[midway, left, align=right] {Manufacturés\\ (intrants, textiles)} (0.8,-2.0);
% Flèche personnes (bleu) - double flèche
\draw[<->, thick, popblue] (ville.east) -- node[midway, above] {Navetteurs} (campagne.east);
% Flèche capitaux (violet)
\draw[->, thick, poppurple, dashed] (ville.west) -- node[midway, above] {Transferts} (campagne.west);
% Légende intégrée
\node[anchor=north west, align=left, font=\footnotesize] at (3.2,0.6) {
\textcolor{popgreen}{\rule{0.5cm}{2.5pt}} Vivriers vers la ville\\[2pt]
\textcolor{poporange}{\rule{0.5cm}{2.5pt}} Manufacturés vers la campagne\\[2pt]
\textcolor{popblue}{\rule{0.5cm}{2.5pt}} Personnes (navetteurs)\\[2pt]
\textcolor{poppurple}{\rule{0.5cm}{2.5pt}} Capitaux (transferts d'argent)
};
\end{tikzpicture}
\end{center}
\legende{Schéma simplifié des flux échangés entre une ville et sa campagne environnante. L'épaisseur des flèches rend compte de l'importance relative des flux.}
\end{popfigure}

\begin{popfigure}
\begin{center}
\begin{tikzpicture}[scale=0.8, every node/.style={font=\small}]
% Zones d'influence (cercles concentriques)
\draw[popgreen!40, dashed, thick] (0,0) circle (4.5);
\draw[popgreen!60, dashed, thick] (0,0) circle (2.5);
% Routes principales
\draw[thick, popline] (0,0) -- (30:5.2) node[above right, align=center, font=\footnotesize] {Route vers\\ Douala};
\draw[thick, popline] (0,0) -- (150:5.2) node[above left, align=center, font=\footnotesize] {Route vers\\ Yaoundé};
\draw[thick, popline] (0,0) -- (-90:5.2) node[below, align=center, font=\footnotesize] {Route vers\\ Foumban};
% Villages / marchés
\node[draw, fill=popgreen!25, circle, minimum size=0.55cm, font=\footnotesize] (v1) at (30:2) {M1};
\node[draw, fill=popgreen!25, circle, minimum size=0.55cm, font=\footnotesize] (v2) at (150:2) {M2};
\node[draw, fill=popgreen!25, circle, minimum size=0.55cm, font=\footnotesize] (v3) at (-90:2) {M3};
\node[draw, fill=popgreen!25, circle, minimum size=0.55cm, font=\footnotesize] (v4) at (30:4) {M4};
\node[draw, fill=popgreen!25, circle, minimum size=0.55cm, font=\footnotesize] (v5) at (150:4) {M5};
\node[draw, fill=popgreen!25, circle, minimum size=0.55cm, font=\footnotesize] (v6) at (-90:4) {M6};
% Flèches flux vivriers (vert) depuis villages vers ville
\foreach \v in {v1,v2,v3,v4,v5,v6}
  \draw[->, thick, popgreen] (\v) -- (0,0);
% Ville centrale Bafoussam
\node[draw, fill=popblue!35, circle, minimum size=1.3cm, font=\footnotesize\bfseries] (baf) at (0,0) {Bafoussam};
% Légende
\node[anchor=north west, align=left, font=\footnotesize] at (-5.6,-5.3) {
\textcolor{popgreen}{\rule{0.5cm}{2.5pt}} Vivriers (manioc, plantain, maïs) vers Bafoussam\\
\textcolor{popline}{\rule{0.5cm}{2.5pt}} Routes principales\\
Cercles pointillés : limites approximatives de l'aire d'influence\\
(proche : environ 25 km ; élargie : environ 45 km)
};
% Nord
\draw[->, thick] (-5,4.2) -- ++(0,0.9) node[above] {N};
\end{tikzpicture}
\end{center}
\legende{Croquis guidé de l'aire d'influence de Bafoussam sur sa campagne. Les villages-marchés (M1 à M6) approvisionnent la ville en vivriers (flèches vertes) le long des routes principales. Les cercles pointillés matérialisent deux cercles d'influence : quotidien (environ 25 km) et hebdomadaire (environ 45 km). Croquis schématique, sans précision cartographique.}
\end{popfigure}

\courssub{Dépendance mutuelle et pistes d'équilibre}

La relation ville-campagne n'est pas à sens unique. La ville dépend de la campagne pour son approvisionnement alimentaire (vivriers, bétail, produits forestiers) et pour une main-d'œuvre peu coûteuse. Inversement, la campagne dépend de la ville pour l'accès aux services (santé, éducation, banques), aux marchés de vente, aux intrants agricoles et aux transferts financiers.

\begin{attention}[Relation à double sens]
Ne pas réduire la relation ville-campagne à une simple « exploitation » de la campagne par la ville. Les transferts d'argent des migrants urbains financent souvent la construction d'écoles, de cases, de forages et l'achat d'intrants dans les villages d'origine. Ignorer ce retour de capitaux conduit à des diagnostics erronés et à des politiques déséquilibrées.
\end{attention}

\begin{savaistu}[Transferts d'argent et développement villageois]
Au Cameroun, les envois de fonds des citadins vers leurs villages d'origine représentent des sommes considérables chaque année. Ces transferts financent une part importante des dépenses des familles rurales (scolarité, santé, construction) et permettent la réalisation de nombreux projets villageois (points d'eau, salles de classe, cases). Ils constituent un levier majeur, bien que souvent informel, de développement local.
\end{savaistu}

\textbf{Pistes pour rééquilibrer les relations}
\begin{itemize}
\item \cle{Développement rural intégré} : appui aux groupements agricoles, vulgarisation de variétés améliorées, formation technique.
\item \cle{Électrification rurale} : accès à l'énergie pour la transformation locale (mouture, séchage, conservation) et la création d'activités génératrices de revenus.
\item \cle{Routes rurales et pistes} : désenclavement des bassins de production pour réduire les coûts de transport et les pertes post-récolte.
\item \cle{Agropoles et pôles de transformation} : implantation d'unités de transformation (farine de manioc, chips de plantain, jus de fruits) à proximité des zones de production, créant de l'emploi non agricole en milieu rural.
\item \cle{Services de proximité} : décentralisation des services de santé (centres de santé intégrés), d'enseignement (lycées techniques agricoles) et financiers (caisses d'épargne villageoises).
\end{itemize}

\courssub{Exercice résolu}

\begin{exoresolu}[Analyse quantitative des flux vivriers vers Douala]
Énoncé : D'après une enquête de terrain, les camions des commerçants \emph{bayam-sellam} acheminent chaque semaine \num{12000} tonnes de vivriers de la région de l'Ouest vers Douala. La production vivrière totale de la région de l'Ouest est estimée à \num{1200000} tonnes par an. Quel est le volume annuel écoulé vers Douala ? Quel pourcentage de la production annuelle cela représente-t-il ?

\tcblower
Solution détaillée :
\begin{align*}
\text{Volume hebdomadaire} &= 12\,000\ \text{t} \\
\text{Volume annuel vers Douala} &= 12\,000 \times 52 = 624\,000\ \text{t} \\
\text{Part de la production} &= \frac{624\,000}{1\,200\,000} \times 100 = 52\ \% .
\end{align*}
Ainsi, d'après ces données, plus de la moitié de la production vivrière de l'Ouest serait écoulée chaque année sur le seul marché de Douala, ce qui confirme la forte dépendance alimentaire de la métropole économique vis-à-vis de sa campagne périphérique.
\end{exoresolu}

\courssub{Exercices d'application}

\begin{exercice}[Relations ville-campagne dans votre département]
\begin{enumerate}
\item Identifiez la ville principale de votre département et dessinez, sur une feuille A4, un croquis de son aire d'influence en suivant la méthode présentée ci-dessus (flèches proportionnelles, code couleur, légende).
\item À partir de données locales (statistiques agricoles, enquêtes auprès de commerçants), estimez le tonnage hebdomadaire de manioc et de plantain qui quitte votre zone rurale vers la ville la plus proche.
\item Calculez la part de cette production dans la production totale départementale (donnée fournie par la délégation départementale de l'Agriculture).
\item Discutez, en 150 mots maximum, les conséquences de ce flux sur l'équilibre alimentaire local et sur l'économie des ménages ruraux.
\end{enumerate}
\end{exercice}

\begin{corrige}[Exercice n°1 -- Relations ville-campagne dans votre département]
\begin{enumerate}
\item Le croquis doit montrer la ville au centre, les principaux axes routiers, les villages fournisseurs (flèches vertes) et les flux de produits manufacturés (flèches oranges). L'épaisseur des flèches reflète les tonnages ou le nombre de navetteurs. La légende précise l'échelle (ex. : 1 mm = 100 t/semaine) et le code couleur.
\item Exemple de calcul : si \num{3500} t de manioc et \num{2200} t de plantain partent chaque semaine, le total hebdomadaire est de $3\,500 + 2\,200 = 5\,700$ t.
\item Production annuelle départementale (ex. : \num{800000} t). Volume annuel vers la ville : $5\,700 \times 52 = 296\,400$ t. Part : $\dfrac{296\,400}{800\,000} \times 100 \approx 37\ \%$.
\item Réponse type : « Ce flux assure l'approvisionnement de la ville mais réduit la disponibilité locale, faisant monter les prix pour les ruraux non producteurs. Les revenus des vendeurs augmentent, pourtant la dépendance aux intrants achetés en ville (engrais, carburant) crée un cercle vicieux. Le développement de la transformation locale (farine, chips) permettrait de retenir une partie de la valeur ajoutée au village. »
\end{enumerate}
\end{corrige}

\coursec{L'agriculture au Cameroun}

Dans la section précédente, nous avons vu que les villes et les campagnes entretiennent des relations étroites : les campagnes nourrissent les villes, les villes fournissent aux campagnes des services et des débouchés. Mais qui produit cette nourriture, et comment ? C'est toute la question de l'\cle{agriculture}, activité fondamentale des campagnes camerounaises. Sans agriculture performante, pas de villes bien nourries, pas de devises pour l'État, pas de développement équilibré. Étudions donc la place de l'agriculture dans l'économie, ses deux grandes formes, sa répartition sur le territoire, ses difficultés et les efforts de modernisation en cours.

\courssub{Définition et place de l'agriculture dans l'économie camerounaise}

\begin{definition}[Agriculture]
L'\cle{agriculture} est l'ensemble des activités qui consistent à cultiver la terre pour produire des végétaux destinés à l'alimentation des hommes et des animaux ou à l'industrie (textiles, huiles, caoutchouc). Associée à l'élevage, elle constitue le \cle{secteur primaire} de l'économie.
\end{definition}

Au Cameroun, l'agriculture occupe une place centrale :

\begin{itemize}
\item elle emploie environ \textbf{40 \% de la population active}, ce qui en fait le premier pourvoyeur d'emplois du pays ;
\item elle contribue de manière importante au \cle{Produit Intérieur Brut} (PIB) et fournit une part essentielle des recettes d'exportation (cacao, café, coton, banane, caoutchouc, huile de palme) ;
\item elle assure l'essentiel de l'\cle{autosuffisance alimentaire} : le Cameroun importe relativement peu de produits vivriers par rapport à ses voisins, ce qui lui vaut le surnom de « grenier de l'Afrique centrale » ;
\item elle alimente les villes en denrées et les industries locales en matières premières (savonneries, huileries, chocolateries, filatures).
\end{itemize}

\begin{aretenir}[L'agriculture, pilier de l'économie]
Avec environ 40 \% de la population active et une contribution majeure au PIB et aux exportations, l'agriculture est la base de l'économie camerounaise et le lien concret entre la ville et la campagne.
\end{aretenir}

\courssub{L'agriculture vivrière traditionnelle}

\begin{definition}[Agriculture vivrière]
L'\cle{agriculture vivrière} est une agriculture dont la production est destinée en priorité à la consommation de la famille du producteur. Le surplus éventuel est vendu sur les marchés locaux pour acheter d'autres biens (sel, huile, savon, scolarité des enfants).
\end{definition}

\textbf{Les principales cultures vivrières} varient selon les régions : le \cle{manioc}, le \cle{plantain}, l'\cle{igname}, le \cle{macabo} et les légumes dominent dans le Sud forestier humide ; le \cle{maïs}, le \cle{mil}, le \cle{sorgho}, l'\cle{arachide} et le \cle{niébé} dominent dans les savanes de l'Ouest, de l'Adamaoua et du Grand Nord.

\textbf{Les techniques traditionnelles} sont simples et demandent beaucoup de main-d'œuvre :

\begin{enumerate}
\item \textbf{Le brûlis} : on défriche une parcelle de forêt ou de savane, on coupe la végétation, on la laisse sécher puis on la brûle. Les cendres fertilisent temporairement le sol.
\item \textbf{La culture à la houe} : le travail du sol se fait à la main, avec la houe et la machette, sans tracteur ni animal de trait.
\item \textbf{La jachère} : après deux ou trois ans de culture, le sol s'épuise ; on abandonne la parcelle plusieurs années pour qu'elle se régénère naturellement, et on défriche ailleurs.
\end{enumerate}

Ce système, appelé \cle{culture itinérante sur brûlis}, convient tant que les terres sont abondantes. Mais la croissance démographique réduit la durée des jachères, ce qui appauvrit les sols et fait baisser les rendements.

\courssub{L'agriculture de rente (les plantations)}

\begin{definition}[Agriculture de rente]
L'\cle{agriculture de rente} (ou agriculture de plantation, agriculture commerciale) est une agriculture dont la production est destinée à la vente, surtout à l'exportation. Elle rapporte des devises au pays et un revenu monétaire au planteur.
\end{definition}

Les principales cultures de rente camerounaises sont :

\begin{itemize}
\item le \cle{cacao} et le \cle{café}, cultivés par des centaines de milliers de petits planteurs du Centre, du Sud, de l'Ouest et du Littoral ;
\item la \cle{banane} d'exportation, produite dans la plaine du Littoral ;
\item le \cle{palmier à huile}, dans le Littoral et le Sud-Ouest ;
\item l'\cle{hévéa} (caoutchouc naturel), dans le Sud et le Littoral ;
\item le \cle{coton}, dans le Grand Nord (Nord, Adamaoua, Extrême-Nord).
\end{itemize}

De \textbf{grandes sociétés agro-industrielles} possèdent d'immenses plantations et des usines de transformation : la \textbf{CDC} (Cameroon Development Corporation, banane, palme, hévéa dans le Sud-Ouest), la \textbf{PHP} (Plantations du Haut Penja, banane au Littoral), la \textbf{SOSUCAM} (Société sucrière du Cameroun, canne à sucre à Nkoteng et Mbandjock), la \textbf{SODECOTON} (Société de développement du coton, dans le Grand Nord).

\begin{exemplebox}[Un champion mondial : le cacao camerounais]
Le Cameroun est le \textbf{5\textsuperscript{e} producteur mondial de cacao}, derrière la Côte d'Ivoire, le Ghana, l'Indonésie et l'Équateur. Sa production annuelle dépasse couramment \num{250000} tonnes. Le cacao est la première culture d'exportation agricole du pays et fait vivre des centaines de milliers de familles de planteurs du Centre-Sud.
\end{exemplebox}

\begin{attention}[Ne pas opposer vivrier et rente]
Erreur fréquente : présenter l'agriculture vivrière et l'agriculture de rente comme deux mondes séparés. En réalité, elles \textbf{coexistent souvent sur la même exploitation} : un planteur du Centre cultive du cacao pour la vente \emph{et} du manioc, du plantain et du macabo pour nourrir sa famille. Le vivrier assure la subsistance, la rente procure le revenu monétaire.
\end{attention}

\courssub{La répartition spatiale des grandes cultures}

La répartition des cultures obéit d'abord aux conditions du \cle{milieu naturel} (climat, sols), puis aux infrastructures (routes, ports) et à l'histoire (politique coloniale des plantations).

\begin{itemize}
\item \textbf{Centre, Sud, Est et Ouest} (zone forestière humide, deux saisons des pluies) : cacao, café robusta, palmier à huile, hévéa ;
\item \textbf{Hauts plateaux de l'Ouest} (altitude, climat tempéré) : café arabica, cultures maraîchères ;
\item \textbf{Littoral} (plaine très humide, proximité du port de Douala) : banane d'exportation, palme, hévéa ;
\item \textbf{Grand Nord} (savane sèche, une saison des pluies) : coton, arachide, mil, sorgho.
\end{itemize}

\begin{popfigure}
\begin{center}
\begin{tikzpicture}[scale=1.05]
% Contour simplifié du Cameroun
\draw[thick,popdark,fill=popgreenL] (2.2,5.6) -- (3.4,6.4) -- (3.1,7.6) -- (3.6,9.2) -- (4.6,10.2) -- (5.4,9.6) -- (5.2,8.4) -- (6.4,7.4) -- (7.4,6.6) -- (7.2,5.2) -- (6.6,3.8) -- (6.2,2.2) -- (5.0,0.6) -- (3.4,0.4) -- (2.2,1.4) -- (1.2,2.0) -- (0.8,3.2) -- (1.4,4.4) -- cycle;
% Zones culturales
\fill[popgold!55] (3.6,9.2) -- (4.6,10.2) -- (5.4,9.6) -- (5.2,8.4) -- (6.4,7.4) -- (6.2,6.6) -- (4.4,6.8) -- (3.1,7.6) -- cycle;
\fill[poporange!60] (1.0,2.6) -- (0.8,3.2) -- (1.4,4.4) -- (2.2,5.6) -- (2.8,4.6) -- (2.4,3.2) -- cycle;
\fill[popblue!45] (1.6,4.6) -- (2.2,5.6) -- (3.4,6.4) -- (3.1,7.6) -- (4.4,6.8) -- (4.0,5.4) -- (2.8,4.6) -- cycle;
\fill[poppurple!35] (4.4,6.8) -- (6.2,6.6) -- (7.4,6.6) -- (7.2,5.2) -- (6.6,3.8) -- (6.2,2.2) -- (5.0,0.6) -- (3.4,0.4) -- (2.6,1.6) -- (2.8,3.2) -- (4.0,5.4) -- cycle;
% Villes
\fill[popink] (2.9,3.6) circle (0.09) node[right=1pt,font=\small]{Yaoundé};
\fill[popink] (1.3,2.9) circle (0.09) node[left=1pt,font=\small]{Douala};
\fill[popink] (4.7,9.3) circle (0.09) node[right=1pt,font=\small]{Maroua};
\fill[popink] (2.6,4.9) circle (0.08) node[above left=0pt,font=\small]{Bafoussam};
% Nord
\draw[->,thick,popdark] (0.6,8.6) -- (0.6,9.6) node[above,font=\small]{N};
% Légende interne
\node[font=\small\bfseries,popdark] at (4.4,7.9) {Coton};
\node[font=\small\bfseries,popdark] at (1.7,3.9) {Banane};
\node[font=\small\bfseries,popdark] at (3.1,5.9) {Café};
\node[font=\small\bfseries,popdark] at (5.2,4.0) {Cacao, palme, hévéa};
\end{tikzpicture}
\end{center}
\legende{Carte schématique de la répartition des grandes cultures de rente au Cameroun. 1 : zone cotonnière (Grand Nord, savane sèche) ; 2 : bananier et palme (Littoral humide) ; 3 : café (Ouest et hauts plateaux) ; 4 : cacao, palmier à huile, hévéa (forêt du Centre-Sud-Est). Échelle indicative ; contour très simplifié.}
\end{popfigure}

\begin{methode}[Lire une carte de répartition agricole]
\begin{enumerate}
\item \textbf{Lire le titre} : que représente la carte (quelle culture, quelle échelle) ?
\item \textbf{Étudier la légende} : identifier chaque couleur ou symbole et ce qu'il signifie.
\item \textbf{Localiser} : situer les zones de culture par rapport aux régions, aux villes et aux grands repères (fleuves, relief, ports).
\item \textbf{Décrire} : dire où la culture est concentrée, où elle est absente.
\item \textbf{Expliquer par le milieu} : relier la répartition au climat (pluies, saisons), au sol, au relief, puis aux facteurs humains (routes, ports d'exportation, main-d'œuvre, histoire des plantations).
\end{enumerate}
\end{methode}

\begin{popfigure}
\begin{center}
\begin{tikzpicture}
\begin{axis}[
ybar, bar width=22pt,
width=0.85\linewidth, height=6.5cm,
symbolic x coords={2015,2018,2021,2023},
xtick=data,
ylabel={Production (milliers de tonnes)},
ymin=0, ymax=320,
legend style={at={(0.02,0.98)},anchor=north west,font=\small},
nodes near coords, every node near coord/.append style={font=\scriptsize},
axis lines=left, enlarge x limits=0.18]
\addplot[fill=popgold!80,draw=popdark] coordinates {(2015,232) (2018,250) (2021,290) (2023,295)};
\addplot[fill=popblue!70,draw=popdark] coordinates {(2015,32) (2018,28) (2021,25) (2023,24)};
\legend{Cacao, Café}
\end{axis}
\end{tikzpicture}
\end{center}
\legende{Production indicative de cacao et de café au Cameroun (en milliers de tonnes). Le cacao progresse et domine largement ; le café recule (vieillissement des vergers, concurrence). Données arrondies, à titre pédagogique.}
\end{popfigure}

\courssub{Les difficultés de l'agriculture camerounaise}

Malgré son importance, l'agriculture camerounaise reste fragile :

\begin{itemize}
\item \textbf{Faible mécanisation} : la houe domine encore ; les tracteurs sont rares, ce qui limite les surfaces cultivées et les rendements.
\item \textbf{Enclavement} : beaucoup de pistes rurales sont impraticables en saison des pluies ; les produits pourrissent faute de pouvoir rejoindre les villes et les ports.
\item \textbf{Fluctuation des prix mondiaux} : le prix du cacao ou du café est fixé sur les marchés internationaux ; quand il s'effondre, le revenu du planteur s'effondre aussi, sans qu'il puisse rien y faire.
\item \textbf{Vieillissement des planteurs} : les jeunes quittent les villages pour les villes (exode rural) ; l'âge moyen des exploitants augmente et les plantations ne sont pas renouvelées.
\item \textbf{Autres contraintes} : faible accès au crédit, parcelles morcelées, maladies des cultures (pourriture brune du cacao), conflits fonciers.
\end{itemize}

\courssub{Les efforts de modernisation}

Face à ces difficultés, l'État et ses partenaires mènent des actions :

\begin{itemize}
\item \textbf{L'encadrement des producteurs} : la \textbf{SODECOTON} fournit aux cotonculteurs semences, engrais et pesticides à crédit, encadre les paysans et achète la récolte à prix garanti ; l'\textbf{IRAD} (Institut de recherche agricole pour le développement) crée des variétés améliorées et résistantes aux maladies.
\item \textbf{La distribution d'intrants} : semences sélectionnées, engrais, produits phytosanitaires.
\item \textbf{Les agropoles} : pôles agro-industriels qui regroupent production, transformation et commercialisation pour créer de la valeur sur place.
\item \textbf{Le désenclavement} : réhabilitation des routes rurales et des voies ferrées.
\item \textbf{L'agriculture périurbaine} : production vivrière et maraîchère intensive autour des grandes villes.
\end{itemize}

\begin{savaistu}[L'agriculture périurbaine nourrit les villes]
Autour de Yaoundé et de Douala, des milliers de maraîchers cultivent tomates, carottes, choux, poireaux, piments et légumes feuilles sur les bas-fonds et les coteaux périphériques. Cette \cle{agriculture périurbaine} ravitaille directement les marchés urbains (marché Mokolo, marché Mboppi) en produits frais, souvent livrés le jour même. C'est l'illustration concrète des relations ville-campagne étudiées dans le dossier précédent : la ville offre un marché immense, la campagne proche fournit la nourriture.
\end{savaistu}

\begin{exoresolu}[Calculer une part et interpréter un graphique]
\textbf{Énoncé.} En 2021, le Cameroun a produit environ \num{290000} tonnes de cacao. La production mondiale la même année était d'environ \num{5000000} tonnes. (1) Calcule la part du Cameroun dans la production mondiale. (2) À l'aide du diagramme en barres, compare l'évolution du cacao et du café entre 2015 et 2023.
\tcblower
\textbf{Solution.} (1) Part du Cameroun :
\[
\frac{290000}{5000000}\times 100 = 5,8\ \%.
\]
Le Cameroun produit donc environ 5,8 \% du cacao mondial, ce qui confirme son rang de 5\textsuperscript{e} producteur mondial.
(2) Entre 2015 et 2023, la production de cacao est passée d'environ \num{232000} à \num{295000} tonnes : elle est en nette progression (environ $+27\%$). À l'inverse, le café est passé d'environ \num{32000} à \num{24000} tonnes : il recule (environ $-25\%$), à cause du vieillissement des caféiers et de la baisse des prix. Le cacao confirme sa place de première culture de rente du pays.
\end{exoresolu}

\begin{exercice}[Pour s'entraîner]
(1) Définis agriculture vivrière et agriculture de rente en une phrase chacune. (2) Recopie et complète : le coton se cultive surtout dans \dots ; la banane d'exportation dans \dots ; le cacao dans \dots. (3) Cite trois difficultés de l'agriculture camerounaise et, pour chacune, une solution en cours. (4) Explique en quelques lignes pourquoi l'agriculture périurbaine illustre les relations ville-campagne.
\end{exercice}

\begin{corrige}[Exercice]
(1) L'agriculture vivrière produit d'abord pour la consommation de la famille ; l'agriculture de rente produit pour la vente et l'exportation. (2) Le Grand Nord (Nord, Adamaoua, Extrême-Nord) ; le Littoral (plaine de Penja, Tiko) ; le Centre-Sud-Ouest. (3) Exemples : faible mécanisation $\rightarrow$ agropoles et équipement ; enclavement $\rightarrow$ réhabilitation des pistes rurales ; fluctuation des prix $\rightarrow$ prix garantis (SODECOTON), transformation locale ; vieillissement des planteurs $\rightarrow$ installation des jeunes, accès au crédit. (4) Les maraîchers installés autour des villes utilisent le marché urbain comme débouché direct : la ville consomme, la campagne proche produit et livre rapidement des denrées fraîches ; c'est un échange quotidien de biens et de revenus entre les deux milieux.
\end{corrige}

\begin{aretenir}[L'essentiel]
L'agriculture camerounaise emploie environ 40 \% de la population active. Elle combine une agriculture vivrière traditionnelle (manioc, plantain, igname, maïs, arachide ; brûlis, houe, jachère) et une agriculture de rente (cacao, café, banane, palme, hévéa, coton) organisée autour de grandes sociétés (CDC, PHP, SOSUCAM, SODECOTON). Sa répartition suit le milieu : cacao-café au Centre-Sud-Ouest, banane au Littoral, coton au Nord. Freinée par la faible mécanisation, l'enclavement, les prix mondiaux et le vieillissement des planteurs, elle se modernise grâce à l'encadrement (SODECOTON, IRAD), aux intrants, aux agropoles et à l'agriculture périurbaine.
\end{aretenir}

\coursec{L'élevage au Cameroun}

\courssub{Transition : de l'agriculture à l'élevage}

Après avoir analysé les systèmes agraires et les cultures qui structurent l'espace rural camerounais, il est naturel de nous tourner vers la seconde composante majeure du secteur primaire : \textbf{l'élevage}. Au Cameroun comme dans la plupart des pays du Sud, agriculture et élevage sont intimement liés : les résidus de cultures nourrissent les animaux, la fumure animale entretient la fertilité des sols, et la traction animale (bœufs de labour) équipe encore de nombreuses exploitations familiales, notamment dans les zones cotonnières du Nord. Cette complémentarité explique pourquoi les deux activités sont souvent menées par les mêmes ménages, selon des modalités très variables selon les régions. L'élevage n'est pas une simple annexe de l'agriculture : il représente une source essentielle de protéines animales, de revenus monétaires et de sécurité alimentaire pour des millions de Camerounais.

\begin{definition}[L'élevage]
L'élevage est l'ensemble des techniques d'élevage, de reproduction et d'exploitation des animaux domestiques (bovins, ovins, caprins, porcins, volailles, etc.) en vue d'obtenir des produits (viande, lait, œufs, cuir, force de travail) et des services (fumure, traction, épargne sur pied). Il se distingue par le mode de conduite du troupeau (extensif ou intensif) et par son insertion dans l'espace rural.
\end{definition}

\courssub{Importance économique, sociale et nutritionnelle de l'élevage}

L'élevage occupe une place stratégique dans l'économie camerounaise. Il contribue à hauteur d'environ \textbf{12 à 15 \% du PIB agricole} et emploie directement ou indirectement plusieurs millions de personnes (éleveurs, bouchers, transporteurs, vétérinaires, transformateurs). Sur le plan nutritionnel, il fournit la quasi-totalité de la viande rouge consommée dans le pays (bœuf, mouton, chèvre) et une part croissante de la viande blanche (volaille) et des œufs. Les produits laitiers, bien que encore marginaux dans la consommation moyenne, se développent autour des bassins de production.

\begin{exemplebox}[Le cheptel bovin camerounais : un chiffre clé]
Selon les données du MINEPIA (Ministère de l'Élevage, des Pêches et des Industries Animales) et de la FAO, le cheptel bovin national est estimé à \textbf{plus de 6 millions de têtes} (environ 6,2 millions en 2022). Il est constitué à plus de 90 \% de races zébus (\emph{Bos indicus}) : Goudali, Bokolodji, White Fulani, Shuwa Arab. Ce cheptel est \textbf{concentré à plus de 70 \% dans la région de l'Adamaoua}, qui constitue le véritable « poumon bovin » du pays. Les régions du Nord et de l'Extrême-Nord en détiennent ensemble environ 20 \%, le reste étant dispersé dans les zones de pâturage de l'Ouest, du Nord-Ouest et du Centre.
\end{exemplebox}

Au-delà de la production marchande, l'élevage remplit des fonctions sociales vitales : \textbf{épargne sur pied} (capital mobile, assurant contre les aléas climatiques), \textbf{prestige et statut social} (taille du troupeau), \textbf{dotations matrimoniales}, \textbf{force de travail} (bœufs de trait dans les zones cotonnières du Nord) et \textbf{fumure organique} (indispensable au maintien de la fertilité dans les systèmes de culture sur brûlis ou en jachère courte).

\courssub{Les deux grands systèmes d'élevage : extensif (traditionnel) et intensif (moderne)}

La dualité structurelle de l'élevage camerounais repose sur l'opposition entre un \textbf{élevage extensif pastoral}, hérité des traditions peules et arabes, et un \textbf{élevage intensif periurbain}, né de la demande urbaine croissante en protéines animales. Cette dualité se lit dans l'espace, dans les techniques, dans les acteurs et dans les performances.

\begin{definition}[Élevage extensif (pastoral traditionnel)]
Système d'élevage caractérisé par une \textbf{faible densité animale par hectare}, l'utilisation de \textbf{parcours naturels non aménagés}, une \textbf{main-d'œuvre familiale}, des \textbf{intrants externes quasi nuls} (aliments achetés, soins vétérinaires curatifs seulement), et une \textbf{mobilité spatiale} (transhumance) pour suivre la ressource fourragère et l'eau. Les races sont locales (zébus, petits ruminants rustiques). La productivité par animal est faible (âge à l'abattage tardif, taux de réforme élevé, production laitière faible), mais la productivité par unité de travail ou par franc investi peut être compétitive.
\end{definition}

\begin{definition}[Élevage intensif (moderne, hors-sol ou semi-hors-sol)]
Système d'élevage caractérisé par une \textbf{forte densité animale}, la \textbf{stabulisation permanente} (bâtiments), l'apport d'\textbf{aliments concentrés achetés} (maïs, soja, tourteaux, prémix), une \textbf{prophylaxie rigoureuse} (vaccinations, traitements préventifs), l'utilisation de \textbf{races améliorées ou hybrides} (poulets de chair Cobb/Hubbard, porcs Large White/Landrace, vaches laitières Holstein/Gir), et une \textbf{recherche de productivité maximale par animal et par m²}. Il est localisé à proximité des centres de consommation (villes) et d'approvisionnement en intrants.
\end{definition}

\begin{definition}[Transhumance]
Déplacement \textbf{saisonnier, cyclique et organisé} d'une partie ou de la totalité d'un troupeau (souvent bovin) entre deux aires géographiques complémentaires : une \textbf{zone de saison sèche} (bas-fonds, berges de fleuves, zones de pâturages secs, points d'eau permanents) et une \textbf{zone de saison des pluies} (plateaux, hautes altitudes, parcours herbagers riches mais sans eau permanente en saison sèche). La transhumance suit des \textbf{couloirs de passage traditionnels} (couloirs de transhumance) et des \textbf{dates codifiées} par la coutume et, de plus en plus, par la réglementation administrative. Elle n'est \textbf{pas un nomadisme erratique} : l'éleveur a un terroir d'attache, des droits d'usage coutumiers, et le parcours est planifié.
\end{definition}

\begin{attention}[Transhumance $\neq$ Nomadisme]
Ne confondez jamais \textbf{transhumance} (migration saisonnière aller-retour entre deux points fixes, troupeau conduit par des bergers, famille sédentarisée au camp de base) et \textbf{nomadisme} (déplacement permanent de l'ensemble de la cellule familiale et du troupeau sans base fixe). Au Cameroun, les éleveurs peuls sont très majoritairement \textbf{transhumants semi-sédentarisés} (camp de base en saison des pluies, camp de saison sèche en bas-fonds). Le nomadisme pur a quasiment disparu.
\end{attention}

\courssub{L'élevage extensif : le domaine du zébu et de la transhumance}

L'élevage extensif concerne essentiellement les \textbf{bovins} (zébus) et les \textbf{petits ruminants} (moutons Djallonké, chèvres du Sahel). Il s'étend sur la \textbf{zone soudano-sahélienne} (Adamaoua, Nord, Extrême-Nord) et les \textbf{plateaux de l'Ouest et du Nord-Ouest} (zones de pâturages d'altitude).

\begin{itemize}
    \item \textbf{L'Adamaoua, cœur du système extensif :} Cette région de plateau (1000--1300 m d'altitude) bénéficie d'une pluviométrie de 1500--1600 mm/an concentrée sur 6--7 mois, d'une herbe abondante en saison des pluies (Andropogon, Hyparrhenia, Pennisetum) et de nombreux points d'eau permanents (lacs de cratère, rivières). C'est la zone de \textbf{saison des pluies} par excellence. Les troupeaux y séjournent de mai à novembre.
    \item \textbf{La transhumance vers le Nord et l'Extrême-Nord :} En saison sèche (novembre--avril), l'herbe sèche et l'eau se raréfie sur le plateau. Les troupeaux descendent vers les \textbf{bas-fonds de la Bénoué, du Logone, du Chari, du Mayo Rey}, où persistent des pâturages verts (végétation hygrophile) et de l'eau. C'est la \textbf{transhumance descendante} (ou « transhumance sèche »). Le retour (transhumance montante) s'opère avec les premières pluies.
    \item \textbf{Les petits ruminants :} Moutons et chèvres accompagnent souvent les troupeaux bovins ou sont conduits par les femmes et les enfants autour des campements. Ils sont plus rustiques, se reproduisent plus vite (intervalle de mise bas 6--8 mois) et constituent l'épargne de premier niveau.
    \item \textbf{Productivité :} Âge à l'abattage 4--5 ans (vs 18--24 mois en intensif), taux de réforme 8--10 \%, production laitière 1--2 L/jour/vache (vs 15--25 L en intensif amélioré). Mais coût de production très faible (quasi nul en intrants achetés).
\end{itemize}

\begin{savaistu}[Les conflits agriculteurs-éleveurs du Nord-Ouest : une concurrence pour l'espace]
Dans la région du Nord-Ouest (plateaux de Bamenda, Wum, Nkambe), l'expansion des cultures vivrières (maïs, haricot, pomme de terre) et de rente (café, palmier à huile) a grignoté les parcours traditionnels et les couloirs de transhumance. Les éleveurs mbororo (peuls du Nord-Ouest) voient leurs itinéraires bloqués par les champs, tandis que les agriculteurs subissent la divagation des troupeaux sur leurs récoltes. Ces tensions, exacerbées par la crise anglophone depuis 2016, ont causé des centaines de morts, des destructions de villages et un déplacement massif de populations. Elles illustrent la \textbf{concurrence pour l'espace rural} entre deux modes d'appropriation du sol : l'appropriation par la mise en culture (droit coutumier « mise en valeur = propriété ») et l'appropriation par le parcours (droit d'usage coutumier « premier occupant = droit de pâture »). La loi n° 74-1 du 6 juillet 1974 sur le régime foncier et la loi n° 98-005 du 14 avril 1998 sur l'élevage tentent de réguler ces conflits (création de couloirs de transhumance, de ranchs, commissions de conciliation), mais l'application reste faible.
\end{savaistu}

\courssub{L'élevage intensif : la réponse periurbaine à la demande urbaine}

Face à l'explosion de la population urbaine (Yaoundé \textasciitilde 4,5 millions, Douala \textasciitilde 4 millions, Bafoussam \textasciitilde 500 000 hab. en 2024) et à l'insuffisance de l'offre extensive (viande dure, saisonnalité, circuits longs), s'est développé depuis les années 1980 un \textbf{élevage intensif periurbain}, porté par des entrepreneurs urbains (fonctionnaires, commerçants, retraités) et des jeunes formés. Ce système illustre parfaitement les \textbf{relations ville-campagne} : la ville (demande, capital, main-d'œuvre qualifiée) commande l'organisation de la campagne periurbaine (production, foncier, déchets urbains valorisés).

\begin{itemize}
    \item \textbf{Aviculture industrielle :} Élevage de \textbf{poulets de chair} (cycles de 35--42 jours, 2,2--2,5 kg) et de \textbf{poules pondeuses} (cycles de 52--72 semaines, 300--330 œufs/poule). Localisé dans les ceintures de 20--50 km autour des grandes villes : \textbf{Mbankomo, Nkolbisson, Obala, Mfou} (autour de Yaoundé) ; \textbf{Manoka, Bonabéri, Nkongsamba, Penja} (autour de Douala) ; \textbf{Bamougoum, Bayangam, Bandjoun} (autour de Bafoussam). Les élevages comptent de 500 à 50 000 sujets. Alimentation : maïs (60--65 \%), tourteau de soja (25--30 \%), prémix vitaminiques. Biosécurité : clôture, sas de désinfection, vaccination (Newcastle, Gumboro, Bronchite infectieuse).
    \item \textbf{Porcherie intensive :} Élevage de \textbf{porcs hybrides} (Large White $\times$ Landrace $\times$ Piétrain) en maternité, post-sevrage, engraissement. Cycles de 5--6 mois pour 100--110 kg. Même localisation periurbaine. Alimentation : son de blé, maïs, tourteau, déchets de brasseries (drêches), déchets alimentaires urbains (récupération). Gestion des effluents : lagunage, biométhanisation naissante.
    \item \textbf{Élevage laitier intensif naissant :} Fermes laitières (vaches Holstein, Gir, croisés Girolando) à \textbf{Nkolbisson, Obala, Dschang, Maroua}. Production de 15--25 L/jour/vache. Transformation sur place (yaourts, fromages frais, lait pasteurisé) pour circuits courts urbains.
\end{itemize}

\begin{exemplebox}[L'aviculture periurbaine de Yaoundé : un bassin de 15 millions de poulets/an]
La ceinture de Yaoundé concentre environ \textbf{60 \% de la production nationale de poulets de chair du secteur moderne} (estimée à 20--25 millions de têtes/an). Les fermes de \textbf{Mbankomo, Nkolbisson, Obala} approvisionnent quotidiennement les marchés de Mfoundi, Mokolo, Nkolndongo. Un cycle de 42 jours produit un poulet de 2,3 kg avec un indice de consommation (IC) de 1,7--1,8 kg d'aliment/kg de poids vif. Le coût de production (aliment + poussins + charges) tourne autour de \textbf{1 100--1 200 FCFA/kg vif}, pour un prix de vente ferme de 1 400--1 500 FCFA/kg, laissant une marge nette sensible mais vulnérable à la volatilité du prix du maïs et du soja (importé à 40--50 \%).
\end{exemplebox}

\begin{popfigure}
\begin{tikzpicture}[scale=0.85]
% Contour très simplifié du Cameroun
\draw[popdark, thick] (0,0) -- (1,0.2) -- (2,0.5) -- (3,0.3) -- (4,0) -- (4.5,-0.5) -- (4.2,-1.5) -- (3.5,-2.2) -- (2.5,-2.5) -- (1.5,-2.2) -- (0.5,-1.5) -- (0,-0.5) -- cycle;
% Régions clés
\fill[popgreen!30] (1.5,0.5) ellipse (0.8 and 0.6); % Adamaoua
\node[popdark, font=\small\bfseries] at (1.5,0.5) {Adamaoua};
\fill[popblue!20] (0.8,-1.8) ellipse (0.7 and 0.5); % Nord/Extrême-Nord
\node[popdark, font=\small] at (0.8,-1.8) {Nord / Extrême-Nord};
\fill[poporange!30] (3.2,0.2) ellipse (0.6 and 0.4); % Ouest/Nord-Ouest
\node[popdark, font=\small] at (3.2,0.2) {Ouest / NO};
% Villes principales
\node[circle, fill=popink, inner sep=1.5pt, label=above:\tiny Yaoundé] (Y) at (2.2,-0.8) {};
\node[circle, fill=popink, inner sep=1.5pt, label=right:\tiny Douala] (D) at (3.8,-0.5) {};
\node[circle, fill=popink, inner sep=1.5pt, label=above:\tiny Bafoussam] (B) at (3.0,0.5) {};
\node[circle, fill=popink, inner sep=1.5pt, label=left:\tiny Ngaoundéré] (N) at (1.2,0.8) {};
\node[circle, fill=popink, inner sep=1.5pt, label=below:\tiny Maroua] (M) at (0.5,-2.0) {};
% Flèches transhumance (Adamaoua -> Bas-fonds)
\draw[popgreen, very thick, ->, >=Stealth] (1.5,0.0) -- (1.0,-1.2) node[midway, left, font=\tiny, popdark] {Transhumance sèche};
\draw[popgreen, very thick, ->, >=Stealth] (1.5,0.0) -- (0.7,-1.8) node[midway, left, font=\tiny, popdark] {Transhumance sèche};
\draw[popgreen, dashed, ->, >=Stealth] (0.8,-1.8) -- (1.5,0.5) node[midway, right, font=\tiny, popdark] {Retour pluies};
% Ceintures periurbaines (cercles pointillés)
\draw[poporange, dashed, thick] (Y) circle (0.6);
\draw[poporange, dashed, thick] (D) circle (0.6);
\draw[poporange, dashed, thick] (B) circle (0.45);
\node[poporange, font=\tiny] at (2.2,-1.5) {Bassin avicole/porcin};
\node[poporange, font=\tiny] at (4.2,-1.2) {Bassin avicole/porcin};
\node[poporange, font=\tiny] at (3.0,1.1) {Bassin laitier/avicole};
% Nord
\draw[->, >=Stealth] (0, -2.8) -- (0, -2.5) node[below, font=\small] {N};
\end{tikzpicture}
\legende{Croquis schématique : organisation de l'espace par l'élevage au Cameroun. En vert : axes de transhumance (zébu extensif) entre l'Adamaoua (saison des pluies) et les bas-fonds du Nord/Extrême-Nord (saison sèche). En orange : ceintures periurbaines d'élevage intensif (aviculture, porcherie, lait) autour des métropoles (Yaoundé, Douala, Bafoussam).}
\end{popfigure}

\begin{methode}[Comparer les deux systèmes d'élevage (type Bac/Probatoire)]
\textbf{Consigne :} « À l'aide d'un tableau, mettez en évidence les différences structurelles entre l'élevage extensif pastoral et l'élevage intensif periurbain au Cameroun. »
\begin{enumerate}
    \item \textbf{Identifiez les critères de comparaison :} Localisation / Milieu / Acteurs / Races / Alimentation / Santé / Productivité / Commercialisation / Impact environnemental.
    \item \textbf{Remplissez le tableau en mobilisant les données du cours :} Chiffres clés (cheptel, productivité), noms de races, noms de localités.
    \item \textbf{Concluez sur la complémentarité :} L'extensif fournit le cheptel de base (reproduction, rusticité) et la viande rouge ; l'intensif répond à la demande urbaine rapide (volaille, porc, lait) et valorise les sous-produits urbains.
\end{enumerate}
\begin{center}
\begin{tabularx}{\linewidth}{|l|X|X|}\hline
\textbf{Critère} & \textbf{Élevage extensif (pastoral)} & \textbf{Élevage intensif (periurbain)} \\ \hline
\textbf{Localisation} & Adamaoua, Nord, Extrême-Nord, Plateaux Ouest/NO & Ceintures 20-50 km autour de Yaoundé, Douala, Bafoussam \\ \hline
\textbf{Acteurs} & Éleveurs peuls/mbororo (familiaux, traditionnels) & Entrepreneurs urbains, jeunes formés, sociétés commerciales \\ \hline
\textbf{Races} & Zébus (Goudali, Bokolodji), Djallonké, Chèvre du Sahel & Hybrides : Cobb/Hubbard (poulet), Large White (porc), Holstein/Gir (lait) \\ \hline
\textbf{Alimentation} & Parcours naturels, résidus de culture, browse & Concentrés achetés (maïs, soja, prémix), drêches, déchets urbains \\ \hline
\textbf{Sanitaire} & Curatif, médecine traditionnelle, vaccination campagnes & Prophylaxie rigoureuse, vétérinaire privé, biosécurité \\ \hline
\textbf{Productivité} & Faible/an (1-2 L lait, abattage 4-5 ans) mais faible coût & Élevée/an (15-25 L lait, abattage 42 j poulet) mais coût intrants fort \\ \hline
\textbf{Circuit} & Long (pistes, marchés ruraux, abattoirs urbains) & Court (ferme $\to$ grossiste/transformateur $\to$ marché urbain) \\ \hline
\textbf{Enjeu env.} & Dégradation parcours, feux de brousse, conflits fonciers & Pollution effluents (azote, phosphore), pression foncière periurbaine \\ \hline
\end{tabularx}
\end{center}
\end{methode}

\begin{exoresolu}[Sujet type Probatoire : L'élevage face à l'urbanisation]
\textbf{Énoncé :} « L'urbanisation rapide du Cameroun transforme les systèmes d'élevage. Montrez comment l'élevage intensif periurbain s'adapte à la demande urbaine et quels en sont les limites. » \tcblower
\textbf{Solution détaillée :}
\begin{enumerate}
    \item \textbf{Adaptation à la demande (Offre) :}
    \begin{itemize}
        \item \textbf{Proximité des marchés :} Localisation stratégique (Mbankomo, Manoka, Bamougoum) pour réduire coûts de transport et pertes (chaîne du froid faible).
        \item \textbf{Spécialisation produit :} Volaille (cycle court 42 j, protéines abordables), Porc (valorisation déchets urbains/brasseries), Lait (produits frais, yaourts, fromages pour classe moyenne urbaine).
        \item \textbf{Intensification technique :} Races à croissance rapide, alimentation calibrée, biosécurité $\to$ standardisation du produit (poids, qualité) exigée par les consommateurs urbains et la restauration rapide.
    \end{itemize}
    \item \textbf{Valorisation des relations ville-campagne :}
    \begin{itemize}
        \item \textbf{Intrants :} Maïs/soja du Nord/Centre $\to$ usines d'aliments (SMA, PROVIMI) $\to$ fermes periurbaines.
        \item \textbf{Déchets urbains :} Drêches (brasseries), son (minoteries), déchets alimentaires (marchés) $\to$ aliment porcs/poules (économie circulaire).
        \item \textbf{Main-d'œuvre/Capital :} Fonctionnaires/Commerçants urbains investissent (épargne salariale) ; main-d'œuvre jeune locale ou migrante.
    \end{itemize}
    \item \textbf{Limites et vulnérabilités :}
    \begin{itemize}
        \item \textbf{Dépendance aux importations :} 40-50 \% tourteau de soja, prémix, matériel génétique (poussins d'un jour, semences porcines), équipement $\to$ sensibilité au taux de change et cours mondiaux.
        \item \textbf{Coût de l'aliment :} 65-70 \% charges de production. Flambée prix maïs/soja $\to$ compression marges, arrêt fermes (crise 2021-2023).
        \item \textbf{Pression foncière/environnementale :} Urbanisation galopante grignote terres agricoles periurbaines ; gestion effluents (fientes, lisiers) souvent défaillante $\to$ pollution nappes/eaux de surface (ex: Mfoundi, Wouri).
        \item \textbf{Sanitaire :} Concentration animale = risque épizootique majeur (Influenza aviaire, Peste porcine africaine) $\to$ biosécurité coûteuse.
    \end{itemize}
    \item \textbf{Conclusion :} L'élevage intensif periurbain est la réponse \textbf{de marché} à l'urbanisation (réactivité, standardisation), mais sa viabilité dépend de la \textbf{souveraineté alimentaire en intrants} (maïs/soja locaux) et de la \textbf{régulation foncière/environnementale} periurbaine.
\end{enumerate}
\end{exoresolu}

\begin{exercice}[Entraînement : Les relations ville-campagne par l'élevage]
\begin{enumerate}
    \item Sur une carte muette du Cameroun, localisez et nommez : (a) La région concentrant 70 \% du cheptel bovin. (b) Deux bassins d'aviculture periurbaine (autour de Yaoundé et Douala). (c) Un couloir de transhumance majeur (flèche orientée).
    \item Calculez le taux de réforme annuel moyen d'un troupeau extensif de 100 têtes si 9 bêtes sont vendues/abattues dans l'année. Interprétez ce chiffre face à un troupeau intensif (taux de réforme 30-40 \%).
    \item Expliquez pourquoi le prix du poulet de chair periurbain reste sensible au cours mondial du soja, alors que le bœuf « sur pied » de l'Adamaoua ne l'est pas directement.
\end{enumerate}
\end{exercice}

\begin{corrige}[Exercice n° Entraînement]
\begin{enumerate}
    \item (a) \textbf{Région de l'Adamaoua} (chef-lieu Ngaoundéré). (b) \textbf{Bassin de Yaoundé :} Mbankomo, Nkolbisson, Obala, Mfou. \textbf{Bassin de Douala :} Manoka, Bonabéri, Nkongsamba, Penja. (c) Flèche \textbf{Adamaoua $\to$ Bas-fonds Bénoué/Logone} (saison sèche, nov.-avr.) et flèche retour (saison des pluies, mai-juin).
    \item Taux de réforme = (Nombre de sorties / Effectif moyen) $\times$ 100 = (9 / 100) $\times$ 100 = \textbf{9 \%}. Interprétation : Taux faible caractéristique de l'extensif (renouvellement lent, vaches réformées vieilles, faible productivité numérique). En intensif (volaille/porc), taux 30-40 \% (renouvellement rapide des lots, rotation du capital rapide).
    \item Le poulet de chair consomme des aliments concentrés intégrant du \textbf{tourteau de soja importé à 40-50 \%} (Argentine, Brésil, USA). Le prix de l'aliment (65-70 \% du coût) suit le cours mondial (CBOT Chicago) et le taux de change FCFA/USD. Le bœuf extensif se nourrit \textbf{quasi exclusivement de parcours naturels gratuits} (herbe, browse) et de résidus de culture locaux ; son coût de production est découplé du marché mondial des grains, seul le prix de vente final suit l'inflation urbaine.
\end{enumerate}
\end{corrige}

\courssub{Vers la synthèse : Villes, campagnes et développement}

L'analyse de l'élevage révèle la \textbf{profonde interdépendance ville-campagne} au Cameroun. La ville (demande de protéines, capital, déchets, main-d'œuvre) structure l'espace rural : elle aspire les flux de bétail sur pied depuis l'Adamaoua (circuits longs, informels) et fait naître des ceintures d'élevage intensif high-tech dans sa périphérie immédiate. Inversement, la campagne nourrit la ville mais subit la pression foncière (periurbain) et la dégradation des ressources (parcours, eau). La \textbf{sécurité alimentaire} du pays repose sur la capacité à \textbf{articuler ces deux systèmes} : moderniser l'extensif (aménagement parcours, santé animale, sédentarisation partielle, croisement) pour sécuriser l'approvisionnement en viande rouge, et sécuriser l'intensif (filières maïs/soja locales, gestion effluents, foncier agricole protégé) pour maîtriser le coût des protéines blanches. C'est tout l'enjeu du \textbf{Document de Stratégie de Développement du Secteur de l'Élevage (DSDSE)} et de la \textbf{Politique Nationale de Développement de l'Élevage (PNDE)}.

\coursec{Synthèse : villes, campagnes et développement}

\courssub{Transition : de l'élevage aux dynamiques territoriales}

Après avoir analysé les systèmes d'élevage au Cameroun — de l'élevage extensif sahélien aux unités intensives périurbaines — nous comprenons que la production animale ne s'arrête pas à la ferme. Elle s'insère dans des circuits commerciaux qui relient les zones rurales productrices aux centres urbains consommateurs. Cette section finale replace l'agriculture et l'élevage dans leur rôle structurant : nourrir les villes, animer les campagnes, et, par leur interdépendance, construire le développement national. Nous mobiliserons les notions de \cle{sécurité alimentaire}, d'\cle{aménagement du territoire} et de \cle{développement durable} pour lire les territoires camerounais comme un système cohérent.

\courssub{Le rôle de l'agriculture et de l'élevage dans l'approvisionnement des villes}

\begin{definition}[Sécurité alimentaire]
La \cle{sécurité alimentaire} existe lorsque « tous les êtres humains, à tout moment, ont un accès physique, social et économique à une alimentation suffisante, saine et nutritive, qui répond à leurs besoins énergétiques et à leurs préférences alimentaires pour une vie active et saine » (FAO, 1996). Elle repose sur quatre piliers : \textbf{disponibilité} (production + importations + stocks), \textbf{accessibilité} (pouvoir d'achat, infrastructures, marchés), \textbf{utilisation} (qualité nutritionnelle, eau potable, santé) et \textbf{stabilité} (résilience aux chocs climatiques, économiques, sécuritaires).
\end{definition}

Au Cameroun, les villes concentrent plus de 58 \% de la population (environ 16 millions d'habitants en 2024) mais ne produisent qu'une fraction de leur alimentation. L'approvisionnement repose sur trois piliers complémentaires :

\begin{center}
\begin{tabularx}{\linewidth}{|l|X|X|}\hline
\textbf{Filière} & \textbf{Origine principale} & \textbf{Flux vers les villes} \\ \hline
Céréales (maïs, mil, sorgho, riz) & Nord (Adamaoua, Nord, Extrême-Nord), Ouest & Camions bâchés, trains (Ngaoundéré–Douala), pirogues (Logone) \\ \hline
Tubercules (manioc, igname, macabo, patate) & Centre, Sud, Est, Ouest, Sud-Ouest & Véhicules légers, motos, marchés de gros (Mfoundi, Nkolfoulou) \\ \hline
Produits animaux (viande, lait, œufs, poisson) & Nord (bovins), Ouest/Adamaoua (lait), zones côtières et lacs (poisson) & Pied (transhumance), camions réfrigérés, filières courtes urbaines \\ \hline
\end{tabularx}
\end{center}

\begin{exemplebox}[Le corridor Nord–Douala : l'exemple du maïs]
Chaque année, plus de 300 000 tonnes de maïs quittent les plaines de la Bénoué et du Diamaré vers Douala et Yaoundé. Le trajet (\SIrange{800}{1000}{km}) dure 3 à 5 jours par camion de \SI{30}{t}. Le coût de transport représente 25 à 35 \% du prix final au détail. Les pertes post-récolte (manutention, humidité, parasites) atteignent 15–20 \%. Ce corridor illustre la \cle{vulnérabilité logistique} des villes : une rupture (crise carburant, insécurité routière, inondation du pont sur la Bénoué) fait flamber les prix en 48 h.
\end{exemplebox}

\begin{popfigure}
\begin{tikzpicture}[scale=0.9, every node/.style={font=\small}]
% Campagne (gauche)
\begin{scope}[xshift=0cm]
\draw[popgreen!80!popdark, thick, fill=popgreen!10, rounded corners=8pt] (0,0) rectangle (5,6);
\node[popgreen!80!popdark, font=\bfseries] at (2.5,5.5) {CAMPAGNE};
\node[align=center] at (2.5,4.8) {Production};
\node[align=center] at (2.5,4.1) {• Agriculture : céréales, tubercules,\\  légumes, fruits, café, cacao};
\node[align=center] at (2.5,3.1) {• Élevage : bovins, ovins, caprins,\\  porcs, volailles, poisson};
\node[align=center] at (2.5,2.1) {• Main-d'œuvre rurale,\\  savoirs locaux, biodiversité};
\node[align=center] at (2.5,1.1) {• Besoins : intrants, équipements,\\  services, formation, santé};
\end{scope}
% Flèches centrale
\draw[poporange, very thick, -{Stealth[length=4mm]}] (5,4.5) -- (7,4.5) node[midway, above, align=center] {Produits agricoles,\\denrées, matières premières};
\draw[popblue, very thick, -{Stealth[length=4mm]}] (7,2.5) -- (5,2.5) node[midway, below, align=center] {Capitaux, intrants,\\services, revenus, innovations};
% Ville (droite)
\begin{scope}[xshift=9cm]
\draw[poppurple!80!popdark, thick, fill=poppurple!10, rounded corners=8pt] (0,0) rectangle (5,6);
\node[poppurple!80!popdark, font=\bfseries] at (2.5,5.5) {VILLE};
\node[align=center] at (2.5,4.8) {Consommation \& Transformation};
\node[align=center] at (2.5,4.1) {• Marchés de gros et de détail};
\node[align=center] at (2.5,3.4) {• Industries agroalimentaires\\  (minoteries, huileries, abattoirs,\\  laiteries, brasseries)};
\node[align=center] at (2.5,2.4) {• Population urbaine,\\  demande diversifiée};
\node[align=center] at (2.5,1.4) {• Capitaux financiers,\\  banques, assurances};
\node[align=center] at (2.5,0.7) {• Services : santé, éducation,\\  recherche, vulgarisation};
\end{scope}
% Boucle de rétroaction (bas)
\draw[popgreen!60!popdark, thick, -{Stealth[length=3mm]}] (2.5,-0.5) -- (2.5,-1.2) -| (11.5,-1.2) -- (11.5,0);
\node[popgreen!60!popdark, align=center, font=\bfseries\small] at (7,-1.5) {CIRCUIT VERTUEUX :\\revenus ruraux $\rightarrow$ demande urbaine $\rightarrow$ investissements ruraux};
% Légende
\node[align=left, font=\tiny] at (7,-2.3) {
\textbullet\ Flèche orange = flux physiques (biomasse, denrées)\\
\textbullet\ Flèche bleue = flux financiers, techniques, humains\\
\textbullet\ Boucle verte = réinvestissement et développement partagé
};
\end{tikzpicture}
\legende{Schéma de synthèse : interdépendance ville–campagne. La campagne produit denrées et matières premières ; la ville consomme, transforme, et renvoie capitaux, services et innovations. Un développement durable suppose que la boucle de rétroaction (flèche verte) soit effective et équitable.}
\end{popfigure}

\courssub{L'interdépendance ville–campagne comme moteur du développement national}

L'interdépendance ne se limite pas aux flux physiques. Elle structure l'économie nationale :

\begin{itemize}
\item \textbf{Demande urbaine $\rightarrow$ spécialisation rurale :} La croissance des villes (Yaoundé +4,2 \%/an, Douala +3,8 \%/an) crée des débouchés stables pour le riz de la plaine du Logone, le maïs de l'Adamaoua, le lait de l'Ouest, le poisson du Lac Tchad et de la côte. Les agriculteurs adaptent leurs assolements (ex. : maïs $\rightarrow$ maïs + soja pour l'alimentation animale).
\item \textbf{Revenus ruraux $\rightarrow$ demande de biens urbains :} La vente des productions agricoles finance l'achat de motos, téléphones, matériaux de construction, carburant, services de santé — produits et services urbains ou importés.
\item \textbf{Main-d'œuvre circulante :} La migration saisonnière (jeunes ruraux vers chantiers urbains en saison sèche, retour aux champs en saison des pluies) diffuse compétences et revenus. Les transferts de fonds (remittances) des migrants urbains vers les villages financent l'éducation, la santé, l'équipement agricole.
\item \textbf{Innovation descendante :} Variétés améliorées (IRAD), intrants (engrais, aliments composés), conseil agricole (MINADER, ONG), mécanisation (location de tracteurs, décortiqueuses) partent des villes vers les campagnes.
\end{itemize}

\begin{aretenir}[Interdépendance ville–campagne au Cameroun]
\begin{itemize}
\item Les villes ne peuvent exister sans l'approvisionnement alimentaire rural (disponibilité).
\item Les campagnes ne peuvent se moderniser sans les marchés, capitaux et services urbains (accessibilité, innovation).
\item Le développement national repose sur la \textbf{fluidité des flux} (routes, rail, énergie, télécoms) et l'\textbf{équité de la répartition de la valeur ajoutée} (prix à la production vs prix à la consommation).
\item \textbf{Indicateur clé :} le ratio prix producteur / prix consommateur. Pour le maïs, il est de 0,45–0,55 ; pour le bovin sur pied, 0,60–0,70. L'amélioration de ce ratio (stockage, transformation locale, circuits courts) est un levier majeur de réduction de la pauvreté rurale.
\end{itemize}
\end{aretenir}

\courssub{Les défis communs : sécurité alimentaire, emploi des jeunes, aménagement du territoire}

\begin{definition}[Aménagement du territoire]
L'\cle{aménagement du territoire} est l'action publique qui organise l'espace national pour corriger les déséquilibres, valoriser les potentialités de chaque région et assurer la cohésion sociale. Ses instruments : \textbf{schémas directeurs} (national, régionaux), \textbf{zonages} (agricoles, urbains, industriels, protégés), \textbf{infrastructures structurantes} (routes, barrages, fibre optique), \textbf{pôles de développement} (agropoles, zones économiques spéciales), \textbf{incitations fiscales} et \textbf{services publics décentralisés}.
\end{definition}

Trois défis majeurs lient villes et campagnes :

\begin{center}
\begin{tabularx}{\linewidth}{|l|X|X|}\hline
\textbf{Défi} & \textbf{Manifestation ville–campagne} & \textbf{Enjeu pour le développement} \\ \hline
Sécurité alimentaire & Dépendance aux importations (riz : 65 \% des besoins ; blé : 100 \% ; poisson : 40 \%) ; volatilité des prix ; malnutrition urbaine (15 \% enfants < 5 ans en surpoids/obésité, 29 \% retard de croissance) & Réorienter la production vers les filières d'import-substitution (riz, maïs, soja, lait, poisson) ; structurer les filières (stockage, transformation, normes) ; protéger les terres agricoles périurbaines \\ \hline
Emploi des jeunes & 60 \% des < 25 ans ; chômage urbain (jeunes diplômés) + sous-emploi rural (accès à la terre, financement) ; migration risquée & Créer des emplois \textbf{dans} les filières agricoles (production, transformation, logistique, services numériques) ; valoriser l'entrepreneuriat agropastoral ; former aux métiers verts \\ \hline
Aménagement du territoire & Macrocephalie (Douala–Yaoundé = 45 \% PIB, 30 \% pop.) ; déprise rurale (Est, Adamaoua, Nord) ; étalement urbain anarchique (bidonvilles, zones inondables) & Développer \textbf{pôles secondaires} (Bertoua, Ngaoundéré, Maroua, Bamenda, Bafoussam, Ebolowa, Kribi) ; corridors de développement (autoroute Yaoundé–Douala, route du Nord, corridor Kribi–Lolodorf) ; agropoles intégrées \\ \hline
\end{tabularx}
\end{center}

\begin{exemplebox}[Le paradoxe de l'Extrême-Nord : grenier et région la plus pauvre]
L'Extrême-Nord produit 40 \% du mil/sorgho national, 60 \% du coton, abrite le cheptel bovin le plus important. Pourtant, le taux de pauvreté y dépasse 74 \% (ECAM 4, 2014). Causes : \textbf{faible transformation locale} (coton exporté brut, bétail exporté sur pied), \textbf{enclavement} (routes dégradées, pas de rail), \textbf{chocs climatiques} (sécheresses récurrentes, inondations), \textbf{insécurité} (Boko Haram). L'interdépendance ville–campagne y est brisée : la valeur ajoutée fuit vers Douala (coton, bétail) ou l'étranger (Nigéria), sans retour d'investissement local.
\end{exemplebox}

\courssub{Les politiques de développement : décentralisation, agropoles, programmes ruraux}

Depuis 2010, le Cameroun déploie une stratégie territoriale articulée autour de trois piliers :

\begin{methode}[Stratégie territoriale intégrée : les 3 piliers]
\begin{enumerate}
\item \textbf{Décentralisation (Loi 2019/024) :} Transfert de compétences (agriculture, élevage, environnement, aménagement local, marchés) aux 360 communes et 10 régions. \textbf{Levier :} budgets communaux (Fonds d'équipement communal, fiscalité locale) pour infrastructures rurales (pistes, magasins de stockage, abattoirs, marchés).
\item \textbf{Agropoles (Agropoles de 2e génération, 2020–2030) :} Zones d'aménagement agricole intégrées associant \textbf{production}, \textbf{transformation}, \textbf{logistique} et \textbf{services} sur un même périmètre. Objectif : capter la valeur ajoutée près des zones de production. \textbf{5 agropoles prioritaires :} Nord (riz, maïs, coton), Adamaoua (lait, viande), Ouest (légumes, pomme de terre, poulet), Sud-Ouest (palmier, cacao, plantain), Est (cacao, manioc, poisson).
\item \textbf{Programmes ruraux ciblés :}
\begin{itemize}
\item \textbf{PNDP (Programme National de Développement Participatif) :} micro-projets communautaires (eau, pistes, écoles, centres de santé, magasins de stockage).
\item \textbf{PIA (Programme d'Investissement Agricole) :} intrants subventionnés, matériel, appui-conseil, financement (Fonds de développement agricole).
\item \textbf{PRODER (Programme de Développement Rural) :} bassins de production, pistes rurales, organisation professionnelle (coopératives, OP).
\item \textbf{Projet d'urgence pour la sécurité alimentaire (PUSA) :} riposte aux chocs (COVID, conflits, climatique).
\end{itemize}
\end{enumerate}
\end{methode}

\begin{savaistu}[Les agropoles du Cameroun : transformer près des champs]
Le concept d'agropole (agriculture + pôle) vise à briser le schéma « production rurale $\rightarrow$ transformation urbaine $\rightarrow$ consommation ». Exemple : l'\textbf{Agropole du Nord} (Garoua–Lagdo) intègre \SI{50000}{ha} aménagés (riz irrigué), une minoterie de \SI{200}{t/j}, une huilerie de coton, un abattoir moderne, un centre de formation, une zone logistique (gare routière, futur rail). Résultat attendu : 150 000 emplois directs/indirects à l'horizon 2030, substitution de \SI{200000}{t} de riz importé. Le défi : garantir l'accès à la terre pour les jeunes et les femmes, et la viabilité économique des unités de transformation (énergie, maintenance, marchés).
\end{savaistu}

\courssub{Complément utile au Bac : notion de développement durable appliquée aux villes africaines}

\begin{definition}[Développement durable (Brundtland, 1987) appliqué aux villes africaines]
Un développement \textbf{qui répond aux besoins du présent sans compromettre la capacité des générations futures à répondre aux leurs}. Pour une ville africaine en croissance rapide, cela signifie concilier \textbf{trois piliers} :
\begin{itemize}
\item \textbf{Économique :} productivité, emploi décent, finances locales saines, attractivité investissements.
\item \textbf{Social :} accès au logement, eau, assainissement, santé, éducation, transport pour tous ; inclusion des migrants ruraux, des jeunes, des femmes.
\item \textbf{Environnemental :} maîtrise de l'étalement (densification, corridors verts), gestion des déchets (économie circulaire), qualité de l'air, adaptation climatique (îlots de chaleur, inondations), préservation des terres agricoles périphériques.
\end{itemize}
\end{definition}

\begin{attention}[Piège fréquent au Bac : développement durable $\neq$ écologie seule]
Ne réduisez pas le développement durable à la protection de l'environnement. Une ville qui protège ses espaces verts mais où 60 \% de la population vit dans l'insalubrité sans emploi n'est \textbf{pas} durable. L'examinateur attend l'articulation des trois piliers. Exemple de réponse attendue : « Douala développe un BRT (Bus Rapid Transit) : \textit{économique} (productivité accrue, coûts transport réduits), \textit{social} (accès mobilité pour pauvres, gain de temps), \textit{environnemental} (baisse émissions CO$_2$, pollution atmosphérique). »
\end{attention}

\begin{exoresolu}[Analyse d'un document : extrait du Discours du Président de la République, 10 février 2023 — « L'agropole, levier de l'import-substitution »]
\textbf{Énoncé :} À partir du document ci-dessous, montrez comment la politique des agropoles traduit une vision de développement durable intégrant les dimensions économique, sociale et environnementale.

\textit{Document :} « Les agropoles que nous déployons dans nos cinq zones agroécologiques majeures ne sont pas de simples zones agricoles. Ce sont des écosystèmes productifs intégrés : irrigation maîtrisée, énergie solaire, unités de transformation de première et seconde génération, centres de formation aux métiers de l'agrobusiness, habitats groupés pour les jeunes agriculteurs, connectivité numérique. Elles visent à substituer nos importations alimentaires (riz, blé, poisson, lait, huile) par une production nationale compétitive, créatrice d'emplois décents pour notre jeunesse, respectueuse de nos sols et de notre eau. »

\textbf{Solution détaillée :}
\begin{enumerate}
\item \textbf{Identification du sujet :} Politique des agropoles (aménagement du territoire, sécurité alimentaire, emploi jeunes).
\item \textbf{Analyse par pilier du développement durable :}
\begin{itemize}
\item \textit{Économique :} « substituer nos importations » (balance commerciale, souveraineté), « production nationale compétitive » (productivité, valeur ajoutée locale), « agrobusiness » (entrepreneuriat, revenus).
\item \textit{Social :} « emplois décents pour notre jeunesse » (insertion professionnelle, lutte contre migration risquée), « centres de formation » (capital humain), « habitats groupés » (conditions de vie, fixation des populations).
\item \textit{Environnemental :} « irrigation maîtrisée » (gestion durable eau, adaptation climatique), « énergie solaire » (énergie renouvelable, bas carbone), « respectueuse de nos sols et de notre eau » (agroécologie, lutte contre dégradation).
\end{itemize}
\item \textbf{Synthèse :} L'agropole est un \textbf{dispositif territorial intégré} qui articule production, transformation, logistique, formation, habitat et énergie. Elle incarne l'interdépendance ville–campagne : la ville (capitaux, technologie, marché) investit la campagne pour y créer de la valeur partagée, dans une logique de circularité (énergie solaire $\rightarrow$ irrigation $\rightarrow$ production $\rightarrow$ transformation $\rightarrow$ emploi $\rightarrow$ revenus $\rightarrow$ demande locale).
\item \textbf{Conclusion argumentée :} La politique des agropoles traduit une vision de développement durable \textbf{si et seulement si} la gouvernance foncière garantit l'accès des jeunes/femmes, si les unités de transformation sont viables économiquement (énergie, maintenance, débouchés), et si les externalités environnementales (salinisation, biodiversité) sont monitorées. C'est un levier structurel, non une solution magique.
\end{enumerate}
\end{exoresolu}

\courssub{Bilan des savoir-faire : analyser un document, construire un croquis, rédiger une conclusion argumentée}

\begin{methode}[Rédiger une conclusion de composition géographique (Bac)]
\begin{enumerate}
\item \textbf{Rappeler la thèse (1 phrase) :} Répondez directement au sujet. Ex. : « L'interdépendance ville–campagne, si elle est structurée par des infrastructures et des politiques inclusives, est le levier principal du développement durable au Cameroun. »
\item \textbf{Faire le bilan synthétique (2–3 phrases) :} Reprenez les 2–3 arguments majeurs du développement sans les répéter mot pour mot. Ex. : « L'approvisionnement des villes repose sur la productivité rurale ; la transformation locale (agropoles) capte la valeur ajoutée ; la décentralisation donne aux territoires les moyens d'agir. »
\item \textbf{Ouvrir (1 phrase) :} Perspective, condition de réussite, comparaison, enjeu futur. Ex. : « Le défi reste la gouvernance foncière et la viabilité énergétique des agropoles pour que la boucle ville–campagne devienne vertueuse pour tous les territoires. »
\item \textbf{Ne JAMAIS :} introduire un nouvel argument, citer un nouvel exemple, faire de l'affectif (« j'espère que... »), poser une question rhétorique.
\end{enumerate}
\end{methode}

\begin{attention}[Au Bac, une conclusion ne doit pas introduire de notion nouvelle non traitée]
L'examinateur sanctionne toute idée qui n'a pas été développée dans le corps de la composition. Si vous n'avez pas parlé de « changement climatique » dans le développement, n'en parlez pas en conclusion. L'ouverture doit \textbf{prolonger} la réflexion (condition, perspective, échelle supérieure), pas l'élargir artificiellement. Exemple d'ouverture acceptée : « La réussite de cette interdépendance suppose une gouvernance métropolitaine associant communes urbaines et rurales. » (si la gouvernance a été évoquée). Exemple refusé : « Le changement climatique menace l'agriculture. » (si absent du développement).
\end{attention}

\begin{exercice}[Entraînement Bac : sujet type]
\textbf{Sujet :} « Les relations ville–campagne au Cameroun : atouts et défis pour la sécurité alimentaire et l'emploi des jeunes. »
\textbf{Consignes :} 
\begin{enumerate}
\item Proposez un plan détaillé (introduction, 2–3 parties équilibrées, conclusion).
\item Réalisez un croquis légendé (format A4) illustrant les flux ville–campagne et les agropoles.
\item Rédigez l'introduction (accroche, définition, sujet, thèse, annonce du plan) et la conclusion (thèse, bilan, ouverture).
\end{enumerate}
\end{exercice}

\begin{corrige}[Exercice n°1 — Éléments de correction]
\textbf{Plan suggéré :}
\begin{itemize}
\item \textbf{Introduction :} Accroche (urbanisation rapide + dépendance alimentaire) ; définitions (ville, campagne, sécurité alimentaire) ; thèse : « L'interdépendance ville–campagne est un atout majeur sous-exploité ; sa structuration par les agropoles et la décentralisation est la condition de la sécurité alimentaire et de l'emploi jeune. » ; plan en 3 parties.
\item \textbf{I. Des flux vitaux mais fragiles (atouts/faiblesses) :} A. Approvisionnement : diversité des bassins, complémentarité écologique. B. Vulnérabilités : logistique, pertes, chocs, insécurité.
\item \textbf{II. L'emploi jeune : le maillon manquant :} A. Poussée démographique + chômage urbain / sous-emploi rural. B. Potentiel des filières (production, transformation, logistique, services numériques). C. Obstacles : foncier, financement, formation, image.
\item \textbf{III. Politiques territoriales : structurer l'interdépendance :} A. Décentralisation : compétences, budgets, proximité. B. Agropoles : intégration production-transformation-services. C. Programmes ruraux : infrastructures, organisation, résilience.
\item \textbf{Conclusion :} Thèse rappelée ; bilan (productivité, transformation, gouvernance) ; ouverture (gouvernance métropolitaine, financement climat, jeunesse comme actrice).
\end{itemize}
\textbf{Croquis :} Carte du Cameroun (contour simplifié) : 5 agropoles (symboles usine+champ), 2 villes primaires (Douala, Yaoundé), 7 pôles secondaires, axes routiers structurants (flèches proportionnelles), bassins de production (hachures colorées), flux export (port de Douala/Kribi). Légende organisée (signes ponctuels, linéaires, surfaciques).
\textbf{Introduction type :} « En 2024, 58 \% des Camerounais vivent en ville, mais 70 \% de l'alimentation provient des campagnes. Cette dichotomie apparente masque une interdépendance profonde : les villes consomment ce que les campagnes produisent, et les campagnes ont besoin des marchés, capitaux et services urbains pour se moderniser. Pourtant, le pays importe 65 \% de son riz et 40 \% de son poisson, tandis que 60 \% des jeunes sont au chômage ou en sous-emploi. \textbf{Thèse :} L'interdépendance ville–campagne, si elle est structurée par des infrastructures, des agropoles intégrées et une décentralisation effective, constitue le levier principal pour assurer la sécurité alimentaire et créer des emplois décents pour la jeunesse. \textbf{Plan :} Nous analyserons d'abord les flux vitaux mais fragiles (I), puis le défi de l'emploi jeune dans les filières agropastorales (II), enfin les politiques territoriales qui tentent de structurer cette interdépendance (III). »
\textbf{Conclusion type :} « L'interdépendance ville–campagne est bien le moteur du développement camerounais : elle assure l'approvisionnement des métropoles, génère de la valeur dans les bassins ruraux, et porte les espoirs d'emploi pour la jeunesse. Les agropoles, la décentralisation et les programmes ruraux en sont les instruments concrets. \textbf{Condition sine qua non :} une gouvernance foncière équitable, une énergie fiable et bas-carbone, et une planification métropolitaine associant communes urbaines et rurales pour que la boucle de rétroaction — revenus ruraux $\rightarrow$ demande urbaine $\rightarrow$ investissements ruraux — devienne vertueuse sur l'ensemble du territoire. »
\end{corrige}

\begin{aretenir}[Synthèse finale : 5 idées-forces pour le Bac]
\begin{enumerate}
\item \textbf{Villes et campagnes ne sont pas deux mondes séparés} mais les deux pôles d'un même système : flux physiques (nourriture, matières), flux financiers (revenus, investissements), flux humains (migration, compétences).
\item \textbf{La sécurité alimentaire urbaine dépend de la productivité et de la connectivité rurales} (routes, stockage, transformation, information marchés).
\item \textbf{L'emploi des jeunes se joue dans l'articulation des filières :} production $\rightarrow$ transformation $\rightarrow$ logistique $\rightarrow$ services. L'agropole est le dispositif spatial qui les rassemble.
\item \textbf{L'aménagement du territoire (décentralisation, agropoles, corridors) est l'outil politique} pour rééquilibrer la macrocephalie Douala–Yaoundé et valoriser tous les bassins de production.
\item \textbf{Le développement durable urbain africain} exige l'articulation économique (productivité, emploi), sociale (inclusion, services) et environnementale (énergie propre, sols préservés, adaptation climatique) — pas l'écologie seule.
\end{enumerate}
\end{aretenir}

\newpage

\coursec{Activité d'intégration}

\coursec{Leçon 3 : Les villes — Activité d'intégration : Le cas de Makepe-Missoke (Douala)}

\courssub{Objectifs de l'activité}
\begin{itemize}
    \item Mobiliser les connaissances sur l'\cle{urbanisation} au Cameroun, les \cle{problèmes urbains}, les \cle{relations ville-campagne}, l'\cle{agriculture} et l'\cle{élevage} pour analyser une situation réelle.
    \item Exploiter un dossier documentaire hétérogène (photographie, tableau statistique, carte de localisation) pour en extraire l'information pertinente.
    \item Construire des représentations graphiques rigoureuses : diagramme en barres (statistiques urbaines) et \cle{croquis de synthèse} (flux ville-campagne avec \cle{flèches proportionnelles} et légende organisée).
    \item Rédiger une argumentation structurée (environ 10 lignes) proposant des solutions d'aménagement fondées sur l'analyse géographique.
    \item Travailler en équipe : répartir les tâches, confronter les interprétations, préparer une restitution orale collective.
\end{itemize}

\courssub{Situation}
\label{sit:makepe}
La ville de Douala, capitale économique du Cameroun, compte plus de \num{3,5} millions d'habitants (estimation 2023) et concentre près de 60~\% de l'activité industrielle nationale. Son extension spatiale rapide, souvent non maîtrisée, a engendré la prolifération de quartiers précaires en périphérie. Le quartier de \textbf{Makepe-Missoke}, situé dans le 3\textsuperscript{e} arrondissement, à l'est de l'aéroport international, illustre ces dynamiques.

\vspace{0,5em}
\textbf{Dossier documentaire fourni à chaque groupe :}

\begin{enumerate}
    \item \textbf{Document 1 : Photographie aérienne (extrait) de Makepe-Missoke.}
    \emph{Description pour l'exercice :} L'image montre un tissu urbain dense et anarchique. Les toitures en tôles ondulées dominent, serrées les unes contre les autres, sans alignement visible. Aucun espace vert ni équipement public (école, centre de santé) n'est identifiable dans le cadre. Les ruelles sont étroites, sinueuses, non bitumées, impraticables pour les véhicules de secours ou de collecte des ordures. En bordure, un cours d'eau (le fleuve Wouri ou un affluent) sert visiblement de réceptacle aux eaux usées et aux déchets solides. Quelques hangars artisanaux (menuiserie, soudure) s'intercalent entre les habitations.
    
    \item \textbf{Document 2 : Tableau de données socio-économiques et environnementales (année 2022).}
    \begin{center}
    \begin{tabularx}{\linewidth}{|l|X|}\hline
    \rowcolor{popblueL} \textbf{Indicateur} & \textbf{Valeur pour Makepe-Missoke} \\ \hline
    Population résidente & \num{42 500} habitants \\ \hline
    Taux d'accès à l'eau potable (branchement privé ou borne fontaine à < 200 m) & 18~\% \\ \hline
    Part de l'emploi informel dans la population active (15-64 ans) & 76~\% \\ \hline
    Production quotidienne de déchets solides ménagers & \num{18,5} tonnes \\ \hline
    Taux de collecte des déchets par les services municipaux (HYSACAM) & 22~\% \\ \hline
    Densité de population estimée & \num{480} hab./ha \\ \hline
    \end{tabularx}
    \end{center}
    \legende{Source : Enquête simulée inspirée des données INS, MINHDU, HYSACAM Douala 2022.}
    
    \item \textbf{Document 3 : Carte de localisation des bassins de production approvisionnant Douala (données 2022).}
    \emph{Données à reporter sur un fond de carte simplifié du Cameroun (fourni vierge) :}
    \begin{itemize}
        \item \textbf{Productions vivrières (tubercules, plantain, maïs) :} Région de l'Ouest (Bafoussam, Dschang, Mbouda) : \num{420 000} t/an ; Région du Littoral rural (Nkongsamba, Yabassi, Edéa) : \num{180 000} t/an ; Région du Sud-Ouest (Kumba, Mamfé) : \num{95 000} t/an.
        \item \textbf{Élevage (bovins, ovins, caprins, volailles) :} Région de l'Adamaoua (Ngaoundéré, Tibati) : \num{85 000} têtes équivalents bovins/an ; Région du Nord (Garoua, Poli) : \num{60 000} têtes équivalents bovins/an ; Péri-urbain de Douala (poulaillers industriels, porcheries) : \num{12 000} t/an de viandes blanches.
    \end{itemize}
    Les axes routiers principaux sont : Douala--Bafoussam (N5, bitumée, \num{280} km), Douala--Nkongsamba (N5, bitumée, \num{160} km), Douala--Yaoundé--Ngaoundéré (N3/N1, bitumée, \num{650} km jusqu'à Ngaoundéré). Le chemin de fer (Camrail) relie Douala à Ngaoundéré via Yaoundé (fret conteneurisé).
\end{enumerate}

\vspace{0,5em}
\textbf{Contexte de la mission :} Vous êtes une équipe de géographes consultants mandatée par la Communauté Urbaine de Douala (CUD) pour réaliser un \emph{diagnostic rapide} et proposer des \emph{pistes d'aménagement prioritaires} pour Makepe-Missoke, en intégrant la dimension « relations ville-campagne » (approvisionnement alimentaire, gestion des déchets, emploi).

\courssub{Consignes}
\textbf{Travail en groupe (4 élèves) — Durée : 2 séances (2h) + 1 séance restitution (1h).}

\begin{enumerate}
    \item \textbf{Analyse de la photographie (Document 1) :} Décrivez l'organisation de l'espace (morphologie, habitat, voirie, réseaux). Identifiez et nommez \textbf{trois problèmes urbains majeurs} visibles, en les reliant aux notions du cours (ex. : habitat précaire, insalubrité, absence d'équipements).
    \item \textbf{Traitement statistique (Document 2) :} 
    \begin{enumerate}
        \item Calculez le tonnage journalier de déchets \emph{non collectés} qui s'accumulent dans le quartier.
        \item Construisez \textbf{deux diagrammes en barres} distincts comparant les six indicateurs du tableau : un pour les trois taux (en \%), un pour les trois grandeurs absolues (valeurs brutes, unités adaptées). Titre, unités, source obligatoires.
    \end{enumerate}
    \item \textbf{Croquis de synthèse : Relations ville-campagne (Document 3 + cours) :} Sur le fond de carte du Cameroun fourni (A4, contour simplifié, villes repères : Douala, Yaoundé, Bafoussam, Ngaoundéré, Garoua, Nkongsamba, Kumba), réalisez un croquis montrant les \textbf{flux d'approvisionnement vers Douala}.
    \begin{itemize}
        \item Tracez des \textbf{flèches proportionnelles} (largeur $\propto$ racine carrée du tonnage ou nombre de têtes) pour les flux vivriers (vert) et les flux d'élevage (rouge).
        \item Représentez les deux axes de transport principaux (route N5, route/rail N3-N1) par des traits distincts.
        \item Construisez une \textbf{légende organisée} intégrée au croquis (titre, signatures graphiques, échelons de proportionnalité, auteur, date, source).
    \end{itemize}
    \item \textbf{Synthèse écrite (individuelle, 10 lignes max) :} À partir de l'ensemble de vos analyses, rédigez une conclusion argumentée proposant \textbf{deux solutions d'aménagement prioritaires et justifiées} pour Makepe-Missoke. L'une doit relever de la \textbf{gestion urbaine interne} (voirie, eau, déchets, habitat), l'autre de la \textbf{valorisation des relations ville-campagne} (emploi, filières agricoles, gestion des déchets organiques).
    \item \textbf{Restitution orale (5 min/groupe) :} Présentez le croquis et la synthèse. Débat collectif : les solutions sont-elles compatibles ? Durables ? Finançables ?
\end{enumerate}

\courssub{Ressources à mobiliser}

\begin{definition}[Notions clés du cours]
\textbf{\cle{Urbanisation} / \cle{Ville} :} Concentration de population, d'activités, de fonctions de commande. \textbf{\cle{Problèmes urbains} :} \cle{Habitat précaire} (bidonvilles), déficit d'infrastructures (eau, assainissement, électricité, voirie), chômage/sous-emploi informel, gestion des déchets, pollution, risques (inondation, érosion). \textbf{\cle{Relations ville-campagne} :} Flux matériels (produits agricoles vers la ville, intrants/équipements vers la campagne), flux humains (migration, navettage), flux financiers (transferts, investissements). \textbf{\cle{Agriculture} camerounaise :} Vivrière (tubercules, plantain, maïs) dominant dans l'Ouest, le Littoral, le Sud-Ouest ; rentes (cacao, café, coton, banane, palme à huile). \textbf{\cle{Élevage} :} Bovins au Nord/Adamaoua (systèmes extensifs) ; volailles/porcs en péri-urbain (systèmes intensifs).
\end{definition}

\begin{methode}[Construire un diagramme en barres]
\begin{enumerate}
    \item Choisir le type : barres verticales (histogramme) pour comparer des grandeurs de même nature.
    \item Définir l'échelle : valeur max $\rightarrow$ hauteur max (ex. 10 cm). Calculer le coefficient : $k = \frac{H_{max}}{V_{max}}$.
    \item Tracer les axes : abscisses = indicateurs (qualitatifs), ordonnées = valeurs (quantitatives, unité).
    \item Dessiner les barres : même largeur, espacement régulier, couleur unie ou nuancée par catégorie.
    \item Légender : titre précis, source, unité, échelle si nécessaire.
\end{enumerate}
\end{methode}

\begin{methode}[Réaliser un croquis de flux (flèches proportionnelles)]
\begin{enumerate}
    \item Choisir la variable à représenter (tonnage, valeur, nombre de têtes).
    \item Définir l'échelle de proportionnalité : largeur de flèche $l = k \times \sqrt{V}$ (racine carrée pour garder des largeurs lisibles) ou $l = k \times V$ (linéaire si écart faible). Exemple : $1 \text{ mm} \leftrightarrow \sqrt{10 000} = 100$ t $\rightarrow$ $k = 0,1$ mm/$\sqrt{\text{t}}$.
    \item Tracer les flèches du lieu de production $\rightarrow$ lieu de consommation (Douala). Courber légèrement pour éviter les chevauchements.
    \item Différencier les natures de flux par la couleur (vert vivrier, rouge élevage) ou le motif.
    \item Légende obligatoire intégrée : titre, signature graphique (ex. \tikz\draw[popgreen,line width=2pt] (0,0)--(1,0); = 100 000 t/an), auteur, date, source.
\end{enumerate}
\end{methode}

\begin{attention}[Pièges fréquents]
\begin{itemize}
\item Ne pas confondre \emph{taux d'accès à l'eau} (pourcentage) et \emph{volume d'eau consommé} (m$^3$). Ne pas mélanger pourcentages et tonnes sur un même diagramme sans double échelle claire (deux graphiques distincts sont préférables).
\item Flèches proportionnelles : la largeur doit refléter l'ordre de grandeur, pas la valeur exacte au pixel près. Vérifier que la plus grosse flèche n'écrase pas la carte.
\item Conclusion : « Construire des routes » ou « Donner de l'eau » ne sont pas des réponses géographiques argumentées. Il faut citer les acteurs (CUD, MINHDU, HYSACAM, MINADER, MINEPIA, ONG, population), les mécanismes (participatif, foncier, filière), le lien avec le diagnostic (ex. : déchets organiques $\rightarrow$ compost $\rightarrow$ agriculture péri-urbaine).
\end{itemize}
\end{attention}

\courssub{Démarche et corrigé commenté}

\begin{exoresolu}[Corrigé type de l'activité d'intégration — Makepe-Missoke]

\textbf{Étape 1 : Analyse de la photographie (Document 1)}
\begin{itemize}
    \item \textbf{Morphologie :} Tissu dense, parcellaire irrégulier, absence de trame viaire structurante. Habitat spontané, matériaux de récupération (tôles, bois), mitoyenneté forte.
    \item \textbf{Problème 1 — \cle{Habitat précaire} et insécurité foncière :} Constructions non conformes aux normes (POS/PLU), risque d'effondrement, d'incendie (ruelles étroites), absence de titres fonciers $\rightarrow$ frein à l'investissement locataire/propriétaire.
    \item \textbf{Problème 2 — Déficit d'infrastructures et de services urbains :} Voirie inexistante (non bitumée, pas de caniveaux) $\rightarrow$ inaccessibilité secours/collecte, érosion hydrique, stagnation eaux pluviales/usées. Absence d'équipements publics de proximité (école, santé, marché structuré).
    \item \textbf{Problème 3 — \cle{Insalubrité} environnementale et pollution :} Rejet direct eaux usées/déchets solides dans le cours d'eau $\rightarrow$ pollution aval, risques sanitaires (choléra, paludisme), nuisances olfactives. Décharges sauvages visibles en lisière.
\end{itemize}
\emph{Justification :} Ces trois problèmes recouvrent les trois piliers du « problème urbain » au programme : habitat, équipements/réseaux, environnement/santé.

\textbf{Étape 2 : Traitement statistique (Document 2)}
\begin{itemize}
    \item \textbf{Calcul déchets non collectés :} Production = 18,5 t/j. Collecte = 22~\% $\rightarrow$ $18,5 \times 0,22 = 4,07$ t/j collectées. Non collectées = $18,5 - 4,07 = \mathbf{14,43 \text{ t/j}}$ (soit 78~\% du volume). Cela représente $\approx \num{5 267}$ t/an qui s'accumulent dans le quartier ou sont brûlés/jettés dans le cours d'eau.
    \item \textbf{Diagrammes en barres (deux graphiques distincts) :}
\end{itemize}

\begin{popfigure}
\begin{tikzpicture}[scale=0.9]
% Graphique A : Taux en %
\draw[popline, thick] (0,0) -- (7,0) node[right] {Indicateurs (taux)};
\draw[popline, thick] (0,0) -- (0,8.5) node[above] {\%};
\foreach \i/\val/\label in {1/18/Eau, 2/76/Informel, 3/22/Collecte} {
    \fill[popblue!70] (\i*1.8-0.4,0) rectangle (\i*1.8+0.4,\val/10);
    \node[below, font=\small] at (\i*1.8,-0.3) {\label};
    \node[above, font=\small] at (\i*1.8,\val/10+0.15) {\val\%};
}
\node[font=\small, align=center] at (3.5, -1.2) {\textbf{Graphique A : Taux de précarité à Makepe-Missoke (2022)}};
\node[font=\tiny] at (3.5, -1.6) {Source : INS/MINHDU/HYSACAM};
\end{tikzpicture}
\legende{Graphique A : Trois indicateurs de taux (accès à l'eau, emploi informel, collecte déchets). Échelle : 1 cm pour 10~\%.}
\end{popfigure}

\begin{popfigure}
\begin{tikzpicture}[scale=0.9]
% Graphique B : Grandeurs absolues (échelles adaptées par indicateur)
\draw[popline, thick] (0,0) -- (7,0) node[right] {Indicateurs (valeurs absolues)};
\draw[popline, thick] (0,0) -- (0,5.5) node[above] {Valeurs (échelles propres)};
% Population (42500 hab -> hauteur 4.25 cm, échelle 1 cm = 10 000 hab)
\fill[poporange!70] (1.8-0.4,0) rectangle (1.8+0.4,4.25);
\node[below, font=\small] at (1.8,-0.3) {Population};
\node[above, font=\small] at (1.8,4.4) {42 500 hab};
% Déchets (18.5 t -> hauteur 3.7 cm, échelle 1 cm = 5 t)
\fill[poporange!70] (3.6-0.4,0) rectangle (3.6+0.4,3.7);
\node[below, font=\small] at (3.6,-0.3) {Déchets/j};
\node[above, font=\small] at (3.6,3.85) {18,5 t};
% Densité (480 hab/ha -> hauteur 4.8 cm, échelle 1 cm = 100 hab/ha)
\fill[poporange!70] (5.4-0.4,0) rectangle (5.4+0.4,4.8);
\node[below, font=\small] at (5.4,-0.3) {Densité};
\node[above, font=\small] at (5.4,4.95) {480 hab/ha};
\node[font=\small, align=center] at (3.5, -1.2) {\textbf{Graphique B : Grandeurs absolues (échelles distinctes par indicateur)}};
\node[font=\tiny] at (3.5, -1.6) {Source : INS/MINHDU/HYSACAM};
\end{tikzpicture}
\legende{Graphique B : Trois indicateurs en valeurs absolues. Chaque barre a sa propre échelle (indiquée en légende), nécessaire car les unités et ordres de grandeur diffèrent.}
\end{popfigure}

\textbf{Étape 3 : Croquis de synthèse — Flux ville-campagne vers Douala}
\begin{itemize}
    \item \textbf{Choix d'échelle (racine carrée pour lisibilité) :} 
    \begin{align*}
        \text{Vivriers : } & \text{Ouest } 420\,000 \rightarrow \sqrt{420000} \approx 648 ; \text{ Littoral } 180\,000 \rightarrow 424 ; \text{ Sud-Ouest } 95\,000 \rightarrow 308. \\
        \text{Élevage : } & \text{Adamaoua } 85\,000 \rightarrow 292 ; \text{ Nord } 60\,000 \rightarrow 245 ; \text{ Péri-urbain } 12\,000 \rightarrow 110.
    \end{align*}
    Coefficient $k = 0,003$ cm/$\sqrt{\text{t}}$ $\rightarrow$ largeurs : Ouest 1,94 cm, Littoral 1,27 cm, Sud-Ouest 0,92 cm, Adamaoua 0,88 cm, Nord 0,74 cm, Péri-urbain 0,33 cm. Arrondir au 0,5 mm pour le tracé.
    \item \textbf{Fond de carte :} Contour Cameroun simplifié (polygone 8-10 points). Villes repères : Douala (côte), Yaoundé (centre), Bafoussam (Ouest), Nkongsamba (Littoral), Ngaoundéré (Adamaoua), Garoua (Nord), Kumba (Sud-Ouest).
    \item \textbf{Flèches :} 
    \begin{itemize}
        \item Vert (vivriers) : Bafoussam $\rightarrow$ Douala (large, N5), Nkongsamba $\rightarrow$ Douala (moyenne, N5), Kumba $\rightarrow$ Douala (fine, route côtière).
        \item Rouge (élevage) : Ngaoundéré $\rightarrow$ Yaoundé $\rightarrow$ Douala (moyenne, rail/route N1-N3), Garoua $\rightarrow$ Ngaoundéré (fine), Péri-urbain $\rightarrow$ Douala (très fine, boucle locale).
    \end{itemize}
    \item \textbf{Transport :} Route N5 (trait continu épais popblue), Route N3/N1 (trait continu popblue), Rail Camrail (trait pointillé poppurple).
\end{itemize}

\begin{popfigure}
\begin{tikzpicture}[scale=0.75]
% Contour Cameroun très simplifié
\draw[popdark, thick] (0,0) -- (8,0) -- (9,2) -- (9,6) -- (7,8) -- (5,7) -- (3,7) -- (1,5) -- (0,3) -- cycle;
% Villes repères (nodes pour coordonnées flèches)
\node[popblue, circle, fill=popblue, inner sep=2pt] (D) at (1.5,1.2) {};
\node[popdark, circle, fill=popdark, inner sep=1.5pt] (Y) at (4.5,3.5) {};
\node[popgreen, circle, fill=popgreen, inner sep=1.5pt] (B) at (6.5,4.5) {};
\node[popgreen, circle, fill=popgreen, inner sep=1.5pt] (Nk) at (5.5,2.5) {};
\node[poppink, circle, fill=poppink, inner sep=1.5pt] (Ng) at (5.5,6.5) {};
\node[poppink, circle, fill=poppink, inner sep=1.5pt] (Ga) at (7.5,7.2) {};
\node[popgreen, circle, fill=popgreen, inner sep=1.5pt] (Ku) at (3.0,1.5) {}; % Kumba approx
% Noms des villes
\node[below, font=\tiny] at (D) {Douala};
\node[left, font=\tiny] at (Y) {Yaoundé};
\node[above, font=\tiny] at (B) {Bafoussam};
\node[right, font=\tiny] at (Nk) {Nkongsamba};
\node[above, font=\tiny] at (Ng) {Ngaoundéré};
\node[right, font=\tiny] at (Ga) {Garoua};
\node[below left, font=\tiny] at (Ku) {Kumba};
% Flux vivriers (vert) - Largeurs : 1.94, 1.27, 0.92 cm -> 19.4, 12.7, 9.2 pt (1cm=10pt approx? Non, TikZ pts. 1cm = 28.45pt. Donc 1.94cm = 55pt. Trop gros pour line width. On utilise mm.)
% line width en mm: 1.94cm = 19.4mm. TikZ line width en mm.
\draw[popgreen, line width=1.9mm, -{Stealth[length=4mm]}] (B) .. controls (5,3) .. (D);
\draw[popgreen, line width=1.3mm, -{Stealth[length=3.5mm]}] (Nk) .. controls (4,2) .. (D);
\draw[popgreen, line width=0.9mm, -{Stealth[length=3mm]}] (Ku) .. controls (2.5,1.5) .. (D);
% Flux élevage (rose/rouge) - Largeurs : 0.88, 0.74, 0.33 cm -> 8.8, 7.4, 3.3 mm
\draw[poppink, line width=0.9mm, -{Stealth[length=3.5mm]}] (Ng) .. controls (4.5,5) and (3,3) .. (Y);
\draw[poppink, line width=0.7mm, -{Stealth[length=2.5mm]}] (Ga) -- (Ng);
\draw[poppink, line width=0.3mm, -{Stealth[length=2mm]}] (2.5,1.5) -- (D); % Péri-urbain approx
% Transport
\draw[popblue, thick] (D) -- (Nk) -- (B); % N5
\draw[popblue, thick] (D) -- (Y); % N3
\draw[popblue, thick] (Y) -- (Ng); % N1 (route)
\draw[poppurple, thick, dashed] (Y) -- (Ng); % Rail (superposé route)
% Légende intégrée (coin bas gauche)
\begin{scope}[shift={(-0.5,-0.5)}]
\node[anchor=south west, font=\small, text=popdark, fill=popcream, fill opacity=0.8, text opacity=1, inner sep=4pt, rounded corners, align=center] (legbox) at (0,0){
    \begin{tabular}{ll}
    \textbf{Douala : Flux d'approvisionnement (2022)} & \\
    \tikz\draw[popgreen,line width=1.9mm] (0,0)--(1,0); & Flux vivriers (larg. $\propto \sqrt{\text{tonnage}}$) \\
    \tikz\draw[poppink,line width=0.9mm] (0,0)--(1,0); & Flux élevage (larg. $\propto \sqrt{\text{têtes équiv.}}$) \\
    \tikz\draw[popblue,thick] (0,0)--(1,0); & Route bitumée (N5, N3/N1) \\
    \tikz\draw[poppurple,thick,dashed] (0,0)--(1,0); & Chemin de fer (Camrail) \\
    \multicolumn{2}{l}{\tikz\draw[popgreen,line width=1.9mm] (0,0)--(0.5,0); = 400 000 t/an (vivriers)} \\
    \multicolumn{2}{l}{\tikz\draw[poppink,line width=0.9mm] (0,0)--(0.5,0); = 85 000 têtes équiv. (élevage)} \\
    \multicolumn{2}{l}{Auteur : Élève 1re Tech \quad Date : 2025} \\
    \multicolumn{2}{l}{Source : MINADER, MINEPIA, INS}
    \end{tabular}
};
\end{scope}
% Flèche Nord
\draw[popmuted, ->] (0.5,7) -- node[above, font=\tiny, rotate=45] {Nord} (1,7.5);
\end{tikzpicture}
\legende{Croquis de synthèse : Flux d'approvisionnement vivrier (flèches vertes) et d'élevage (flèches roses) vers Douala. Largeurs proportionnelles à la racine carrée des tonnages/têtes. Routes bitumées (bleu), chemin de fer (violet pointillé). Légende intégrée avec échelons.}
\end{popfigure}

\textbf{Étape 4 : Synthèse écrite — Proposition de deux solutions (modèle)}

\begin{quote}
\textbf{Diagnostic :} Makepe-Missoke concentre les maux de l'\cle{urbanisation} non maîtrisée : habitat précaire (480 hab/ha), exclusion des réseaux (18~\% eau), économie de survie (76~\% informel), insalubrité critique (14,4 t/j de déchets non collectés polluant le Wouri).

\textbf{Solution 1 — Aménagement interne : \cle{Opération de résorption programmée (ORP)} participative.} Sous maîtrise d'ouvrage CUD/MINHDU, avec financement BAD/Banque Mondiale (projet PDVIR). Actions : (1) Régularisation foncière (certificats d'occupation) pour sécuriser les occupants et inciter l'auto-réhabilitation ; (2) Ouverture de voiries structurantes (6 m) et secondaires (4 m) avec caniveaux, permettant l'accès camions HYSACAM et pompiers ; (3) Extension du réseau Eau (bornes fontaines tous les 200 m) et électricité (branchements sociaux) ; (4) Mise en place d'un tri à la source (bacs biodéchets/recyclables) géré par des \cle{GIC} de jeunes formés, créant de l'emploi formel. \emph{Justification :} S'attaque aux 3 problèmes identifiés (habitat, réseaux, déchets) par une approche intégrée, reproductible (ex. quartier Ndogpassi).

\textbf{Solution 2 — Valorisation ville-campagne : Filière « \cle{Déchets organiques} $\rightarrow$ \cle{Compost} $\rightarrow$ \cle{Agriculture péri-urbaine} ».} Les 14,4 t/j de biodéchets (60~\% des déchets) sont une ressource. Création d'une plateforme de compostage semi-industriel (1 ha) en zone inondable inconstructible en bordure de Makepe, gérée par une coopérative (femmes, jeunes). Compost certifié vendu aux maraîchers du Littoral rural (Nkongsamba, Yabassi) et aux poulaillers péri-urbains (Document 3 : 12 000 t/an viandes blanches). Retours : revenus pour le quartier, réduction flux déchets vers décharge, amélioration sols/sécurité alimentaire locale. \emph{Justification :} Boucle l'\cle{économie circulaire}, lie le problème urbain (déchets) à la solution rurale (intrant agricole), renforce l'approvisionnement de Douala (flux vert sur croquis), crée de l'emploi non délocalisable.
\end{quote}

\textbf{Étape 5 : Restitution orale — Points d'attention pour le débat}
\begin{itemize}
    \item \textbf{Compatibilité :} La voirie (Sol 1) dessert la plateforme compostage (Sol 2). Le tri à la source (Sol 1) alimente le compostage (Sol 2). Cohérence forte.
    \item \textbf{Durabilité :} Nécessite une gouvernance locale (comité de quartier, GIC) et des contrats d'approvisionnement pluriannuels avec les agriculteurs/éleveurs (Sol 2). Risque : fluctuation prix engrais chimiques $\rightarrow$ compétitivité compost.
    \item \textbf{Financement :} Sol 1 : investissement lourd (voirie, réseaux) $\rightarrow$ dette publique/partenaires. Sol 2 : investissement modéré (plateforme, broyeurs) $\rightarrow$ fonds climat/\cle{économie circulaire} (Fonds Vert, AFD) + autofinancement par vente compost.
    \item \textbf{Échelle :} Makepe-Missoke = 42 500 hab. Solutions doivent être pensées à l'échelle des 150+ quartiers précaires de Douala (massification).
\end{itemize}
\end{exoresolu}

\courssub{Bilan de l'activité}
Cette activité d'intégration a permis de mettre en évidence plusieurs acquis majeurs du programme de Première Technique :

\begin{itemize}
    \item \textbf{Articulation des échelles :} Les élèves sont passés de l'échelle locale (quartier, photographie, données micro) à l'échelle régionale/nationale (flux d'approvisionnement, carte du Cameroun), comprenant que les problèmes urbains ne se résolvent pas \emph{in situ} seulement, mais en interaction avec l'espace rural productif.
    \item \textbf{Transversalité des savoirs :} La mobilisation simultanée des chapitres « Problèmes urbains », « Relations ville-campagne », « Agriculture » et « Élevage » a montré leur interdépendance concrète : l'insalubrité urbaine (déchets) peut devenir ressource agricole (compost) ; l'emploi informel urbain peut se structurer en filières agricoles rurales-périurbaines.
    \item \textbf{Rigueur graphique et quantitative :} La construction des diagrammes en barres (choix d'échelles, double représentation taux/absolu) et du croquis à flèches proportionnelles (calcul racine carrée, légende normalisée) a exercé des savoir-faire techniques exigeants (mathématiques appliquées, sémiologie graphique) évaluables au \cle{Probatoire}.
    \item \textbf{Argumentation géographique :} La rédaction de la synthèse (10 lignes, 2 solutions justifiées) a contraint à la concision, à l'usage du vocabulaire précis (ORP, foncier, GIC, économie circulaire, intrants, sécurité alimentaire) et à la justification par les données du dossier (chiffres à l'appui).
    \item \textbf{Posture professionnelle :} Le jeu de rôle « consultants pour la CUD » a situ l'élève en situation de responsabilité, d'expertise et de communication orale, préparant aux attendus de l'insertion professionnelle (BTS, DUT, concours administratifs) et à la citoyenneté active.
\end{itemize}

\vspace{0,5em}
\noindent \textbf{Prolongements possibles (Projets d'établissement / Année supérieure) :} Étude d'impact environnemental simplifiée de la plateforme de compostage ; Simulation de concertation publique (rôles : riverains, CUD, HYSACAM, maraîchers, ONG) ; Comparaison avec d'autres métropoles africaines (Abidjan, Lagos, Dakar) via le réseau « Villes Africaines » de l'ONU-Habitat.

\newpage

\coursec{Méthodes et exercices résolus}

\courssub{Méthodes et exercices résolus}

\begin{methode}[Analyser l'évolution de l'urbanisation à partir de données statistiques]
\textbf{Objectif :} Calculer le taux d'urbanisation, le taux de croissance urbaine et interpréter l'explosion urbaine camerounaise.
\textbf{Étapes :}
\begin{enumerate}
    \item \textbf{Identifier les données :} Population totale ($P_{tot}$), population urbaine ($P_{urb}$), population rurale ($P_{rur}$) à deux dates ($t_0$ et $t_1$).
    \item \textbf{Calculer le taux d'urbanisation (TU) :} $TU = \frac{P_{urb}}{P_{tot}} \times 100$ (en \%). Il mesure le poids de la ville dans le pays.
    \item \textbf{Calculer le taux de croissance annuel moyen (TCAM) urbain :} $TCAM = \left( \sqrt[n]{\frac{P_{urb}(t_1)}{P_{urb}(t_0)}} - 1 \right) \times 100$ avec $n$ = nombre d'années.
    \item \textbf{Comparer :} Confronter TCAM urbain et TCAM national/rural. Si TCAM urbain $>$ TCAM national $\rightarrow$ urbanisation rapide (exode rural, croissance naturelle forte).
    \item \textbf{Conclure :} Lier les chiffres aux causes (exode rural, extension administrative, attractivité économique) et conséquences (pression sur l'habitat, l'emploi, les services).
\end{enumerate}
\textbf{Réflexe Bac :} Toujours préciser l'unité (\% pour TU, \%/an pour TCAM). Ne pas confondre \emph{taux d'urbanisation} (niveau à un instant T) et \emph{vitesse d'urbanisation} (variation du TU dans le temps).
\end{methode}

\begin{exoresolu}[Évolution de l'urbanisation au Cameroun (1987-2005)]
\textbf{Document :} Recensements généraux de la population et de l'habitat (RGPH 1987 et 2005, BUCREP).
\begin{center}
\begin{tabularx}{\linewidth}{|l|X|X|}\hline
\textbf{Indicateur} & \textbf{1987 (RGPH)} & \textbf{2005 (RGPH)} \\ \hline
Population totale & 10 493 655 & 17 463 836 \\ \hline
Population urbaine & 4 000 000 (38,1\%) & 9 100 000 (52,1\%) \\ \hline
Population rurale & 6 493 655 & 8 363 836 \\ \hline
\end{tabularx}
\end{center}
\textbf{Questions :}
\begin{enumerate}
    \item Vérifier le taux d'urbanisation pour 1987 et 2005.
    \item Calculer le TCAM de la population urbaine sur la période 1987-2005.
    \item Calculer le TCAM de la population totale sur la même période.
    \item Interpréter l'écart entre les deux TCAM et en déduire la dynamique de l'urbanisation au Cameroun.
\end{enumerate}
\tcblower
\textbf{Solution détaillée :}
\begin{enumerate}
    \item \textbf{Taux d'urbanisation (TU) :}
    \begin{itemize}
        \item 1987 : $TU_{1987} = \frac{4\,000\,000}{10\,493\,655} \times 100 \approx \textbf{38,1\%}$ (vérifié).
        \item 2005 : $TU_{2005} = \frac{9\,100\,000}{17\,463\,836} \times 100 \approx \textbf{52,1\%}$ (vérifié).
    \end{itemize}
    \item \textbf{TCAM population urbaine (1987-2005) :}
    $n = 2005 - 1987 = 18$ ans.
    \[TCAM_{urb} = \left( \sqrt[18]{\frac{9\,100\,000}{4\,000\,000}} - 1 \right) \times 100 = \left( \sqrt[18]{2,275} - 1 \right) \times 100\]
    $\ln(2,275) \approx 0,8222$ ; $\frac{0,8222}{18} \approx 0,04568$ ; $e^{0,04568} \approx 1,0467$
    \[TCAM_{urb} \approx (1,0467 - 1) \times 100 = \textbf{4,7\% par an}.\]
    \item \textbf{TCAM population totale (1987-2005) :}
    \[TCAM_{tot} = \left( \sqrt[18]{\frac{17\,463\,836}{10\,493\,655}} - 1 \right) \times 100 = \left( \sqrt[18]{1,664} - 1 \right) \times 100\]
    $\ln(1,664) \approx 0,5094$ ; $\frac{0,5094}{18} \approx 0,02830$ ; $e^{0,02830} \approx 1,0287$
    \[TCAM_{tot} \approx (1,0287 - 1) \times 100 = \textbf{2,9\% par an}.\]
    \item \textbf{Interprétation :}
    Le TCAM urbain (\textbf{4,7\%}) est nettement supérieur au TCAM total (\textbf{2,9\%}) et donc au TCAM rural (inférieur à 2,9\%).
    \textbf{Conclusion :} L'urbanisation au Cameroun est \textbf{rapide et accélérée}. Elle résulte d'un double phénomène : une forte \emph{croissance naturelle} en ville (fécondité encore élevée, baisse de la mortalité) et un \emph{exode rural} massif (recherche d'emploi et de services, extension administrative des communes). Le pays a basculé vers une majorité urbaine : le seuil des 50\% a été franchi au début des années 2000 (52,1\% au RGPH 2005).
\end{enumerate}
\end{exoresolu}

\begin{methode}[Caractériser et hiérarchiser le réseau urbain camerounais]
\textbf{Objectif :} Identifier les métropoles, pôles régionaux et villes moyennes ; expliquer leur fonction et leur rayonnement.
\textbf{Étapes :}
\begin{enumerate}
    \item \textbf{Classer par taille (population) :}
    \begin{itemize}
        \item \textbf{Métropoles (poids lourds) :} Douala (capitale économique, port, environ 3,5 à 4 millions d'habitants), Yaoundé (capitale politique et administrative, environ 3 à 3,5 millions d'habitants). Fonctions : commandement, décision, commerce international, industrie.
        \item \textbf{Pôles régionaux (chefs-lieux de région) :} Bamenda, Bafoussam, Garoua, Maroua, Ngaoundéré, Bertoua, Ebolowa, Buea. Fonctions : administration décentralisée, commerce de gros, services de santé et d'éducation supérieurs, relais vers les campagnes.
        \item \textbf{Villes moyennes / secondaires :} chefs-lieux de département, villes carrefours (ex : Foumban, Dschang, Kumba, Limbe, Yagoua, Kousseri). Fonctions : marché hebdomadaire, administration de base, services de proximité.
    \end{itemize}
    \item \textbf{Analyser la primatie :} Calculer l'indice de primatie (population de la 1\iere{} ville / population de la 2\ieme{} ville). Si $> 2$ $\rightarrow$ primatie forte (déséquilibre). Au Cameroun : Douala/Yaoundé $\approx 1$ (bipôle), mais Douala/3\ieme{} ville $> 5$ $\rightarrow$ macrocéphalie du couple dominant.
    \item \textbf{Localiser sur croquis :} Positionner les chefs-lieux de région + Douala/Yaoundé. Tracer les axes structurants (Transcamerounais, axe Douala-Yaoundé, axes Nord-Sud). Montrer la concentration littorale et de l'Ouest face au peuplement urbain plus lâche de l'Est et de l'Extrême-Nord.
    \item \textbf{Justifier la hiérarchie :} Histoire (colonisation, capitales), géographie physique (relief, climat, ressources), infrastructures (port, rail, routes, aéroports), politique (décentralisation).
\end{enumerate}
\textbf{Formule clé :} Indice de primatie $I_p = \frac{P_1}{P_2}$.
\end{methode}

\begin{exoresolu}[Hiérarchie urbaine et macrocéphalie au Cameroun]
\textbf{Données (estimations) :} Douala : 3 800 000 hab. ; Yaoundé : 3 500 000 hab. ; Bamenda : 500 000 hab. ; Garoua : 450 000 hab. ; Bafoussam : 400 000 hab.
\textbf{Questions :}
\begin{enumerate}
    \item Calculer l'indice de primatie de Douala sur Yaoundé ($I_{p1}$) et sur Bamenda ($I_{p2}$).
    \item Qualifier la structure du réseau urbain (primatie, bipôle, équilibré).
    \item Expliquer pourquoi Douala et Yaoundé concentrent plus de la moitié de la population urbaine totale alors qu'elles ne sont que deux villes parmi des dizaines.
    \item Sur un croquis simplifié du Cameroun, positionner les 5 villes ci-dessus et tracer l'axe routier Douala-Yaoundé-Bafoussam-Bamenda-Garoua.
\end{enumerate}
\tcblower
\textbf{Solution détaillée :}
\begin{enumerate}
    \item \textbf{Indices de primatie :}
    \begin{align*}
        I_{p1} (\text{Douala/Yaoundé}) &= \frac{3\,800\,000}{3\,500\,000} \approx \textbf{1,09} \\
        I_{p2} (\text{Douala/Bamenda}) &= \frac{3\,800\,000}{500\,000} = \textbf{7,6}
    \end{align*}
    \item \textbf{Structure du réseau :}
    $I_{p1} \approx 1$ $\rightarrow$ \textbf{bipôle fort} (Douala et Yaoundé de taille comparable).
    $I_{p2} > 2$ (et même $> 5$) $\rightarrow$ \textbf{macrocéphalie} marquée du couple dominant par rapport au troisième rang. Le réseau est \textbf{déséquilibré} : deux \og têtes \fg{} géantes, puis un grand vide avant le troisième échelon.
    \item \textbf{Explication de la concentration (Douala + Yaoundé $\approx$ 7,3 M sur environ 13 M d'urbains, soit plus de 55\%) :}
    \begin{itemize}
        \item \textbf{Histoire :} Douala (premier port, première capitale allemande) ; Yaoundé (capitale administrative coloniale, confirmée à l'indépendance).
        \item \textbf{Économie :} Douala = Port autonome (l'essentiel du trafic maritime national), zone industrielle (Bassa, Bonabéri), siège des banques et entreprises. Yaoundé = siège des institutions (Présidence, ministères, Assemblée nationale, ambassades), fonction publique massive, tertiaire supérieur.
        \item \textbf{Infrastructures :} aéroports internationaux, nœuds routiers et ferroviaires (Transcamerounais), universités, grands hôpitaux.
        \item \textbf{Effet d'entraînement :} attraction des migrants (exode rural), concentration des investissements publics et privés $\rightarrow$ cercle cumulatif de la concentration.
    \end{itemize}
    \item \textbf{Croquis (description pour réalisation) :}
    \begin{itemize}
        \item Contour du pays : triangle pointant au Nord (Lac Tchad), base au Sud (Golfe de Guinée).
        \item Villes : Douala (côte, estuaire du Wouri), Yaoundé (Centre-Sud), Bafoussam (Ouest, hautes terres), Bamenda (Nord-Ouest), Garoua (Nord, sur la Bénoué).
        \item Axe : flèche continue Douala $\rightarrow$ Yaoundé (N3) $\rightarrow$ Bafoussam (N4) $\rightarrow$ Bamenda (N6) $\rightarrow$ Garoua (N1).
        \item Légende : points proportionnels (Douala/Yaoundé gros, autres moyens), flèche du Nord, échelle graphique indicative.
    \end{itemize}
\end{enumerate}
\end{exoresolu}

\begin{methode}[Analyser les problèmes urbains et proposer des réponses adaptées]
\textbf{Objectif :} Lier les dysfonctionnements (habitat, transport, déchets, emploi, eau/énergie) aux causes structurelles et aux politiques publiques.
\textbf{Étapes :}
\begin{enumerate}
    \item \textbf{Lister les problèmes majeurs (mots-clés Bac) :}
    \begin{itemize}
        \item \textbf{Habitat :} habitat spontané et précaire (quartiers sous-équipés : Nylon à Douala, Mvog-Ada à Yaoundé), insalubrité, surpeuplement, insécurité foncière.
        \item \textbf{Mobilité :} embouteillages chroniques, insuffisance du transport en commun (taxis collectifs et motos-taxis \og bendskins \fg{} dominants), voirie dégradée, absence de transport de masse.
        \textbf{Services et réseaux :} coupures d'eau (Camwater) et d'électricité (ENEO) fréquentes, raccordements illégaux, pertes techniques et commerciales.
        \item \textbf{Gestion des déchets :} dépôts sauvages, décharges à ciel ouvert, faible tri et recyclage, inondations par bouchage des caniveaux.
        \item \textbf{Emploi et pauvreté :} secteur informel dominant (plus de 80\% des emplois urbains), sous-emploi, chômage des diplômés, pauvreté urbaine croissante.
    \end{itemize}
    \item \textbf{Identifier les causes structurelles :} urbanisation rapide non planifiée (la vitesse dépasse la capacité de gestion), faiblesse des ressources financières des communes, gouvernance foncière complexe (droit coutumier et droit moderne), absence de plans d'urbanisme actualisés et appliqués.
    \item \textbf{Citer les réponses et politiques (vocabulaire technique) :}
    \begin{itemize}
        \item \textbf{Résorption de l'habitat précaire :} restructuration des quartiers spontanés, régularisation foncière, ouverture de voiries primaires.
        \item \textbf{Mobilité :} réhabilitation de la voirie, projets de transport en commun (bus à haut niveau de service à Douala et Yaoundé), parkings relais.
        \item \textbf{Gestion des déchets :} partenariat public-privé (HYSACAM), centres de tri et de valorisation (compost, biogaz), sensibilisation.
        \item \textbf{Gouvernance :} décentralisation (transferts de compétences et de ressources aux communes), numérisation de la gestion foncière.
    \end{itemize}
    \item \textbf{Structurer la réponse :} Problème $\rightarrow$ Cause $\rightarrow$ Conséquence $\rightarrow$ Solution existante ou piste $\rightarrow$ Limites.
\end{enumerate}
\textbf{Réflexe Bac :} Ne pas se contenter de lister les problèmes. Le sujet demande souvent \og Quelles solutions ? \fg{} ou \og Évaluez l'efficacité \fg{}. Mobiliser les acteurs et outils nommés : communes urbaines, HYSACAM, plans d'urbanisme, partenariat public-privé.
\end{methode}

\begin{exoresolu}[Gestion des déchets solides à Douala : le cas HYSACAM]
\textbf{Contexte :} Douala produit environ 2 500 tonnes de déchets par jour. La société HYSACAM (Hygiène et Salubrité du Cameroun) assure la collecte en régie déléguée (partenariat public-privé). Les déchets sont enfouis dans la décharge technique de Nkolfoulou.
\textbf{Questions :}
\begin{enumerate}
    \item Calculer la production annuelle de déchets à Douala (en tonnes). Si la capacité de la décharge est de 1 500 000 tonnes, quelle déduction sur la durée de vie du site ?
    \item Citer 3 avantages et 3 limites du modèle de partenariat public-privé HYSACAM pour la propreté urbaine.
    \item Expliquer pourquoi, malgré HYSACAM, des dépôts sauvages persistent dans les quartiers périphériques (ex : Nylon, New-Bell).
    \item Proposer 2 mesures complémentaires (technique et citoyenne) pour tendre vers l'économie circulaire.
\end{enumerate}
\tcblower
\textbf{Solution détaillée :}
\begin{enumerate}
    \item \textbf{Calculs :}
    Production journalière = 2 500 t/j.
    Production annuelle = $2\,500 \times 365 = \textbf{912 500 tonnes/an}$.
    Capacité de la décharge = 1 500 000 tonnes.
    Durée de vie théorique = $\frac{1\,500\,000}{912\,500} \approx \textbf{1,6 an}$ (soit environ 1 an et 7 mois).
    \textbf{Déduction :} la capacité est \textbf{insuffisante} à court terme. Il faut prévoir une extension ou un nouveau site, et surtout réduire le volume enfoui (tri, valorisation).
    \item \textbf{Modèle HYSACAM :}
    \textbf{Avantages :} (1) professionnalisation et équipements (bennes, compacteurs) ; (2) régularité de la collecte dans les centres-villes et sur les axes principaux ; (3) création d'emplois (agents de propreté, chauffeurs, cadres).
    \textbf{Limites :} (1) coût élevé pour la commune (redevances) et impayés fréquents ; (2) couverture incomplète : quartiers spontanés et périphériques mal desservis (rues étroites, non-paiement des redevances par les usagers) ; (3) valorisation faible : l'essentiel est enfoui, peu de tri à la source ni d'unité industrielle de compostage à grande échelle.
    \item \textbf{Persistance des dépôts sauvages (quartiers périphériques) :}
    \begin{itemize}
        \item \textbf{Accessibilité :} voirie inexistante ou impraticable pour les bennes (rues étroites, non loties).
        \item \textbf{Foncier :} zones non loties, conflits fonciers, flou de responsabilité entre commune et opérateur.
        \item \textbf{Comportements :} absence de culture du tri, habitudes de dépôt dans les ravins et cours d'eau, refus de payer la redevance de propreté.
    \end{itemize}
    \item \textbf{Mesures pour l'économie circulaire :}
    \begin{enumerate}
        \item \textbf{Technique :} implanter des \textbf{centres de tri et de transfert de proximité} (dans chaque arrondissement) et des unités de \textbf{valorisation organique} (compostage pour l'agriculture urbaine, méthanisation pour le biogaz) $\rightarrow$ réduction du volume enfoui, création d'emplois verts.
        \item \textbf{Citoyenne :} \textbf{tri à la source} (bacs distincts : organique, recyclable, résiduel), tarification incitative (redevance indexée sur le volume non trié) et éducation environnementale dans les écoles et les quartiers.
    \end{enumerate}
\end{enumerate}
\end{exoresolu}

\begin{methode}[Analyser les relations ville-campagne (Dossier 2) : flux, interdépendance, recomposition]
\textbf{Objectif :} Dépasser l'opposition ville/campagne pour montrer les flux multidirectionnels et la recomposition des territoires ruraux (périurbanisation, rurbanisation).
\textbf{Étapes :}
\begin{enumerate}
    \item \textbf{Identifier les flux :}
    \begin{center}
    \begin{tabularx}{\linewidth}{|l|X|X|}\hline
    \textbf{Flux} & \textbf{Ville $\rightarrow$ Campagne} & \textbf{Campagne $\rightarrow$ Ville} \\ \hline
    \textbf{Humains} & Retraites, week-ends, investisseurs fonciers & Exode rural (jeunes, actifs), étudiants, patients \\ \hline
    \textbf{Économiques} & Intrants (engrais, semences), équipements, crédit, services & Produits vivriers, bois, artisanat, main-d'œuvre saisonnière \\ \hline
    \textbf{Financiers} & Transferts d'argent (mandats, mobile money), investissements & Épargne rurale (tontines), remboursements de crédits \\ \hline
    \textbf{Information} & Techniques nouvelles, formation, médias, administration & Savoirs locaux, produits de terroir \\ \hline
    \end{tabularx}
    \end{center}
    \item \textbf{Analyser l'interdépendance :} la ville \textbf{a besoin} de la campagne (nourriture, eau, bois, main-d'œuvre, espaces de loisir). La campagne \textbf{a besoin} de la ville (débouchés commerciaux, services de santé et d'éducation, équipements, monnaie, innovation).
    \item \textbf{Caractériser la recomposition des espaces ruraux :}
    \begin{itemize}
        \item \textbf{Périurbanisation :} étalement urbain dans le rural proche (0 à 30 km) : lotissements, habitat pavillonnaire, spéculation foncière, perte de terres agricoles de qualité (périphéries de Yaoundé, Douala, Bafoussam).
        \item \textbf{Rurbanisation :} villages adoptant des modes de vie urbains (maisons en parpaing, électricité, télévision, motos, commerces) sans être intégrés administrativement à la ville : les \og villages-villes \fg{}.
        \item \textbf{Agriculture urbaine et périurbaine :} maraîchage dans les bas-fonds (vallée du Mfoundi à Yaoundé ; bas-fonds de Douala), élevage porcin et avicole en ville. Réponse à l'insécurité alimentaire et revenu complémentaire, mais risques sanitaires (eaux usées, pesticides).
    \end{itemize}
    \item \textbf{Évaluer les politiques :} programmes de développement rural, désenclavement (routes rurales), électrification rurale, décentralisation, projets d'agropoles. Limites : financement, gouvernance, insécurité dans certaines régions.
\end{enumerate}
\textbf{Mots-clés Bac :} \cle{Interdépendance}, \cle{Périurbanisation}, \cle{Rurbanisation}, \cle{Flux alimentaires}, \cle{Exode rural}, \cle{Agropoles}.
\end{methode}

\begin{exoresolu}[Flux alimentaires et sécurité alimentaire : le bassin de l'Ouest vers Douala et Yaoundé]
\textbf{Document :} La région de l'Ouest (Bafoussam, Dschang, Bangangté, Foumban) est le \og grenier \fg{} du Cameroun. Productions : pomme de terre, carotte, poireau, chou, tomate, haricot, maïs, banane plantain. Axes routiers : Bafoussam-Yaoundé (N4, environ 280 km) et Bafoussam-Douala (environ 230 km).
\textbf{Questions :}
\begin{enumerate}
    \item Compléter le tableau des flux pour la filière \og légumes frais de l'Ouest vers Douala \fg{}.
    \item Calculer la part du prix final qui revient au producteur si : prix producteur (au champ) = 150 FCFA/kg ; prix de gros (marché de Bafoussam) = 250 FCFA/kg ; prix de détail (marché de Douala) = 600 FCFA/kg.
    \item Expliquer les 3 principaux goulets d'étranglement (contraintes) sur cet axe logistique.
    \item En quoi un projet d'agropole dans l'Ouest peut-il améliorer la relation ville-campagne pour cette filière ?
\end{enumerate}
\tcblower
\textbf{Solution détaillée :}
\begin{enumerate}
    \item \textbf{Tableau des flux de la filière légumes (Ouest $\rightarrow$ Douala) :}
    \begin{center}
    \begin{tabularx}{\linewidth}{|l|X|X|}\hline
    \textbf{Type de flux} & \textbf{Campagne (Ouest) $\rightarrow$ Ville (Douala)} & \textbf{Ville (Douala) $\rightarrow$ Campagne (Ouest)} \\ \hline
    \textbf{Marchandises} & Légumes frais (pomme de terre, carotte, chou...), bananes, tubercules & Intrants : engrais, semences certifiées, pesticides, emballages, carburant \\ \hline
    \textbf{Humains} & Producteurs, transporteurs, commissionnaires, revendeuses (\og bayam-sellam \fg{}) & Techniciens agricoles, acheteurs grossistes, investisseurs fonciers \\ \hline
    \textbf{Financiers} & Recettes des ventes (liquide, mobile money), taxes de marché & Paiement des intrants, salaires journaliers, crédits (microfinance), transferts \\ \hline
    \textbf{Information} & Demandes du marché (calibres, variétés), prix du jour & Prix des intrants, prévisions météo, techniques améliorées, formation \\ \hline
    \end{tabularx}
    \end{center}
    \item \textbf{Part de la valeur revenant au producteur :}
    \[Part = \frac{\text{Prix producteur}}{\text{Prix consommateur final}} \times 100 = \frac{150}{600} \times 100 = \textbf{25\%}.\]
    \textbf{Analyse :} le producteur ne touche qu'un quart de la valeur finale. Les 75\% restants rémunèrent le transport (route dégradée, péages), les intermédiaires multiples (collecteur, grossiste, demi-grossiste, détaillant), les pertes liées à la périssabilité et les taxes de marché.
    \item \textbf{Goulets d'étranglement (logistique) :}
    \begin{enumerate}
        \item \textbf{État des routes :} tronçons dégradés, nids-de-poule, glissements de terrain en saison des pluies $\rightarrow$ temps de trajet allongé, coût de transport élevé, pertes de légumes (chaleur, secousses).
        \item \textbf{Organisation de la filière :} atomisation des producteurs (petites exploitations de moins de 1 ha), faiblesse des organisations professionnelles pour négocier les prix, multiplicité des intermédiaires qui captent la marge.
        \item \textbf{Conservation et transformation :} quasi-absence de chaîne du froid (camions frigorifiques, chambres froides) et d'unités de transformation locale $\rightarrow$ obligation de vendre vite et à bas prix à la récolte, gaspillage post-récolte de 30 à 40\%.
    \end{enumerate}
    \item \textbf{Apport d'une agropole dans l'Ouest :}
    \begin{itemize}
        \item \textbf{Infrastructures :} plateforme logistique (chambres froides, quais de chargement), voirie d'accès, énergie stable, eau.
        \item \textbf{Organisation :} regroupement des producteurs en coopératives fortes, contrats avec grossistes et industriels (traçabilité, prix planchers).
        \item \textbf{Transformation :} unités intégrées (lavage, calibrage, emballage, transformation) $\rightarrow$ valorisation sur place, réduction des pertes, création d'emplois ruraux.
        \item \textbf{Financement :} guichet de crédit agricole au sein de l'agropole.
    \end{itemize}
    \textbf{Conclusion :} l'agropole fait passer la relation d'un rapport \og vendeur faible / acheteur fort \fg{} à un \textbf{partenariat structuré}, sécurisant le revenu paysan et l'approvisionnement des villes.
\end{enumerate}
\end{exoresolu}

\begin{methode}[Caractériser les systèmes agraires et l'agriculture camerounaise (dualité, défis, politiques)]
\textbf{Objectif :} Distinguer agriculture d'exportation (plantations) et agriculture vivrière (paysanne), analyser leurs déterminants, leurs performances et les politiques agricoles.
\textbf{Étapes :}
\begin{enumerate}
    \item \textbf{Poser la dualité structurelle (tableau comparatif) :}
    \begin{center}
    \begin{tabularx}{\linewidth}{|l|X|X|}\hline
    \textbf{Critère} & \textbf{Agriculture d'exportation (plantations)} & \textbf{Agriculture vivrière (paysanne)} \\ \hline
    \textbf{Spéculations} & Cacao, café, coton, banane, hévéa, palmier à huile, ananas & Maïs, manioc, plantain, igname, taro, riz, mil/sorgho, arachide, niébé, légumes \\ \hline
    \textbf{Acteurs} & Sociétés agro-industrielles (CDC, PHP, SOCAPALM, HEVECAM, SODECOTON), grands planteurs & Exploitations familiales (environ 90\% des exploitations, 1 à 3 ha), main-d'œuvre familiale \\ \hline
    \textbf{Techniques} & Mécanisation partielle, intrants chimiques, variétés sélectionnées, encadrement technique & Itinéraires traditionnels, faible mécanisation (houe, machette), intrants rares et chers, semences paysannes, jachère raccourcie \\ \hline
    \textbf{Foncier} & Concessions de longue durée, grands espaces continus & Droit coutumier dominant, morcellement, insécurité foncière \\ \hline
    \textbf{Commercialisation} & Filières organisées, contrats, exportation par le port de Douala & Circuits informels, multiplicité des intermédiaires, marchés locaux, fortes pertes post-récolte \\ \hline
    \textbf{Revenus} & Salaires, devises pour l'État & Autoconsommation (60 à 80\% de la production), vente de l'excédent = principal revenu monétaire rural \\ \hline
    \end{tabularx}
    \end{center}
    \item \textbf{Localiser les bassins de production (croquis mental) :}
    \begin{itemize}
        \item \textbf{Zone forestière humide (Sud, Centre, Est, Littoral, Sud-Ouest) :} cacao, café robusta, palmier à huile, hévéa, banane. Vivriers : plantain, manioc, macabo, taro.
        \item \textbf{Hautes terres (Ouest, Nord-Ouest) :} café arabica, légumes de climat tempéré (pomme de terre, carotte), haricot, maïs, élevage bovin.
        \textbf{Zone soudanienne (Nord, Adamaoua) :} coton, mil, sorgho, maïs, arachide, niébé, oignon, élevage bovin extensif.
        \item \textbf{Zone sahélienne (Extrême-Nord) :} mil, sorgho, riz (plaines inondables du Logone et du Mayo-Danay), coton, arachide, élevage transhumant.
    \end{itemize}
    \item \textbf{Analyser les contraintes transverses :} changement climatique (sécheresses au Nord, irrégularité des pluies), dégradation des sols (érosion, perte de fertilité), accès difficile aux intrants et au financement, vieillissement du paysannat, faible transformation locale (exportation de produits bruts), gouvernance des filières (prix payés aux producteurs faibles et instables).
    \item \textbf{Citer les leviers politiques :} mécanisation, semences certifiées (IRAD), subvention des intrants, irrigation, transformation locale, sécurisation foncière rurale, promotion des jeunes agripreneurs et des agropoles (stratégie nationale de développement SND30).
\end{enumerate}
\textbf{Chiffres clés Bac :} l'agriculture représente environ 15 à 20\% du PIB, emploie environ 60\% de la population active et fournit une part importante des recettes d'exportation (cacao, coton, banane, caoutchouc, café). Autosuffisance alimentaire relative, avec des déficits en riz, blé, poisson, lait et huile.
\end{methode}

\begin{exoresolu}[La filière coton au Nord-Cameroun : enjeux et limites du modèle SODECOTON]
\textbf{Contexte :} Le coton est la principale culture de rente du Septentrion. Il est encadré par la SODECOTON (Société de Développement du Coton), qui assure l'encadrement, la collecte, l'égrenage et la commercialisation. Zone : Nord, Adamaoua, Extrême-Nord. Producteurs : environ 250 000 exploitations familiales.
\textbf{Données :} rendement moyen paysan : 1 200 kg/ha (coton graine). Prix producteur : 245 FCFA/kg. Coût des intrants (engrais, pesticides, semences) fournis à crédit par la SODECOTON : environ 120 000 FCFA/ha.
\textbf{Questions :}
\begin{enumerate}
    \item Calculer le revenu brut (RB) et le revenu net (RN) par hectare pour un producteur cotonnier moyen. Exprimer le RN en dollars US (1 \$ = 600 FCFA).
    \item Expliquer pourquoi on parle de \og modèle intégré \fg{} pour la SODECOTON. Citer 3 services rendus aux producteurs.
    \item Identifier 3 limites structurelles de ce modèle aujourd'hui (économiques, écologiques, sociales).
    \item Proposer 2 pistes de diversification ou d'intensification durable pour sécuriser le revenu paysan.
\end{enumerate}
\tcblower
\textbf{Solution détaillée :}
\begin{enumerate}
    \item \textbf{Calculs économiques (par hectare) :}
    \begin{align*}
        \text{Revenu Brut (RB)} &= \text{Rendement} \times \text{Prix producteur} \\
        &= 1\,200 \text{ kg} \times 245 \text{ FCFA/kg} = \textbf{294 000 FCFA/ha} \\
        \text{Revenu Net (RN)} &= RB - \text{Coût des intrants} = 294\,000 - 120\,000 = \textbf{174 000 FCFA/ha} \\
        \text{RN en \$US} &= \frac{174\,000}{600} = \textbf{290 \$US/ha}
    \end{align*}
    \textbf{Analyse :} revenu net très faible (290 \$/ha/an). Pour une exploitation de 3 ha de coton : environ 870 \$/an, soit moins de 2,5 \$/jour $\rightarrow$ niveau de \textbf{pauvreté}. Dépendance totale au cours mondial du coton et au taux de change.
    \item \textbf{Modèle intégré SODECOTON :} l'entreprise encadre \textbf{toute la chaîne} : recherche variétale $\rightarrow$ distribution des intrants à crédit $\rightarrow$ encadrement technique $\rightarrow$ collecte et transport $\rightarrow$ égrenage (usines de Garoua, Maroua, Kaélé) $\rightarrow$ exportation de la fibre par le port de Douala $\rightarrow$ paiement des producteurs à prix garanti.
    \textbf{3 services :} (1) crédit intrants sans garantie foncière (la récolte future sert de garantie) ; (2) encadrement technique de proximité ; (3) prix plancher garanti qui amortit les chutes du cours mondial.
    \item \textbf{Limites structurelles :}
    \begin{enumerate}
        \item \textbf{Économique :} dépendance à une monoculture (risque lié au cours mondial), coût croissant des intrants importés, rentabilité déclinante face aux cultures vivrières (maïs, niébé, sésame).
        \item \textbf{Écologique :} épuisement des sols (coton sur coton, disparition de la jachère), pollution chimique (pesticides : santé des paysans, eau, biodiversité), déforestation pour ouvrir de nouveaux champs.
        \item \textbf{Sociale :} vieillissement des producteurs (les jeunes partent vers les villes ou l'orpaillage), endettement chronique, conflits fonciers entre éleveurs transhumants et cultivateurs.
    \end{enumerate}
    \item \textbf{Pistes de diversification et d'intensification durable :}
    \begin{enumerate}
        \item \textbf{Rotation coton-légumineuses (niébé, arachide, soja) :} fixation de l'azote atmosphérique $\rightarrow$ réduction des engrais azotés, rupture du cycle des ravageurs, revenu alimentaire et marchand complémentaire.
        \item \textbf{Agroforesterie :} intégration d'arbres fertilisants (Faidherbia albida) ou à valeur économique (karité, moringa) $\rightarrow$ fertilité des sols, bois-énergie, revenus diversifiés, résilience climatique.
    \end{enumerate}
\end{enumerate}
\end{exoresolu}

\begin{methode}[Analyser les systèmes d'élevage et leurs interactions avec l'agriculture et l'espace]
\textbf{Objectif :} Distinguer élevage moderne et élevage traditionnel/pastoral ; analyser les conflits d'usage (transhumance, parcours, eau) et les politiques d'élevage.
\textbf{Étapes :}
\begin{enumerate}
    \item \textbf{Typologier l'élevage (tableau) :}
    \begin{center}
    \begin{tabularx}{\linewidth}{|l|X|X|}\hline
    \textbf{Critère} & \textbf{Élevage traditionnel / pastoral (Nord, Adamaoua, Est)} & \textbf{Élevage moderne / périurbain (grandes villes, Ouest, Littoral)} \\ \hline
    \textbf{Espèces} & Bovins (zébus Goudali, White Fulani), ovins et caprins, ânes, camélins (Extrême-Nord) & Bovins laitiers (races améliorées et croisements), porcs, volailles (poulets de chair, pondeuses) \\ \hline
    \textbf{Système} & Extensif, parcours libres, transhumance saisonnière, exploitation familiale & Intensif (stabulation), aliments composés, bâtiments et équipements, main-d'œuvre qualifiée \\ \hline
    \textbf{Alimentation} & Pâturage naturel, résidus de culture, points d'eau naturels et forages & Aliments concentrés achetés (60 à 70\% des charges), fourrages cultivés, compléments minéraux \\ \hline
    \textbf{Sanitaire} & Campagnes de vaccination publiques, traitements traditionnels, mortalité élevée des jeunes & Prophylaxie rigoureuse (vaccins, biosécurité), suivi vétérinaire privé \\ \hline
    \textbf{Commercialisation} & Marchés à bétail hebdomadaires (Garoua, Maroua, Ngaoundéré), exportation régionale (Nigeria, Gabon, Congo) & Circuits courts (abattoirs modernes, supermarchés, hôtels), contrats \\ \hline
    \textbf{Productivité} & Faible : faible taux d'exploitation du troupeau, premier vêlage tardif, lait 1 à 2 L/j/vache & Élevée : cycles courts (porc, volaille), lait 15 à 25 L/j/vache \\ \hline
    \end{tabularx}
    \end{center}
    \item \textbf{Analyser les conflits pour les ressources et l'espace (cœur du dossier) :}
    \begin{itemize}
        \item \textbf{Transhumance :} descente des troupeaux vers le Sud (Adamaoua, Est, Nord-Ouest) en saison sèche (novembre-mai) à la recherche d'eau et de pâturage ; remontée vers le Nord en saison des pluies (juin-octobre).
        \item \textbf{Conflits agriculteurs-éleveurs :} divagation des troupeaux dans les champs, destruction des cultures, blocage des couloirs de transhumance par l'extension des cultures, pollution des points d'eau, vols de bétail.
        \item \textbf{Causes profondes :} pression démographique et foncière, changement climatique (sécheresses au Nord $\rightarrow$ transhumance plus précoce et plus longue), insécurité, gouvernance foncière défaillante (couloirs non matérialisés, pâturages non délimités), absence d'accords locaux de transhumance.
    \end{itemize}
    \item \textbf{Citer les politiques et réponses :} projets d'appui à l'élevage et au pastoralisme. Actions : aménagement des couloirs, pâturages et points d'eau (bornage, mares, forages), sécurisation de la transhumance (accords locaux, brigades mixtes), amélioration génétique, cultures fourragères et banques de foin, marchés à bétail modernes et abattoirs, santé animale (vaccins, laboratoire LANAVET).
\end{enumerate}
\textbf{Chiffres clés :} cheptel bovin d'environ 6 à 7 millions de têtes (environ 90\% en élevage traditionnel). La production laitière locale ne couvre qu'une partie des besoins $\rightarrow$ importations massives de poudre de lait.
\end{methode}

\begin{exoresolu}[Conflits agriculteurs-éleveurs dans la région de l'Adamaoua : analyse et médiation]
\textbf{Situation :} La région de l'Adamaoua (hautes terres, pâturages naturels, zone de transit de la transhumance) connaît une intensification des conflits. Les agriculteurs étendent les cultures (maïs, igname, soja) dans les bas-fonds et les couloirs de passage. Les éleveurs dénoncent le blocage des couloirs ; les agriculteurs dénoncent la divagation des troupeaux et la destruction des cultures.
\textbf{Questions :}
\begin{enumerate}
    \item Construire un tableau d'analyse des causes (facteurs structurels et facteurs conjoncturels) de ce conflit.
    \item Expliquer le mécanisme de la \og transhumance contrainte \fg{} observée ces dernières années.
    \item En tant que géographe consultant pour le Préfet, proposer une démarche en 4 étapes pour la sécurisation foncière des couloirs de transhumance.
    \item Calculer la capacité de charge (UGB/ha) d'un pâturage naturel de l'Adamaoua si : production fourragère moyenne = 3 000 kg MS/ha/an ; besoin d'une unité de gros bétail (UGB) = 3 000 kg MS/an ; taux d'utilisation préconisé = 50\% (préservation du milieu). Combien d'UGB pour 1 000 ha ? Est-ce suffisant pour un troupeau transhumant de 500 UGB séjournant 6 mois ?
\end{enumerate}
\tcblower
\textbf{Solution détaillée :}
\begin{enumerate}
    \item \textbf{Tableau d'analyse des causes :}
    \begin{center}
    \begin{tabularx}{\linewidth}{|l|X|X|}\hline
    \textbf{Dimension} & \textbf{Facteurs structurels (long terme)} & \textbf{Facteurs conjoncturels (déclencheurs)} \\ \hline
    \textbf{Foncier} & Couloirs et pâturages non classés, morcellement coutumier, ventes de terres à des migrants agriculteurs, absence de titres fonciers pour les éleveurs & Spéculation foncière autour de Ngaoundéré, attribution de titres sur d'anciens parcours \\ \hline
    \textbf{Démographie et climat} & Croissance de la population rurale, extension des cultures vivrières et du coton, sécheresses récurrentes au Nord & Année sèche exceptionnelle, arrivée massive de réfugiés et de déplacés avec leurs troupeaux \\ \hline
    \textbf{Gouvernance} & Faiblesse de l'administration locale, justice lente, absence de commissions de concertation fonctionnelles & Impunité des auteurs de violences, non-paiement des dommages \\ \hline
    \textbf{Socio-économique} & Paupérisation des éleveurs, perte des savoirs traditionnels de médiation, montée des tensions identitaires & Vols de bétail, destruction de récoltes à maturité, blocus des marchés \\ \hline
    \end{tabularx}
    \end{center}
    \item \textbf{Transhumance contrainte :}
    Habituellement : descente Nord $\rightarrow$ Sud en saison sèche, remontée Sud $\rightarrow$ Nord en saison des pluies.
    \textbf{Actuellement :} sécheresse au Nord + insécurité + extension des cultures septentrionales (coton) $\rightarrow$ les éleveurs \textbf{ne remontent plus au Nord} (ou plus tard, ou par petits groupes). Ils \textbf{restent bloqués au Sud} (Adamaoua, Est, Nord-Ouest) toute l'année. Conséquences : \textbf{surpâturage} des parcours méridionaux, destruction de la végétation ligneuse, conflits exacerbés avec les agriculteurs du Sud qui n'attendent pas les troupeaux en saison des pluies. C'est une \textbf{rupture du cycle traditionnel}.
    \item \textbf{Démarche de sécurisation des couloirs (4 étapes) :}
    \begin{enumerate}
        \item \textbf{Diagnostic participatif et cartographie :} inventaire des couloirs existants (tracé GPS, largeur), points d'eau, pâturages, zones cultivées, villages ; concertation large (chefs traditionnels, mairies, services techniques de l'agriculture et de l'élevage, associations d'éleveurs et d'agriculteurs) ; carte unique partagée.
        \item \textbf{Délimitation juridique et matérialisation :} arrêté préfectoral de classement des couloirs de transhumance et des pâturages communaux ; pose de bornes, plantation de haies vives, aménagement de points d'eau (mares, forages solaires) dans le couloir.
        \item \textbf{Contrats locaux de transhumance :} signature d'accords entre mairies, groupements d'éleveurs et groupements d'agriculteurs : calendrier de passage à dates fixes, largeur respectée, indemnisation rapide et forfaitaire des dégâts, sanctions conventionnelles.
        \item \textbf{Suivi-évaluation et pérennisation :} brigade mixte de surveillance, comité de pilotage trimestriel, intégration dans les plans communaux de développement, formation des bergers et des agriculteurs.
    \end{enumerate}
    \item \textbf{Calcul de la capacité de charge (CC) :}
    \begin{align*}
        \text{Production utile (PU)} &= \text{Production totale} \times \text{Taux d'utilisation} \\
        &= 3\,000 \text{ kg MS/ha/an} \times 0,50 = \textbf{1 500 kg MS/ha/an} \\
        \text{Capacité de charge (CC)} &= \frac{\text{PU}}{\text{Besoin UGB/an}} = \frac{1\,500}{3\,000} = \textbf{0,5 UGB/ha/an} \\
        \text{Capacité totale pour 1 000 ha} &= 0,5 \times 1\,000 = \textbf{500 UGB/an} \\
        \text{Besoin du troupeau (500 UGB, 6 mois)} &= 500 \times \frac{6}{12} = \textbf{250 UGB-années} \\
        \text{Comparaison} &: 500 > 250 \rightarrow \textbf{théoriquement suffisant sur l'année}.
    \end{align*}
    \textbf{Nuance cruciale :} la production fourragère est \textbf{saisonnière} (pic de juillet à septembre, quasi nulle de janvier à avril). Si les 500 UGB sont présents pendant la saison sèche (production $\approx$ 0), il y a \textbf{surpâturage massif}, sauf s'il existe des stocks de foin ou une alimentation d'appoint. Le calcul annuel masque le \textbf{déficit saisonnier}. \textbf{Conclusion :} le pâturage est sous-dimensionné pour accueillir le troupeau en saison sèche sans complémentation.
\end{enumerate}
\end{exoresolu}

\begin{methode}[Synthèse : villes, campagnes et développement – rédiger une analyse structurée]
\textbf{Objectif :} Mobiliser l'ensemble des connaissances de l'unité (urbanisation, problèmes urbains, relations ville-campagne, agriculture, élevage) pour répondre à un sujet de type \og Discutez \fg{}, \og Analysez \fg{}, \og Dans quelle mesure \fg{}.
\textbf{Structure type (plan dialectique en 3 parties) :}
\begin{enumerate}
    \item \textbf{Introduction :} accroche (chiffre clé : TU supérieur à 50\%, environ 60\% d'actifs agricoles), définition des termes clés (ville, campagne, développement, aménagement du territoire), analyse du sujet, annonce du plan.
    \item \textbf{Partie 1 : Des territoires en mutation rapide sous l'effet de l'urbanisation (dynamiques).}
        \begin{itemize}
            \item A. Explosion urbaine (métropolisation de Douala et Yaoundé, villes moyennes) et étalement (périurbanisation).
            \item B. Récomposition des campagnes (rurbanisation, perte de terres, diversification des activités non agricoles).
            \item C. Flux intenses (humains, économiques, informationnels) créant une forte \textbf{interdépendance fonctionnelle}.
        \end{itemize}
    \item \textbf{Partie 2 : Des dysfonctionnements majeurs freinant le développement intégré (tensions).}
        \begin{itemize}
            \item A. Ville : pauvreté urbaine, habitat précaire, congestion, dépendance alimentaire aux importations, gestion défaillante des déchets, de l'eau et de l'énergie.
            \item B. Campagne : faible productivité agropastorale, enclavement, vieillissement, conflits fonciers et de transhumance, faible valorisation locale.
            \item C. Relation ville-campagne asymétrique : la ville prélève les ressources (vivriers, main-d'œuvre, eau, bois) sans réinvestir suffisamment (prix bas aux producteurs, infrastructures rurales dégradées).
        \end{itemize}
    \item \textbf{Partie 3 : Vers un développement territorial intégré ? Politiques et leviers (perspectives).}
        \begin{itemize}
            \item A. Aménagement du territoire : polycentrisme (agropoles, pôles régionaux, décentralisation effective), désenclavement rural (routes, numérique, énergie).
            \item B. Transformation structurelle : agro-industrialisation locale (valeur ajoutée, emploi des jeunes), intensification écologique (agroécologie, irrigation, semences et races améliorées), économie circulaire (déchets urbains $\rightarrow$ engrais pour les campagnes).
            \item C. Gouvernance foncière et sociale : sécurisation foncière rurale et urbaine, contrats locaux ville-campagne (transhumance, marchés), protection des espaces agricoles périurbains.
        \end{itemize}
    \item \textbf{Conclusion :} bilan nuancé (interdépendance réelle mais inégalitaire), ouverture (défi démographique à l'horizon 2050, urgence climatique).
\end{enumerate}
\textbf{Conseils de rédaction Bac :} utiliser des \textbf{exemples précis} (noms de villes, régions, projets : SODECOTON, HYSACAM, agropoles, SND30). Chiffrer (\%, FCFA, ha, km). Connecteurs logiques : \emph{certes... mais... ; par ailleurs ; en conséquence ; cette situation s'explique par...}
\end{methode}

\begin{exoresolu}[Sujet type Bac : \og L'urbanisation rapide au Cameroun : chance ou menace pour le développement rural ? \fg{}]
\textbf{Consigne :} En vous appuyant sur des exemples précis pris au Cameroun, montrez que l'urbanisation rapide constitue à la fois un moteur et un défi pour le développement des campagnes, et proposez des pistes pour une articulation vertueuse ville-campagne.
\tcblower
\textbf{Introduction}

[Accroche] Le Cameroun compte aujourd'hui plus d'un urbain sur deux (52,1\% au RGPH 2005, environ 56\% selon les estimations récentes), contre moins de 15\% à l'indépendance. Douala et Yaoundé forment un bipôle d'environ 7,5 millions d'habitants.
[Définitions] Urbanisation = croissance de la part de la population urbaine. Développement rural = amélioration durable des conditions de vie, des revenus et des services en zone rurale.
[Analyse du sujet] L'urbanisation bouleverse les rapports ville-campagne. Chance par les flux (marchés, argent, innovation) ? Menace par les prélèvements (foncier, eau, main-d'œuvre) et la pollution ?
[Plan] 1. L'urbanisation, moteur de dynamiques rurales. 2. L'urbanisation, source de pressions et de vulnérabilités. 3. Construire une articulation vertueuse : le pari de l'aménagement intégré.

\textbf{Développement}

\textbf{1. L'urbanisation, moteur de dynamiques rurales (la chance)}
\begin{itemize}
    \item \textbf{Débouchés massifs pour l'agropastoral :} des millions de citadins $\rightarrow$ demande alimentaire énorme (céréales, tubercules, légumes, viande, lait, poisson). Ex : bassin de l'Ouest (pomme de terre, carotte) vers Douala et Yaoundé ; bassin cotonnier du Nord vers les usines et l'exportation via Douala ; bassin laitier de l'Adamaoua et du Nord-Ouest vers les laiteries urbaines.
    \item \textbf{Flux financiers de retour (transferts) :} les migrants urbains envoient de l'argent (mobile money) pour les intrants, les équipements (motos, pompes), la construction, la santé et l'éducation. Ces transferts constituent un financement essentiel de l'agriculture familiale.
    \item \textbf{Innovation et services :} les villes concentrent la recherche (IRAD, universités), la formation, les services vétérinaires et phytosanitaires, les banques et la microfinance, les médias qui diffusent techniques et prix des marchés.
    \item \textbf{Diversification rurale non agricole :} tourisme rural (monts Mandara, chefferies bamiléké), artisanat vendu en ville, commerce et transport.
\end{itemize}

\textbf{2. L'urbanisation, source de pressions et de vulnérabilités (la menace)}
\begin{itemize}
    \item \textbf{Accaparement foncier et spéculation :} périurbanisation galopante (rayon de 30 à 50 km). Ex : Sa'a, Obala autour de Yaoundé ; périphérie de Douala. Des terres agricoles fertiles (bas-fonds, sols volcaniques) sont loties pour l'habitat. Flambée des prix fonciers $\rightarrow$ exclusion des paysans et des jeunes.
    \item \textbf{Prélèvement des ressources et pollution en retour :} prélèvements d'eau urbains $\rightarrow$ baisse des ressources en aval (irrigation, pêche, eau potable rurale). Rejets d'eaux usées non traitées et déchets solides $\rightarrow$ pollution des sols et des eaux des zones de production périurbaine (maraîchage des bas-fonds du Mfoundi : risques sanitaires).
    \item \textbf{Exode rural sélectif et vieillissement :} départ des jeunes valides $\rightarrow$ main-d'œuvre agricole rare et coûteuse, féminisation et vieillissement du paysannat, jachères raccourcies $\rightarrow$ dégradation des sols.
    \item \textbf{Concurrence des importations :} la ville consomme riz, blé, poisson et lait importés $\rightarrow$ dévalorisation des productions locales (riz du Nord et de l'Extrême-Nord, maïs, manioc).
\end{itemize}

\textbf{3. Vers une articulation vertueuse : le pari de l'aménagement intégré (pistes)}
\begin{itemize}
    \item \textbf{Polycentrisme et agropoles :} sortir de la macrocéphalie Douala-Yaoundé en développant des pôles régionaux (agropoles : coton et viande au Nord, légumes et lait dans l'Ouest, cacao et bois au Sud). Transformer \textbf{sur place} (création de valeur, d'emplois, fixation des populations). Ex : abattoirs modernes régionaux, unités de transformation du cacao et de la tomate.
    \item \textbf{Circuits courts et cantines scolaires :} relier directement les coopératives rurales aux cantines urbaines (écoles, hôpitaux, entreprises) via des plateformes logistiques régionales : prix rémunérateur pour le producteur, prix stable pour le consommateur, traçabilité.
    \item \textbf{Économie circulaire urbain-rural :} valorisation des déchets urbains (compost, biogaz) $\rightarrow$ engrais organiques pour l'agriculture périurbaine et rurale (substitution aux engrais chimiques importés) ; eaux usées traitées $\rightarrow$ irrigation maraîchère contrôlée.
    \item \textbf{Gouvernance foncière concertée :} classement préalable des espaces agricoles, pastoraux et forestiers \textbf{avant} l'extension urbaine ; création de ceintures vertes productives protégées autour des métropoles ; baux ruraux de longue durée pour les jeunes agripreneurs.
    \item \textbf{Investissements structurants ruraux :} routes rurales (désenclavement des bassins de production), électricité (solaire, hydroélectricité), numérique, eau (barrages multi-usages), santé et éducation de qualité pour fixer les jeunes.
\end{itemize}

\textbf{Conclusion}

L'urbanisation camerounaise n'est ni une chance pure ni une menace fatale : c'est un \textbf{fait structurel majeur} qui réorganise le territoire. Actuellement, la relation est \textbf{asymétrique} : la ville prélève la substance de la campagne (vivriers, main-d'œuvre, eau, foncier) sans lui rendre suffisamment de valeur ajoutée, de services et de propreté. Le développement durable du Cameroun (SND30) passe par l'\textbf{inversion de cette pente} : faire de la ville un \textbf{moteur de transformation rurale} (marchés stables, innovation, capitaux, recyclage) et de la campagne un \textbf{territoire de production performant, résilient et vivable}. Cela exige une \textbf{planification spatiale forte}, une \textbf{décentralisation financière réelle} et un \textbf{pacte ville-campagne} (prix justes, services partagés). Le défi démographique des prochaines décennies impose d'agir \textbf{maintenant} sur les couloirs de transhumance, les agropoles et la gouvernance foncière.
\end{exoresolu}

\begin{attention}[Les 5 erreurs qui coûtent des points au Bac – Unité \og Les villes \fg{}]
\begin{enumerate}
    \item \textbf{Confondre \og taux d'urbanisation \fg{} et \og vitesse d'urbanisation \fg{} / \og taux de croissance urbaine \fg{}.}
    \emph{Erreur :} \og Le taux d'urbanisation est de 4,7\% par an. \fg{}
    \emph{Correct :} le \textbf{taux d'urbanisation} est une \textbf{part} (52,1\% en 2005). La \textbf{vitesse d'urbanisation} (ou TCAM urbain) est un \textbf{taux de variation} (4,7\%/an). Ne jamais mélanger \% (niveau) et \%/an (rythme).
    \item \textbf{Traiter \og ville \fg{} et \og campagne \fg{} comme deux mondes étanches (opposition binaire) au lieu d'analyser les flux et l'interdépendance.}
    \emph{Erreur :} \og La ville est moderne, la campagne est traditionnelle. \fg{} / \og L'exode rural vide la campagne. \fg{}
    \emph{Correct :} montrer les \textbf{flux bidirectionnels} (argent, intrants, informations vers le rural ; vivriers, main-d'œuvre, eau, bois vers l'urbain). Parler de \textbf{recomposition} (périurbanisation, rurbanisation, agriculture urbaine) et non de simple vidage.
    \item \textbf{Oublier la dualité agricole (plantations vs paysanne) et la dualité de l'élevage (pastoral vs moderne) dans l'analyse des relations ville-campagne.}
    \emph{Erreur :} parler de \og l'agriculture \fg{} ou de \og l'élevage \fg{} au singulier sans distinguer logiques, acteurs, espaces et productivités.
    \emph{Correct :} systématiser le \textbf{tableau comparatif}. Citer les spéculations, localisations et acteurs précis (SODECOTON, CDC, HYSACAM, bayam-sellam). Chiffrer (rendements, parts du PIB, des exportations, de l'emploi).
    \item \textbf{Ne pas chiffrer et ne pas localiser (croquis mental absent).}
    \emph{Erreur :} \og Beaucoup de gens vont en ville. \fg{} \og Le Nord produit du coton. \fg{}
    \emph{Correct :} \og TU 52,1\% (RGPH 2005), TCAM urbain 4,7\%/an, Douala-Yaoundé environ 7,5 M hab. \fg{} \og Région du Nord : environ 250 000 exploitants cotonniers, 1 200 kg/ha, 245 FCFA/kg. \fg{} Toujours avoir 3 ou 4 chiffres clés et situer sur une carte mentale (régions, bassins, axes routiers).
    \item \textbf{Proposer des solutions \og magiques \fg{} ou hors sol (généralités) au lieu de politiques publiques nommées et évaluées.}
    \emph{Erreur :} \og Il faut moderniser l'agriculture. \fg{} \og Il faut construire des routes. \fg{} \og Il faut sensibiliser. \fg{}
    \emph{Correct :} citer les \textbf{programmes et outils réels} : SND30, agropoles, projets d'appui à l'élevage et au pastoralisme, partenariat HYSACAM, plans d'urbanisme, classement des couloirs de transhumance. En citer les \textbf{actions concrètes} (intrants subventionnés, couloirs bornés, compostage) et leurs \textbf{limites réelles} (financement, gouvernance, insécurité, changement climatique).
\end{enumerate}
\end{attention}

\newpage

\coursec{Exercices}

\courssub{Je vérifie mes connaissances}

\begin{exercice}[QCM : Urbanisation et problèmes urbains]
Parmi les propositions suivantes, une seule est exacte. Recopie la lettre correspondante et justifie brièvement ton choix.
\begin{enumerate}
\item Le taux d'urbanisation du Cameroun en 2020 est estimé à environ :
\begin{itemize}
\item[A.] 30 \%
\item[B.] 58 \%
\item[C.] 75 \%
\item[D.] 90 \%
\end{itemize}
\item Un quartier spontané se caractérise principalement par :
\begin{itemize}
\item[A.] Un tracé orthogonal et des équipements collectifs complets
\item[B.] Une occupation illégale du sol, un habitat précaire et l'absence de voirie
\item[C.] Une forte densité de population aisée et des villas de standing
\item[D.] Une planification rigoureuse par les services d'urbanisme
\end{itemize}
\item La croissance urbaine au Cameroun est essentiellement due à :
\begin{itemize}
\item[A.] L'excédent naturel seul
\item[B.] L'immigration internationale uniquement
\item[C.] L'exode rural et l'accroissement naturel
\item[D.] La création de nouvelles capitales régionales
\end{itemize}
\item La ville de Douala est le premier pôle urbain camerounais car elle est :
\begin{itemize}
\item[A.] La capitale politique
\item[B.] Le principal port et le poumon économique du pays
\item[C.] La ville la plus ancienne
\item[D.] La plus étendue en superficie
\end{itemize}
\end{enumerate}
\end{exercice}

\begin{exercice}[Vrai ou Faux justifié : Relations ville-campagne]
Indique si chaque affirmation est Vraie (V) ou Fausse (F) et justifie ta réponse en une phrase en t'appuyant sur les notions du cours.
\begin{enumerate}
\item Les flux de produits agricoles vont exclusivement de la campagne vers la ville.
\item Les transferts d'argent des migrants urbains vers leurs familles rurales constituent un flux financier ville $\to$ campagne.
\item L'exode rural vide les campagnes de leur main-d'œuvre active, ce qui favorise toujours la modernisation de l'agriculture.
\item Les villes camerounaises approvisionnent les campagnes en produits manufacturiers et en services (santé, éducation).
\end{enumerate}
\end{exercice}

\begin{exercice}[Définitions et distinctions : Agriculture et Élevage]
Donne une définition précise pour chacun des termes suivants et distingue les paires de concepts indiquées.
\begin{enumerate}
\item \textbf{Agriculture vivrière} / \textbf{Agriculture de rente (ou d'exportation)}.
\item \textbf{Élevage extensif} / \textbf{Élevage intensif}.
\item \textbf{Jachère} / \textbf{Rotation culturale}.
\item \textbf{Produit agricole de rente} : cite deux exemples majeurs pour le Cameroun.
\end{enumerate}
\end{exercice}

\begin{exercice}[Calculs directs : Taux d'urbanisation et productivité]
\begin{enumerate}
\item En 1976, le Cameroun comptait 10 millions d'habitants dont 2,5 millions d'urbains. En 2020, la population est de 26 millions dont 15 millions d'urbains. Calcule le taux d'urbanisation pour ces deux dates (arrondi au dixième) et le taux d'accroissement annuel moyen de la population urbaine sur la période (formule simplifiée : $t = \frac{1}{n} \times \ln\left(\frac{P_{final}}{P_{initial}}\right) \times 100$, avec $n$ le nombre d'années).
\item Une exploitation de maïs de 5 hectares a produit 12,5 tonnes de grains. Calcule le rendement à l'hectare (en t/ha). Si l'objectif de l'État est d'atteindre 4 t/ha, cet exploitant atteint-il l'objectif ?
\item Un éleveur possède 200 têtes de bétail sur 500 hectares de parcours. Calcule le chargement (en UBT/ha) en supposant 1 UBT = 1 tête. Ce chargement correspond-il à un élevage extensif ou intensif ? Justifie.
\end{enumerate}
\end{exercice}

\courssub{Je m'entraîne}

\begin{exercice}[Analyse de document : Quartier spontané à Yaoundé]
\textbf{Document :} \emph{Photographie aérienne d'un quartier de la périphérie de Yaoundé (Mvog-Ada). On y observe un tissu urbain dense, anarchique, des ruelles étroites impraticables aux véhicules, des maisons en banco ou en parpaings non enduits, l'absence d'évacuation des eaux usées (ruisseaux à ciel ouvert), des branchements électriques sauvages, une promiscuité extrême entre habitats et activités (petit commerce, ateliers).}
\begin{enumerate}
\item Décris le paysage urbain présenté sur le document (morphologie, habitat, voirie, réseaux).
\item Identifie et explique \textbf{deux problèmes majeurs} auxquels sont confrontés les habitants de ce quartier (sanitaire, sécurité, foncier, environnemental).
\item Propose \textbf{une solution concrète et justifiée} que les pouvoirs publics pourraient mettre en œuvre pour résorber l'un de ces problèmes (ex. : restructuration, lotissement, relogement, équipements).
\end{enumerate}
\end{exercice}

\begin{exercice}[Exploitation de données : Évolution de l'urbanisation au Cameroun (1960-2020)]
\textbf{Document :} \emph{Tableau ci-dessous.}
\begin{center}
\begin{tabular}{|c|c|c|c|c|c|c|c|}
\hline
\textbf{Année} & 1960 & 1976 & 1987 & 2005 & 2010 & 2015 & 2020 \\
\hline
\textbf{Population totale (millions)} & 5,4 & 7,6 & 10,5 & 17,5 & 19,4 & 22,8 & 26,0 \\
\hline
\textbf{Population urbaine (millions)} & 0,7 & 1,5 & 3,5 & 8,8 & 10,5 & 13,0 & 15,1 \\
\hline
\textbf{Taux d'urbanisation (\%)} & 13,0 & 19,7 & 33,3 & 50,3 & 54,1 & 57,0 & 58,1 \\
\hline
\end{tabular}
\end{center}
\begin{enumerate}
\item Vérifie par le calcul les taux d'urbanisation indiqués pour 1960 et 1976 (arrondis au dixième).
\item Construis un \textbf{diagramme en barres} représentant l'évolution du taux d'urbanisation de 1960 à 2020 (échelle : 1 cm pour 10 \%, 1 cm par date).
\item Commente l'évolution de l'urbanisation au Cameroun en \textbf{5 lignes maximum} (tendance, rythme, seuils franchis, conséquences).
\end{enumerate}
\end{exercice}

\begin{exercice}[Croquis commenté : Relations ville-campagne -- Bafoussam et sa région]
Réalise un croquis schématique (format A4 paysage) représentant les relations entre la ville de Bafoussam (chef-lieu de la région de l'Ouest) et son arrière-pays rural.
\begin{enumerate}
\item \textbf{Éléments à représenter :}
\begin{itemize}
\item La ville de Bafoussam (point central, symbole urbain).
\item Les axes routiers principaux (route vers Mbouda et Bamenda au nord, route vers Douala et Yaoundé, axes secondaires vers les départements).
\item Les zones de production agricole (café, maïs, haricot, pomme de terre, maraîchage) : hachures ou couleurs distinctes.
\item Les zones d'élevage (bovins, porcs, volailles) : symboles adaptés.
\item Les marchés ruraux périodiques (points) et le grand marché urbain (symbole fort).
\end{itemize}
\item \textbf{Flux à représenter par des flèches (épaisseur proportionnelle, code couleur) :}
\begin{itemize}
\item Flux de produits agricoles et d'élevage (Campagne $\to$ Ville) : flèches vertes.
\item Flux de produits manufacturiers, intrants (engrais, aliments pour bétail), services (Ville $\to$ Campagne) : flèches rouges.
\item Flux de personnes (exode rural, navetteurs, retours de vacances) : flèches bleues pointillées.
\item Flux financiers (transferts des migrants, crédit, épargne) : flèches orange pointillées.
\end{itemize}
\item \textbf{Légende organisée} (titre, auteur, date, échelle indicative, orientation, signes conventionnels classés par thèmes).
\item \textbf{Commentaire (10 lignes) :} analyse la complémentarité ville-campagne illustrée par ce croquis.
\end{enumerate}
\end{exercice}

\begin{exercice}[Étude de cas : Filière cacao dans la région du Centre-Sud]
\textbf{Document :} \emph{Le cacao est la première culture d'exportation du Cameroun. La région du Centre-Sud (Mbalmayo, Ebolowa, Sangmélima) est un bassin de production historique. La filière associe petits planteurs (1 à 3 ha), coopératives, acheteurs agréés, ports de Douala et Kribi. Les rendements stagnent (300 à 400 kg/ha) face à un potentiel de 1 500 kg/ha. Les planteurs vieillissent, les jeunes préfèrent la moto-taxi en ville. Les pistes sont impraticables en saison des pluies.}
\begin{enumerate}
\item Identifie les \textbf{acteurs} de la filière cacao cités dans le texte et classe-les (production, collecte, transformation/exportation).
\item Relève \textbf{trois contraintes majeures} qui expliquent la faible productivité (rendement/ha) et le vieillissement de la filière.
\item Propose \textbf{deux mesures concrètes} (technique, foncière, commerciale, infrastructurelle) pour relancer cette filière et fixer les jeunes au village.
\end{enumerate}
\end{exercice}

\courssub{Je me prépare au Bac}

\begin{exercice}[Sujet de dissertation type Baccalauréat Technique]
\textbf{Sujet :} \og Les problèmes de la croissance urbaine au Cameroun et les solutions envisageables \fg

\textbf{Consignes :} traite ce sujet en respectant la structure classique de la dissertation géographique : introduction (accroche, définition des termes, analyse du sujet, problématique, annonce du plan), développement en \textbf{deux parties équilibrées} (chacune avec deux arguments illustrés d'exemples camerounais précis), conclusion (bilan, ouverture). La copie doit faire environ deux pages.
\begin{enumerate}
\item \textbf{Introduction :} rédige l'introduction complète.
\item \textbf{Partie I :} les manifestations et les causes des problèmes urbains camerounais (habitat, emploi, environnement, gouvernance).
\item \textbf{Partie II :} les politiques et stratégies d'aménagement pour une ville durable (urbanisme, logement, transport, décentralisation, villes intermédiaires).
\item \textbf{Conclusion :} rédige la conclusion complète.
\end{enumerate}
\end{exercice}

\begin{exercice}[Sujet type Probatoire Technique : Agriculture, élevage et sécurité alimentaire]
\textbf{Sujet :} \og Montrer que l'agriculture et l'élevage camerounais peinent à nourrir les villes, puis proposer des pistes d'amélioration \fg

\textbf{Données d'appui :}
\begin{itemize}
\item Importations annuelles de blé : environ 900 000 t ; riz : environ 1 200 000 t ; poisson : environ 200 000 t ; lait : environ 100 000 t équivalent lait (données MINFI/Douanes récentes).
\item Production nationale de maïs : environ 2,3 Mt (besoin estimé à 3 Mt) ; manioc : environ 5 Mt ; plantain : environ 4,5 Mt.
\item Taux d'autosuffisance alimentaire global estimé à moins de 60 \% pour les céréales.
\item Superficie cultivée inférieure à 20 \% du potentiel arable (environ 9 millions de ha exploités sur près de 50 M ha de terres).
\item Rendement moyen du maïs : 1,6 t/ha (potentiel 5 à 7 t/ha) ; rendement du cacao : 350 kg/ha (potentiel 1 500 kg/ha).
\item Cheptel bovin : environ 6 millions de têtes (dont environ 80 \% dans les régions du Nord), productivité laitière faible (moins de 500 L/vache/an).
\end{itemize}

\textbf{Travail demandé :} rédige une réponse structurée en trois parties :
\begin{enumerate}
\item \textbf{Constat :} démontre, en t'appuyant \textbf{obligatoirement} sur les données chiffrées ci-dessus, le déficit de production (dépendance aux importations, faibles rendements, sous-exploitation du potentiel).
\item \textbf{Explication :} analyse les causes structurelles (itinéraires techniques, foncier, financement, infrastructures, organisation paysanne, changement climatique, exode rural).
\item \textbf{Propositions :} formule \textbf{quatre pistes d'amélioration concrètes et chiffrées si possible} (ex. : objectif de rendement, superficie à aménager, crédit agricole, transformation locale, pistes rurales) pour réduire la facture alimentaire et créer de l'emploi rural.
\end{enumerate}
\end{exercice}

\begin{exercice}[Analyse de documents croisés : Dynamique des systèmes agraires au Nord-Cameroun]
\textbf{Document 1 (texte) :} \emph{Dans le Nord (Adamaoua, Nord, Extrême-Nord), l'agriculture pluviale (sorgho, mil, maïs, coton) et l'élevage bovin extensif cohabitent. La pression démographique (densité supérieure à 100 hab/km$^2$ dans le Diamaré) réduit la jachère (de 10-15 ans à moins de 3 ans), provoquant une dégradation des sols (moins de matière organique, érosion éolienne). Le coton (SODECOTON) assure un revenu monétaire mais l'engrais (formule coton) acidifie les sols s'il n'est pas couplé à du fumier. Les parcours se rétractent, les couloirs de transhumance sont bouchés par les champs. Les conflits agriculteurs-éleveurs sont fréquents.}

\textbf{Document 2 (données) :}
\begin{center}
\begin{tabular}{|l|c|c|}
\hline
\textbf{Indicateur} & \textbf{1990} & \textbf{2020} \\
\hline
Durée moyenne de la jachère (années) & 12 & 2 \\
\hline
Superficie en culture cotonnière (ha) & 150 000 & 250 000 \\
\hline
Cheptel bovin (millions de têtes) & 3,5 & 5,5 \\
\hline
Rendement sorgho/mil (kg/ha) & 900 & 750 \\
\hline
Taux d'équipement en charrues à traction bovine (\%) & 15 & 45 \\
\hline
\end{tabular}
\end{center}

\textbf{Travail demandé :}
\begin{enumerate}
\item À partir des deux documents, \textbf{caractérise l'évolution du système agraire} du Nord-Cameroun entre 1990 et 2020 (intensification ? extensification ? dégradation ?).
\item \textbf{Explique les mécanismes} de la crise du système « agriculture-élevage » (concurrence pour l'espace, fertilité, conflits).
\item \textbf{Propose un aménagement concerté} (zonage, couloirs, fumure, cultures fourragères, organisation villageoise) pour sortir de l'impasse. Présente-le sous forme d'un \textbf{schéma annoté} (croquis simplifié d'un terroir villageois type) + légende + 5 lignes de commentaire.
\end{enumerate}
\end{exercice}



\newpage

\coursec{Corrigés des exercices}

\coursec{Leçon 03 : Les villes}

\courssub{La ville et l'urbanisation au Cameroun}

\begin{corrige}[Exercice 1 --- QCM : Urbanisation et problèmes urbains]
\begin{enumerate}
\item \textbf{B.} Le taux d'urbanisation du Cameroun est estimé à environ 58 \% en 2020 (environ 15 millions d'urbains sur 26 millions d'habitants). Les propositions A (30 \%) et C (75 \%) correspondent à des situations d'autres pays ou d'autres époques ; D est invraisemblable pour un pays africain.
\item \textbf{B.} Le quartier spontané naît d'une occupation illégale ou non planifiée du sol : habitat précaire (banco, parpaings non enduits, tôles), absence de voirie et de réseaux (eau, électricité, assainissement). Les propositions A et D décrivent un quartier planifié, C un quartier résidentiel aisé.
\item \textbf{C.} La croissance urbaine résulte de la combinaison de l'\cle{exode rural} (départ des jeunes vers les villes) et d'un \cle{accroissement naturel} élevé en ville (les urbains sont jeunes et féconds). L'immigration internationale est marginale au Cameroun.
\item \textbf{B.} Douala concentre le premier port du pays (par où transite l'essentiel du commerce extérieur) et l'essentiel du tissu industriel : c'est le \cle{poumon économique} du Cameroun, alors que Yaoundé est la capitale politique (A est donc faux).
\end{enumerate}
\textbf{Point de vigilance :} au QCM, la justification est exigée : recopier la seule lettre sans argument fait perdre la moitié des points.
\end{corrige}

\begin{corrige}[Exercice 4 --- Calculs directs : Taux d'urbanisation et productivité]
\begin{enumerate}
\item Taux d'urbanisation 1976 : $\dfrac{1,5}{7,6} \times 100 \approx {19,7} \%$. Taux d'urbanisation 2020 : $\dfrac{15}{26} \times 100 \approx {57,7} \%$.\\
Taux d'accroissement annuel moyen de la population urbaine (1976-2020, soit $n = 44$ ans) :
\begin{align*}
t &= \frac{1}{44} \times \ln\left(\frac{15}{1,5}\right) \times 100 = \frac{1}{44} \times \ln(10) \times 100 \\
&\approx \frac{2,303}{44} \times 100 \approx {5,2} \% \text{ par an.}
\end{align*}
La population urbaine camerounaise a donc crû en moyenne d'environ 5,2 \% par an, rythme très supérieur à celui de la population totale : il traduit l'ampleur de l'exode rural.
\item Rendement : $\dfrac{12,5}{5} = {2,5}$ t/ha. L'exploitant est en dessous de l'objectif de 4 t/ha : il lui manque 1,5 t/ha, soit un déficit de $\dfrac{1,5}{4} \times 100 = {37,5} \%$ par rapport à l'objectif.
\item Chargement : $\dfrac{200}{500} = {0,4}$ UBT/ha. Ce chargement faible (moins d'une tête par hectare) traduit un \textbf{élevage extensif} : les animaux disposent de vastes parcours naturels, avec une faible productivité par hectare. Un élevage intensif présenterait un chargement de plusieurs UBT/ha, voire un élevage hors sol.
\end{enumerate}
\textbf{Point de vigilance :} vérifier que $n$ est bien le nombre d'années écoulées ($2020 - 1976 = 44$) et non le nombre de dates ; toujours donner l'unité du résultat (\%, t/ha, UBT/ha).
\end{corrige}

\begin{corrige}[Exercice 6 --- Exploitation de données : Évolution de l'urbanisation au Cameroun (1960-2020)]
\begin{enumerate}
\item 1960 : $\dfrac{0,7}{5,4} \times 100 \approx {13,0} \%$ ; 1976 : $\dfrac{1,5}{7,6} \times 100 \approx {19,7} \%$. Les valeurs du tableau sont confirmées.
\item Diagramme en barres : hauteurs proportionnelles aux taux, soit à l'échelle de 1 cm pour 10 \% : 1,3 cm ; 2,0 cm ; 3,3 cm ; 5,0 cm ; 5,4 cm ; 5,7 cm ; 5,8 cm.
\begin{popfigure}
\begin{tikzpicture}[x=0.9cm,y=0.075cm]
\draw[->,thick] (0,0) -- (8.6,0) node[below left,font=\scriptsize]{Années};
\draw[->,thick] (0,0) -- (0,65) node[above,font=\scriptsize]{Taux (\%)};
\foreach \y in {10,20,30,40,50,60}{\draw[popline!50] (0,\y) -- (8.2,\y); \node[left,font=\scriptsize] at (0,\y) {\y};}
\foreach \x/\h/\a/\c in {0.7/13.0/1960/popblue, 1.9/19.7/1976/popblue, 3.1/33.3/1987/popblue, 4.3/50.3/2005/poporange, 5.5/54.1/2010/poporange, 6.7/57.0/2015/poporange, 7.9/58.1/2020/poporange}{
\fill[\c!70,draw=\c] (\x-0.4,0) rectangle (\x+0.4,\h);
\node[below,font=\scriptsize] at (\x,0) {\a};
\node[above,font=\scriptsize] at (\x,\h) {\h};}
\end{tikzpicture}
\legende{Évolution du taux d'urbanisation du Cameroun, 1960-2020 (en \% de la population totale). Échelle verticale : 10 \% par graduation. Source : données du tableau (recensements et estimations).}
\end{popfigure}
\item \textbf{Commentaire :} le taux d'urbanisation du Cameroun est passé de 13 \% en 1960 à plus de 58 \% en 2020 : il a été multiplié par plus de quatre en soixante ans. La croissance est particulièrement rapide entre 1976 et 2005 (de 19,7 \% à 50,3 \%), période du grand exode rural. Le seuil de 50 \% est franchi au milieu des années 2000 : le Cameroun devient un pays majoritairement urbain. Le rythme ralentit ensuite légèrement après 2010. Cette urbanisation rapide, mal maîtrisée, explique l'ampleur des problèmes urbains actuels (habitat spontané, chômage, insalubrité).
\end{enumerate}
\textbf{Barème indicatif (sur 10) :} vérifications 2 pts ; diagramme (échelle, barres, titres des axes) 4 pts ; commentaire 4 pts. \textbf{Point de vigilance :} un diagramme sans titre d'axes ni échelle est sanctionné ; le commentaire doit citer des chiffres du tableau.
\end{corrige}

\courssub{Les problèmes urbains}

\begin{corrige}[Exercice 5 --- Analyse de document : Quartier spontané à Yaoundé]
\begin{enumerate}
\item \textbf{Description :} le quartier présente une morphologie dense et désordonnée, sans plan d'urbanisme : les constructions se touchent, les ruelles étroites ne permettent pas le passage des véhicules (pas de voirie carrossable). L'habitat est précaire (banco, parpaings bruts). Les réseaux sont absents ou informels : eaux usées à ciel ouvert, branchements électriques sauvages, pas d'assainissement. Habitat et activités économiques (commerces, ateliers) se mélangent dans une forte promiscuité.
\item \textbf{Deux problèmes majeurs :}
\begin{itemize}
\item \emph{Problème sanitaire :} l'absence d'évacuation des eaux usées et d'assainissement favorise la stagnation des eaux, la prolifération des moustiques et la diffusion des maladies hydriques (choléra, paludisme, typhoïde).
\item \emph{Problème foncier et de sécurité :} l'occupation illégale du sol prive les habitants de titres fonciers (expulsions possibles) ; les branchements électriques sauvages et l'absence de voirie exposent aux incendies et empêchent le passage des secours.
\end{itemize}
\item \textbf{Solution :} une \emph{restructuration sur place} du quartier : ouverture de voies de desserte, régularisation foncière des occupants, raccordement aux réseaux d'eau et d'électricité, construction d'un système d'assainissement. Cette solution est préférable à la démolition-expulsion car elle préserve le tissu social et économique du quartier tout en améliorant durablement les conditions de vie.
\end{enumerate}
\textbf{Barème indicatif (sur 10) :} description 3 pts ; deux problèmes expliqués 4 pts ; solution justifiée 3 pts. \textbf{Point de vigilance :} ne pas se contenter de recopier le document : il faut \emph{expliquer} le lien entre le constat (eaux stagnantes) et la conséquence (maladies).
\end{corrige}

\begin{corrige}[Exercice 9 --- Sujet de dissertation type Baccalauréat Technique : Les problèmes de la croissance urbaine au Cameroun et les solutions envisageables]
\textbf{Introduction (attendue) :} accroche : le Cameroun est passé de 13 \% d'urbains en 1960 à plus de 58 \% en 2020 et compte aujourd'hui deux métropoles de plusieurs millions d'habitants. Définition : la croissance urbaine est l'augmentation de la population et de l'étendue des villes. Analyse du sujet : il comporte deux termes, les problèmes (constat) et les solutions (remèdes). Problématique : comment le Cameroun peut-il maîtriser une urbanisation rapide qui dépasse ses capacités de planification ? Annonce du plan : nous montrerons d'abord les manifestations et les causes des problèmes urbains, puis les politiques d'aménagement pour une ville durable.

\textbf{Partie I --- Manifestations et causes des problèmes urbains :}
\begin{itemize}
\item \emph{Habitat et environnement :} prolifération des quartiers spontanés (Mvog-Ada à Yaoundé, Nylon à Douala), habitat précaire, insalubrité, inondations dans les quartiers bas de Douala, décharges sauvages, pollution de l'air et des eaux.
\item \emph{Emploi et gouvernance :} chômage des jeunes et explosion du secteur informel (moto-taxis, vente à la sauvette), insécurité, embouteillages permanents. \emph{Causes :} exode rural massif, accroissement naturel élevé, planification urbaine dépassée, moyens financiers limités des communes.
\end{itemize}

\textbf{Partie II --- Politiques et stratégies pour une ville durable :}
\begin{itemize}
\item \emph{Urbanisme et logement :} plans directeurs d'urbanisme, restructuration des quartiers spontanés, lotissements, programmes de logements sociaux (SIC, MAETUR), régularisation foncière.
\item \emph{Transport, décentralisation et villes intermédiaires :} élargissement et entretien des voiries, échangeurs à Yaoundé, projets de transport urbain de masse ; décentralisation et transfert de compétences aux communes ; développement des villes secondaires (Bafoussam, Ngaoundéré, Maroua, Kribi) pour décongestionner Douala et Yaoundé et fixer les populations près des campagnes.
\end{itemize}

\textbf{Conclusion (attendue) :} bilan : la croissance urbaine camerounaise, rapide et mal maîtrisée, engendre de graves problèmes d'habitat, d'emploi et d'environnement, mais des politiques d'aménagement existent et doivent être renforcées et financées. Ouverture : vers la notion de ville durable ou vers le rôle des villes intermédiaires dans l'équilibre ville-campagne.

\textbf{Barème indicatif (sur 20) :} introduction complète (accroche, définitions, problématique, plan) 4 pts ; partie I (deux arguments exemplifiés) 6 pts ; partie II (deux arguments exemplifiés) 6 pts ; conclusion avec ouverture 2 pts ; présentation et langue 2 pts. \textbf{Points de vigilance :} chaque argument doit être illustré d'un exemple camerounais \emph{nommé} ; les deux parties doivent être équilibrées ; interdiction de traiter les solutions dans la partie I.
\end{corrige}

\courssub{Dossier 2 : Les relations ville-campagne}

\begin{corrige}[Exercice 2 --- Vrai ou Faux justifié : Relations ville-campagne]
\begin{enumerate}
\item \textbf{F.} Les flux dominants de produits agricoles vont bien de la campagne vers la ville (approvisionnement alimentaire), mais il existe aussi des flux inverses (produits manufacturés, intrants agricoles, vivres transformés) : les relations ville-campagne sont des \cle{échanges réciproques}.
\item \textbf{V.} Les migrants installés en ville envoient régulièrement de l'argent à leurs familles restées au village : c'est un flux financier ville $\to$ campagne qui finance la scolarisation, la santé et parfois l'investissement agricole.
\item \textbf{F.} L'exode rural prive les campagnes de bras valides, mais cela \emph{freine} le plus souvent l'agriculture (vieillissement des exploitants, friches, baisse des rendements) au lieu de la moderniser : le mot « toujours » rend l'affirmation fausse.
\item \textbf{V.} Les villes concentrent les commerces, les industries, les hôpitaux, les lycées et les administrations qui desservent l'arrière-pays rural : c'est la fonction de \cle{commandement} et de service des villes.
\end{enumerate}
\textbf{Point de vigilance :} se méfier des mots absolus (« exclusivement », « toujours ») : en géographie, ils signalent presque toujours une affirmation fausse.
\end{corrige}

\begin{corrige}[Exercice 7 --- Croquis commenté : Relations ville-campagne -- Bafoussam et sa région]
\textbf{Attendus du croquis :} Bafoussam occupe une position centrale ; les routes rayonnent vers les villes et villages environnants (Bandjoun, Mbouda, Dschang, Foumban, axe vers Douala et Yaoundé) ; les zones caféières à l'ouest et au sud, les zones de maïs-haricot et de maraîchage autour de la ville, les pommes de terre vers les hautes terres de l'Ouest ; les symboles d'élevage (porcs, volailles) près des zones densément peuplées ; les marchés périodiques villageois reliés au grand marché de Bafoussam. Les flèches vertes convergent vers la ville (denrées), les flèches rouges divergent (manufactures, intrants), les pointillés bleus et oranges sont bidirectionnels (personnes, argent).

\begin{popfigure}
\begin{tikzpicture}[scale=0.9,font=\scriptsize]
% Zones de production
\fill[popgreenL] (-4.5,-2.2) rectangle (-1.2,0.4);
\fill[popgoldL] (-4.5,0.6) rectangle (-1.2,2.6);
\fill[poporangeL] (1.2,-2.4) rectangle (4.6,-0.2);
\fill[popgreen!15] (1.2,0.4) rectangle (4.6,2.6);
% Ville
\node[circle,fill=popink,inner sep=2.5pt,label={[font=\small\bfseries]right:Bafoussam}] (B) at (0,0) {};
\node[below=2pt of B,font=\scriptsize\itshape] {grand marché};
% Routes
\draw[thick,popdark] (B) -- (-3,1.6) node[circle,fill=popdark,inner sep=1pt,label=above:Mbouda]{};
\draw[thick,popdark] (B) -- (-3,-1.2) node[circle,fill=popdark,inner sep=1pt,label=below:Bandjoun]{};
\draw[thick,popdark] (B) -- (3,1.5) node[circle,fill=popdark,inner sep=1pt,label=above:Foumban]{};
\draw[thick,popdark] (B) -- (3,-1.4) node[circle,fill=popdark,inner sep=1pt,label=below:Dschang]{};
% Marchés périodiques
\foreach \p in {(-2,0.6),(-1.8,-0.7),(2,0.8),(1.9,-0.6)}{\node[draw=popdark,fill=white,inner sep=1.2pt,circle] at \p {};}
% Flux
\draw[-{Stealth},line width=2.2pt,popgreen] (-2.8,1.4) -- (-0.4,0.25);
\draw[-{Stealth},line width=2.2pt,popgreen] (-2.6,-1.1) -- (-0.4,-0.15);
\draw[-{Stealth},line width=2.2pt,popgreen] (2.6,-1.2) -- (0.4,-0.2);
\draw[-{Stealth},line width=2.2pt,popgreen] (2.7,1.3) -- (0.4,0.2);
\draw[-{Stealth},line width=1.6pt,poppink] (0.2,0.35) -- (2.4,1.2);
\draw[-{Stealth},line width=1.6pt,poppink] (-0.2,-0.35) -- (-2.4,-1.0);
\draw[-{Stealth},dashed,popblue] (-2.2,0.5) -- (-0.35,0.1);
\draw[-{Stealth},dashed,poporange] (0.35,-0.1) -- (2.2,-0.55);
% Nord
\draw[-{Stealth},thick] (5.1,1.6) -- (5.1,2.5); \node[above] at (5.1,2.5) {N};
% Légende graphique
\node[anchor=west] at (-4.6,-3.0) {\textcolor{popgreen}{\rule{14pt}{3pt}} produits agricoles et d'élevage (campagne $\to$ ville)};
\node[anchor=west] at (-4.6,-3.5) {\textcolor{poppink}{\rule{14pt}{3pt}} manufacturiers, intrants, services (ville $\to$ campagne)};
\node[anchor=west] at (1.2,-3.0) {\textcolor{popblue}{- - -} flux de personnes \quad \textcolor{poporange}{- - -} flux financiers};
\node[anchor=west] at (1.2,-3.5) {\textcolor{popgreen!60!black}{$\blacksquare$} café \quad \textcolor{popgold}{$\blacksquare$} maïs-haricot-maraîchage \quad \textcolor{poporange}{$\blacksquare$} élevage \quad \textcolor{popgreen!30}{$\blacksquare$} pomme de terre};
\end{tikzpicture}
\legende{Croquis schématique des relations ville-campagne autour de Bafoussam (Ouest-Cameroun). 1 : ville et grand marché ; 2 : routes ; 3 : marchés périodiques ; 4 : flux de denrées ; 5 : flux de biens et services ; 6 : flux de personnes et d'argent. Échelle indicative : 1 cm $\approx$ 15 km.}
\end{popfigure}

\textbf{Commentaire attendu (10 lignes) :} le croquis illustre la complémentarité fonctionnelle entre Bafoussam et sa campagne. La campagne approvisionne la ville en produits vivriers (maïs, haricots, pommes de terre, légumes) et en produits de rente (café) qui alimentent le grand marché et les filières d'exportation. En sens inverse, la ville fournit à la campagne des biens manufacturés, des intrants agricoles (engrais, semences, aliments pour bétail) et des services (santé, éducation, administration, banques). Les flux de personnes sont intenses : exode rural des jeunes vers Bafoussam, navettes quotidiennes des commerçants, retours des migrants aux périodes de fêtes. Les transferts d'argent des migrants urbains financent l'investissement rural (champs, maisons, scolarisation). Cette interdépendance fait de la ville un pôle de commandement et de la campagne un arrière-pays nourricier : le développement de l'un est impossible sans l'autre, d'où la nécessité de désenclaver les pistes rurales pour fluidifier ces échanges.

\textbf{Barème indicatif (sur 20) :} éléments localisés 5 pts ; flux corrects (sens, couleurs, épaisseurs) 5 pts ; légende organisée (titre, orientation, échelle, signes classés) 5 pts ; commentaire 5 pts. \textbf{Point de vigilance :} un croquis sans légende ni flèche du nord est systématiquement pénalisé ; les flèches doivent respecter le sens des flux.
\end{corrige}

\begin{corrige}[Exercice 11 --- Analyse de documents croisés : Dynamique des systèmes agraires au Nord-Cameroun]
\begin{enumerate}
\item \textbf{Évolution du système agraire :} entre 1990 et 2020, le système connaît à la fois une \emph{extensification} (le coton passe de 150 000 à 250 000 ha, soit $+67 \%$) et une \emph{intensification} partielle (équipement en charrues à traction bovine multiplié par trois, de 15 \% à 45 \%). Mais la jachère s'effondre de 12 à 2 ans et le rendement du sorgho/mil recule de 900 à 750 kg/ha, soit une baisse de $\dfrac{900-750}{900} \times 100 \approx {17} \%$, tandis que le cheptel bovin passe de 3,5 à 5,5 millions de têtes sur des parcours en rétraction. C'est donc une \cle{intensification inachevée et non durable} : on cultive plus, sur plus d'espace, mais la fertilité se dégrade.
\item \textbf{Mécanismes de la crise :} la croissance démographique (densité supérieure à 100 hab/km$^2$ dans le Diamaré) étend les champs au détriment des parcours et des jachères ; sans repos du sol ni fumure organique suffisante, la fertilité baisse et les rendements chutent, ce qui pousse à défricher davantage : c'est un \emph{cercle vicieux}. Les champs bouchent les couloirs de transhumance : les troupeaux, plus nombreux, divaguent dans les cultures, d'où des conflits récurrents entre agriculteurs et éleveurs. L'engrais minéral coton, utilisé sans fumier, acidifie les sols et aggrave la dégradation.
\item \textbf{Aménagement concerté :}
\begin{popfigure}
\begin{tikzpicture}[scale=0.95,font=\scriptsize]
% Terroir
\draw[thick,popdark,rounded corners] (-5,-3) rectangle (5,3);
% Village
\fill[popink!15] (-0.9,-0.7) rectangle (0.9,0.7);
\node[font=\scriptsize\bfseries] at (0,0) {Village};
% Champs permanents proches
\fill[popgreenL] (-2.6,-1.8) rectangle (-1.1,1.8);
\node[rotate=90] at (-1.85,0) {champs permanents fumés};
% Champs coton/céréales
\fill[popgoldL] (1.1,-1.8) rectangle (3.2,1.8);
\node[rotate=90,align=center] at (2.15,0) {coton-céréales\\rotation + fourragères};
% Parcours
\fill[poppurpleL] (-4.8,2.0) rectangle (4.8,2.85);
\node at (0,2.42) {parcours délimité (bovins)};
\fill[poppurpleL] (-4.8,-2.85) rectangle (4.8,-2.0);
\node at (0,-2.42) {parcours délimité};
% Couloir de transhumance
\draw[line width=5pt,poporange!60] (3.5,-2.85) -- (3.5,2.85);
\node[rotate=90,popdark] at (3.85,0) {couloir de transhumance protégé};
% Point d'eau
\node[circle,fill=popblue,inner sep=2pt,label=right:{forage}] at (-3.8,2.42) {};
\node[circle,fill=popblue,inner sep=2pt,label=right:{mare}] at (-3.8,-2.42) {};
% Flèches fumier/résidus
\draw[-{Stealth},thick,popgreen!60!black] (1.1,1.0) -- (-1.1,1.0) node[midway,above]{fumier};
\draw[-{Stealth},thick,popgold!80!black] (-1.1,-1.2) -- (1.1,-1.2) node[midway,below]{résidus de culture};
% Nord
\draw[-{Stealth},thick] (5.4,2.2) -- (5.4,2.9); \node[above] at (5.4,2.9) {N};
\end{tikzpicture}
\legende{Schéma d'un terroir villageois aménagé du Nord-Cameroun. 1 : village ; 2 : champs permanents fumés ; 3 : champs de coton et céréales en rotation avec cultures fourragères ; 4 : parcours délimités ; 5 : couloir de transhumance protégé ; 6 : points d'eau hors des champs ; 7 : échanges fumier/résidus entre élevage et agriculture.}
\end{popfigure}
\emph{Commentaire (5 lignes) :} le zonage concerté sépare les espaces de culture et de parcours tout en les rendant complémentaires : le fumier des troupeaux fertilise les champs (association agriculture-élevage), les résidus de culture et les fourragères nourrissent le bétail, et les couloirs protégés évitent la divagation. La rotation et la fumure organique restaurent la fertilité des sols. Un comité villageois de gestion du terroir, associant agriculteurs et éleveurs, prévient les conflits. Ce système agro-pastoral intégré permet d'intensifier durablement sans dégrader les sols.
\end{enumerate}
\textbf{Barème indicatif (sur 20) :} évolution caractérisée avec calculs 5 pts ; mécanismes expliqués 5 pts ; schéma annoté et légende 5 pts ; commentaire 5 pts. \textbf{Points de vigilance :} croiser explicitement les deux documents (le texte explique les chiffres du tableau) ; le schéma doit montrer la \emph{complémentarité} agriculture-élevage (flèches fumier/résidus), pas seulement la séparation des espaces.
\end{corrige}

\courssub{L'agriculture au Cameroun}

\begin{corrige}[Exercice 3 --- Définitions et distinctions : Agriculture et Élevage]
\begin{enumerate}
\item L'\textbf{agriculture vivrière} produit des cultures destinées essentiellement à l'autoconsommation et au marché local (maïs, manioc, plantain, mil, sorgho) ; l'\textbf{agriculture de rente} produit des cultures destinées à la vente, notamment à l'exportation (cacao, café, coton, banane, hévéa, palmier à huile). La distinction porte sur la \emph{destination} de la production, non sur la taille de l'exploitation.
\item L'\textbf{élevage extensif} laisse les animaux paître librement sur de vastes parcours naturels, avec peu d'investissements et de faibles rendements (élevage bovin traditionnel de l'Adamaoua et du Nord) ; l'\textbf{élevage intensif} concentre les animaux dans des bâtiments ou de petits espaces, avec alimentation complémentée, soins vétérinaires et forte productivité (aviculture et porcheries périurbaines de l'Ouest et du Centre).
\item La \textbf{jachère} est le repos temporaire d'une parcelle (plusieurs années) pour permettre au sol de reconstituer sa fertilité ; la \textbf{rotation culturale} est la succession planifiée de cultures différentes sur une même parcelle d'une année à l'autre pour préserver la fertilité \emph{sans} laisser la terre au repos.
\item Deux produits de rente majeurs du Cameroun : le \textbf{cacao} (première culture d'exportation) et le \textbf{café} ; sont aussi acceptés le coton, la banane ou l'huile de palme.
\end{enumerate}
\textbf{Point de vigilance :} une définition sans exemple camerounais est une définition incomplète au Probatoire : toujours illustrer.
\end{corrige}

\begin{corrige}[Exercice 8 --- Étude de cas : Filière cacao dans la région du Centre-Sud]
\begin{enumerate}
\item \textbf{Acteurs :} production : les petits planteurs (exploitations de 1 à 3 ha) ; collecte : les coopératives et les acheteurs agréés ; exportation : les ports de Douala et de Kribi. (La transformation industrielle, encore limitée au Cameroun, n'est pas citée dans le texte.)
\item \textbf{Trois contraintes majeures :}
\begin{itemize}
\item \emph{Contrainte technique :} vergers vieillissants et faible adoption des intrants, d'où des rendements de 300 à 400 kg/ha, très inférieurs au potentiel de 1 500 kg/ha.
\item \emph{Contrainte humaine :} vieillissement des planteurs et départ des jeunes vers les activités urbaines (moto-taxi), jugées plus rentables : la relève n'est pas assurée.
\item \emph{Contrainte infrastructurelle :} pistes rurales impraticables en saison des pluies, ce qui renchérit l'évacuation de la production et réduit les prix payés au producteur.
\end{itemize}
\item \textbf{Deux mesures :}
\begin{itemize}
\item \emph{Technique et commerciale :} distribution de plants améliorés et d'intrants via les coopératives, encadrement des planteurs et garantie d'un prix plancher rémunérateur pour rendre le cacao attractif face à la moto-taxi.
\item \emph{Infrastructurelle :} réhabilitation et entretien des pistes rurales du bassin Centre-Sud pour désenclaver les villages et réduire les coûts d'évacuation ; on peut ajouter la création d'unités locales de transformation (fermentation, séchage, broyage) créatrices d'emplois pour les jeunes.
\end{itemize}
\end{enumerate}
\textbf{Barème indicatif (sur 10) :} acteurs classés 3 pts ; trois contraintes expliquées 4 pts ; deux mesures justifiées 3 pts. \textbf{Point de vigilance :} chaque contrainte doit être reliée à un élément précis du texte (chiffre ou fait) ; une mesure non justifiée (« il faut aider les planteurs ») ne vaut que la moitié des points.
\end{corrige}

\courssub{L'élevage au Cameroun}

\begin{corrige}[Exercice 3 --- Définitions et distinctions : Agriculture et Élevage (suite)]
\emph{Voir les définitions de l'élevage extensif/intensif, jachère/rotation et produits de rente dans l'exercice 3 ci-dessus (partie Agriculture).}
\end{corrige}

\begin{corrige}[Exercice 4 --- Calculs directs : Taux d'urbanisation et productivité (suite)]
\emph{Voir le calcul de chargement (0,4 UBT/ha) et son interprétation (élevage extensif) dans l'exercice 4 ci-dessus (partie Urbanisation).}
\end{corrige}

\courssub{Synthèse : villes, campagnes et développement}

\begin{corrige}[Exercice 10 --- Sujet type Probatoire Technique : Agriculture, élevage et sécurité alimentaire]
\textbf{1. Constat :} l'agriculture et l'élevage camerounais ne couvrent pas les besoins du pays. Le taux d'autosuffisance céréalière est inférieur à 60 \% : le pays importe chaque année environ 900 000 t de blé (non produit localement), 1 200 000 t de riz, 200 000 t de poisson et 100 000 t équivalent lait. La production nationale de maïs (2,3 Mt) reste inférieure au besoin estimé (3 Mt), soit un déficit d'environ 0,7 Mt. Les rendements sont très faibles : 1,6 t/ha pour le maïs contre un potentiel de 5 à 7 t/ha, et 350 kg/ha pour le cacao contre 1 500 kg/ha possibles. Enfin, moins de 20 \% du potentiel arable est cultivé (environ 9 M ha sur près de 50 M ha de terres) et la vache camerounaise produit moins de 500 L de lait par an. Le pays peine donc à nourrir des villes en pleine croissance.

\textbf{2. Explication :} les causes sont structurelles : techniques culturales traditionnelles (houe, peu d'intrants, semences non améliorées), foncier coutumier peu sécurisé qui décourage l'investissement, accès quasi nul au crédit agricole, pistes rurales dégradées qui isolent les bassins de production, faible organisation des paysans en coopératives, aléas climatiques croissants (sécheresses au Nord, pluies irrégulières) et exode rural qui vieillit la main-d'œuvre agricole.

\textbf{3. Propositions (quatre pistes) :}
\begin{itemize}
\item Doubler le rendement du maïs de 1,6 à 3,2 t/ha en cinq ans par la diffusion de semences améliorées et d'engrais subventionnés : à superficie égale, cela comblerait le déficit de 0,7 Mt.
\item Aménager et mettre en culture 1 million d'hectares supplémentaires (vallées, bas-fonds) avec sécurisation foncière (titres fonciers, certificats d'usage).
\item Créer un fonds de crédit agricole à taux réduit accessible via les coopératives et développer la transformation locale (minoteries, semouleries de maïs et de manioc, laiteries) pour créer des emplois ruraux et fixer les jeunes.
\item Réhabiliter plusieurs milliers de kilomètres de pistes rurales et aménager des couloirs de transhumance au Nord pour sécuriser l'élevage (6 millions de têtes) et réduire les importations de lait et de viande.
\end{itemize}
\textbf{Barème indicatif (sur 20) :} constat chiffré (au moins quatre données exploitées) 6 pts ; explication des causes 6 pts ; quatre propositions concrètes 6 pts ; cohérence et présentation 2 pts. \textbf{Points de vigilance :} la consigne impose d'utiliser \emph{obligatoirement} les chiffres fournis : toute copie sans donnée chiffrée est hors sujet ; les propositions doivent être réalistes et reliées aux causes identifiées.
\end{corrige}

\newpage

\coursec{Fiche bilan}

\begin{popfigure}
\begin{center}
\begin{tikzpicture}[mindmap, grow cyclic,
  every node/.style={concept, font=\small},
  root concept/.append style={concept color=popblue, font=\bfseries},
  level 1/.append style={level distance=3.4cm, sibling angle=60},
  level 2/.append style={level distance=2.2cm, font=\scriptsize}]
\node[root concept]{Les villes}
  child[concept color=poporange]{ node {Ville et urbanisation}
    child { node {Croissance rapide} }
    child { node {Douala, Yaoundé} }
  }
  child[concept color=poppink]{ node {Problèmes urbains}
    child { node {Habitat, chômage} }
    child { node {Insalubrité} }
  }
  child[concept color=popgreen]{ node {Relations ville-campagne}
    child { node {Exode rural} }
    child { node {Approvisionnement} }
  }
  child[concept color=popgold]{ node {Agriculture}
    child { node {Vivrière} }
    child { node {De rente} }
  }
  child[concept color=poppurple]{ node {Élevage}
    child { node {Traditionnel} }
    child { node {Moderne} }
  }
  child[concept color=popblueL]{ node {Développement}
    child { node {Équilibre villes-campagnes} }
  };
\end{tikzpicture}
\end{center}
\legende{Carte mentale de la leçon : les six grandes parties à maîtriser (schéma simplifié).}
\end{popfigure}

\begin{aretenir}[L'essentiel de la leçon : définitions clés]
\begin{itemize}
  \item \cle{Ville} : agglomération concentrant une population nombreuse (au Cameroun, une localité d'au moins \num{5000} habitants est classée urbaine) exerçant des activités surtout tertiaires (commerce, services, administration) et industrielles.
  \item \cle{Urbanisation} : augmentation de la population urbaine et extension de l'espace bâti des villes. Elle se mesure par le \cle{taux d'urbanisation} (part de la population urbaine dans la population totale).
  \item \cle{Exode rural} : départ massif et durable des campagnards, surtout des jeunes, vers les villes.
  \item \cle{Agriculture vivrière} : production destinée d'abord à l'autoconsommation de la famille (manioc, maïs, plantain, igname, arachide).
  \item \cle{Agriculture de rente} (ou de plantation, ou d'exportation) : production destinée à la vente, notamment à l'exportation (cacao, café, coton, banane, hévéa, palmier à huile).
  \item \cle{Élevage traditionnel} : élevage extensif (grands troupeaux en libre parcours), de faible productivité ; \cle{élevage moderne} : élevage intensif, avec animaux sélectionnés, alimentation suivie et soins vétérinaires, plus rentable.
\end{itemize}
\end{aretenir}

\begin{aretenir}[Chiffres et faits à connaître par c\oe ur]
\begin{center}
\begin{tabularx}{\linewidth}{|l|X|}\hline
\rowcolor{popblueL} \textbf{Notion} & \textbf{Repère chiffré ou localisé} \\ \hline
Taux d'urbanisation du Cameroun & plus de la moitié de la population vit en ville (environ 55 \%) \\ \hline
Deux métropoles & Douala (capitale économique, premier port du pays) et Yaoundé (capitale politique et administrative) ; chacune dépasse 3 millions d'habitants \\ \hline
Problèmes urbains majeurs & habitat spontané (bidonvilles), chômage et sous-emploi, insalubrité, insécurité, embouteillages, encombrement des marchés \\ \hline
Agriculture & occupe la majorité de la population active ; deux types : vivrière et de rente \\ \hline
Élevage & bovin surtout dans l'Adamaoua et l'Extrême-Nord ; élevage avicole (poulets) et porcin (porcs) surtout dans l'Ouest et autour des grandes villes \\ \hline
Relations ville-campagne & la campagne nourrit la ville et lui envoie des hommes (exode) ; la ville offre marchés, services (hôpitaux, écoles) et emplois \\ \hline
\end{tabularx}
\end{center}
\end{aretenir}

\begin{methode}[Méthodes express à retenir]
\begin{enumerate}
  \item \textbf{Croquis} : toujours un titre, une légende numérotée, la flèche du nord et une échelle indicative ; peu de couleurs, des traits clairs.
  \item \textbf{Étude de cas} : illustrer chaque notion par un exemple camerounais précis (une ville, une culture, une région d'élevage).
  \item \textbf{Calcul d'un taux de variation} :
  \[ \text{taux de variation} = \frac{\text{valeur finale} - \text{valeur initiale}}{\text{valeur initiale}} \times 100 \]
  Exemple : une ville passe de \num{100000} à \num{130000} habitants : taux $= \dfrac{130000 - 100000}{100000} \times 100 = 30\,\%$.
  \item \textbf{Conclusion} : répondre d'abord au problème posé, puis ouvrir sur une solution (aménagement des villes, appui au monde rural).
\end{enumerate}
\end{methode}

\begin{attention}[Erreurs fréquentes à éviter]
\begin{itemize}
  \item Ne pas confondre \cle{ville} (l'espace bâti) et \cle{urbanisation} (le processus de croissance des villes).
  \item Ne pas écrire que l'exode rural ne concerne que les vieillards : ce sont surtout les \textbf{jeunes} qui partent, ce qui vieillit et appauvrit les campagnes.
  \item Ne pas opposer totalement vivrière et de rente : une même exploitation pratique souvent les deux (agriculture mixte).
  \item Dans un calcul de pourcentage, ne jamais oublier de multiplier par 100 et d'écrire le symbole \%.
\end{itemize}
\end{attention}

\begin{aretenir}[Checklist avant le Probatoire / le Bac technique]
\begin{itemize}
  \item $\square$ Je sais définir ville, urbanisation et exode rural.
  \item $\square$ Je connais les deux plus grandes villes du Cameroun et leurs fonctions.
  \item $\square$ Je cite au moins quatre problèmes urbains avec un exemple chacun.
  \item $\square$ J'explique les causes et les conséquences de l'exode rural.
  \item $\square$ Je distingue agriculture vivrière et agriculture de rente avec des exemples de cultures.
  \item $\square$ Je localise les grandes zones d'élevage du Cameroun.
  \item $\square$ Je différencie élevage traditionnel et élevage moderne.
  \item $\square$ Je sais décrire les échanges (hommes, biens, services) entre ville et campagne.
  \item $\square$ Je sais calculer un pourcentage et un taux de variation.
  \item $\square$ Je sais réaliser un croquis simple avec titre, légende et orientation.
  \item $\square$ Je sais rédiger une démarche justifiée et conclure en répondant au sujet.
\end{itemize}
\end{aretenir}

\findecours

\end{document}
